% Lesson 41 — Grammar: Phrasal verbs — Anglais Tle A — Cours signé OnBuch+
% Programme officiel MINESEC. Compiler avec : tectonic EN41-grammar-phrasal-verbs.tex
\newcommand{\DOCMATIERE}{Anglais}
\newcommand{\DOCNIVEAU}{Tle A}
\newcommand{\DOCMODULE}{Chapter 4 : Citizenship and Human Rights}
\newcommand{\DOCLECON}{EN41}
\newcommand{\DOCTITRE}{Lesson 41 — Grammar: Phrasal verbs}
% OnBuch+ — préambule « cours signés » ENRICHI (compilé avec tectonic / XeLaTeX)
% Système visuel « Pop » d'OnBuch+ : fond crème (jamais blanc), bordures encre
% épaisses, ombres « dures » décalées, titres Archivo Black en capitales.
% Version MATHÉMATIQUES : graphes (pgfplots/tikz), boîtes pédagogiques
% (définition, propriété, méthode, attention, le savais-tu, exercice résolu,
% exercice, TP, bilan), page de garde et signature OnBuch+ ancrée en bas.
%
% Macros attendues dans main.tex AVANT \input{preamble} :
%   \DOCMATIERE  \DOCNIVEAU  \DOCMODULE  \DOCLECON (numéro)  \DOCTITRE
\documentclass[11pt]{article}

\usepackage{fontspec}
\usepackage[english,french]{babel} % français = langue principale ; l'anglais via \eng{..} / textebox
\usepackage[a4paper,margin=2cm,top=3.4cm,bottom=2.8cm,headheight=1.4cm,footskip=1.3cm]{geometry}
\usepackage{amsmath,amssymb,amsfonts}
\usepackage{mathtools} % \xmapsto, \coloneqq... souvent utilisés par les modèles
\usepackage{cancel} % simplifications barrées (bilans de forces)
% \abs{z} (module, valeur absolue) et \norm{u} (norme), utilisés par les modèles
\providecommand{\abs}{}\let\abs\relax\DeclarePairedDelimiter{\abs}{\lvert}{\rvert}
\providecommand{\norm}{}\let\norm\relax\DeclarePairedDelimiter{\norm}{\lVert}{\rVert}
\usepackage[table]{xcolor} % [table] : fournit aussi \rowcolors (alternance de lignes)
\usepackage{pagecolor}
\usepackage{tikz}
\usetikzlibrary{arrows.meta,positioning,calc,shapes.geometric,shapes.misc,shapes.arrows,decorations.pathmorphing,decorations.pathreplacing,decorations.markings,patterns,shadows,mindmap,trees,fit,backgrounds,matrix,intersections,angles,quotes}
\usepackage{pgfplots}
\usepackage[version=4]{mhchem} % équations nucléaires \ce{^{235}_{92}U}
\usepackage[european,siunitx]{circuitikz} % schémas électriques
\pgfplotsset{compat=1.18}
\usepgfplotslibrary{fillbetween,groupplots} % aires entre courbes (calcul intégral)
\usepackage{enumitem}
\usepackage{fancyhdr}
% [most] charge toutes les bibliothèques tcolorbox (listings, minted, raster,
% documentation...) et gonfle inutilement la mémoire TeX sur un document long
% (~300 boîtes) : on ne charge que ce qui est réellement utilisé.
\usepackage[skins,breakable]{tcolorbox}
\usepackage{graphicx}
\usepackage{colortbl}
\usepackage{booktabs}
\usepackage{tabularx}
\usepackage{diagbox} % case d'en-tête diagonale des tableaux à double entrée
% Colonnes extensibles centrée / gauche / droite (C, L, R), souvent utilisées par les modèles
\newcolumntype{C}{>{\centering\arraybackslash}X}
\newcolumntype{L}{>{\raggedright\arraybackslash}X}
\newcolumntype{R}{>{\raggedleft\arraybackslash}X}
\usepackage{array}
\usepackage{multicol}
\usepackage{multirow}
\usepackage{siunitx}
% parse-numbers=false : accepte aussi \SI{2,5 \times 10^{-3}}{...}
\sisetup{locale=FR,output-decimal-marker={,},group-minimum-digits=4,parse-numbers=false}
\DeclareSIUnit\ohm{\ensuremath{\symup{\Omega}}} % U+2126 absent de la police
\DeclareSIUnit\division{div} % réglages d'oscilloscope (ms/div, V/div)
\DeclareSIUnit\tr{tr} % tours (vitesse de rotation, ex. \SI{1500}{\tr\per\minute})
\DeclareSIUnit\molar{M} % concentration molaire (ex. \SI{3}{\molar}, \SI{1}{\micro\molar})
\DeclareSIUnit\an{an} % année (le modèle confond parfois avec \year, un compteur TeX) — ex. \SI{5}{\kilo\meter\per\an}
\DeclareSIUnit\FCFA{FCFA} % franc CFA (ex. \SI{500}{\FCFA\per\kilo\gram})
\DeclareSIUnit\degreeBrix{\textdegree Brix} % teneur en sucre d'un sirop/jus (réfractométrie)
\DeclareSIUnit\month{mois} % le modèle utilise parfois \month, un compteur TeX, comme unité de durée
\usepackage{qrcode}
% \degree : commande fréquemment utilisée par les modèles pour un angle ou une
% température en dehors de \SI{}{} (non fournie par les paquets chargés).
\providecommand{\degree}{^{\circ}}
% \qed (fin de démonstration), utilisé par les modèles sans amsthm
\providecommand{\qed}{\hfill$\square$}
% \barre : double barre des valeurs interdites dans un tableau de variations (tkz-tab)
\providecommand{\barre}{\ensuremath{\|}}
% environnement proof (amsthm non chargé) utilisé par les modèles
\ifdefined\proof\else\newenvironment{proof}[1][Démonstration]{\par\noindent\textit{#1.}\ }{\qed\par}\fi
\usepackage{lastpage}
\usepackage{hyperref}
\hypersetup{hidelinks}

% ============================================================
% Palette « Pop » OnBuch+ — mêmes hex que la classe P (plus_ui.dart)
% ============================================================
\definecolor{poporange}{HTML}{F59321}
\definecolor{popdark}{HTML}{E07A0C}
\definecolor{popcream}{HTML}{FFF6E8}
\definecolor{popink}{HTML}{1C1714}
\definecolor{poppurple}{HTML}{7A5AE0}
\definecolor{popgreen}{HTML}{2E9E6A}
\definecolor{poppink}{HTML}{E0457B}
\definecolor{popblue}{HTML}{2F7DE1}
\definecolor{popgold}{HTML}{C98A12}
\definecolor{popmuted}{HTML}{8A7F76}
\definecolor{popline}{HTML}{EADFD2}
\definecolor{popgray}{HTML}{8A7F76}
\colorlet{popgrey}{popgray}
% Alias tolérants (le modèle nomme parfois les couleurs différemment)
\colorlet{popwhite}{white}
\colorlet{popblack}{popink}
\colorlet{popred}{poppink}
\colorlet{popyellow}{popgold}
% Teintes claires (fonds de tableaux/encadrés) — le modèle nomme parfois
% les couleurs avec un suffixe L (ex. popblueL) sans les avoir définies.
\colorlet{poporangeL}{poporange!20}
\colorlet{popdarkL}{popdark!20}
\colorlet{poppurpleL}{poppurple!20}
\colorlet{popgreenL}{popgreen!20}
\colorlet{poppinkL}{poppink!20}
\colorlet{popblueL}{popblue!20}
\colorlet{popgoldL}{popgold!20}
\colorlet{popredL}{popred!20}
\colorlet{popgrayL}{popgray!20}
% Teintes claires pour fonds de tableaux / zones
\colorlet{popblueL}{popblue!10!white}
\colorlet{poporangeL}{poporange!14!white}
\colorlet{popgreenL}{popgreen!12!white}
\colorlet{poppurpleL}{poppurple!12!white}
\colorlet{poppinkL}{poppink!12!white}
\colorlet{popgoldL}{popgold!14!white}

\pagecolor{popcream}

% ============================================================
% Typographie
% ============================================================
\defaultfontfeatures{Ligatures=TeX}
\setmainfont{PlusJakartaSans-Medium.ttf}[Path=../../../fonts/,BoldFont=PlusJakartaSans-Bold.ttf]
% Maths en sans-serif (Fira Math) avec chiffres et lettres droites de Plus
% Jakarta, pour que les formules se fondent dans le texte.
\usepackage[math-style=ISO,bold-style=ISO]{unicode-math}
\setmathfont{FiraMath-Regular.otf}[Path=../../../fonts/]
\setmathfont{PlusJakartaSans-Medium.ttf}[Path=../../../fonts/,range={up/num,up/{Latin,latin},\mathup}]
\usepackage{icomma} % virgule décimale sans espace en mode math
% Caractères mathématiques Unicode absents des polices de texte (Plus Jakarta,
% Archivo Black) : rendus par leur équivalent mathématique, en texte comme en
% formule — sinon ils s'affichent en carré vide (ex. « Arithmétique dans ℤ »).
\usepackage{newunicodechar}
\newunicodechar{ℝ}{\ensuremath{\mathbb{R}}}
\newunicodechar{ℤ}{\ensuremath{\mathbb{Z}}}
\newunicodechar{ℂ}{\ensuremath{\mathbb{C}}}
\newunicodechar{ℕ}{\ensuremath{\mathbb{N}}}
\newunicodechar{ℚ}{\ensuremath{\mathbb{Q}}}
\newunicodechar{𝔻}{\ensuremath{\mathbb{D}}}
\newunicodechar{α}{\ensuremath{\alpha}}
\newunicodechar{β}{\ensuremath{\beta}}
\newunicodechar{γ}{\ensuremath{\gamma}}
\newunicodechar{δ}{\ensuremath{\delta}}
\newunicodechar{ε}{\ensuremath{\varepsilon}}
\newunicodechar{θ}{\ensuremath{\theta}}
\newunicodechar{λ}{\ensuremath{\lambda}}
\newunicodechar{μ}{\ensuremath{\mu}}
\newunicodechar{σ}{\ensuremath{\sigma}}
\newunicodechar{τ}{\ensuremath{\tau}}
\newunicodechar{φ}{\ensuremath{\varphi}}
\newunicodechar{ψ}{\ensuremath{\psi}}
\newunicodechar{ω}{\ensuremath{\omega}}
\newunicodechar{Σ}{\ensuremath{\Sigma}}
% majuscules produites par \MakeUppercase dans les titres (α-aminés -> Α)
\newunicodechar{Α}{\ensuremath{\alpha}}
\newunicodechar{Β}{\ensuremath{\beta}}
\newunicodechar{Γ}{\ensuremath{\Gamma}}
\newunicodechar{Δ}{\ensuremath{\Delta}}
\newunicodechar{Ε}{\ensuremath{\varepsilon}}
\newunicodechar{Θ}{\ensuremath{\Theta}}
\newunicodechar{Λ}{\ensuremath{\Lambda}}
\newunicodechar{Μ}{\ensuremath{\mu}}
\newunicodechar{Τ}{\ensuremath{\tau}}
\newunicodechar{Φ}{\ensuremath{\Phi}}
\newunicodechar{Ψ}{\ensuremath{\Psi}}
\newunicodechar{Ω}{\ensuremath{\Omega}}
\newunicodechar{↦}{\ensuremath{\mapsto}}
\newunicodechar{∈}{\ensuremath{\in}}
\newunicodechar{∉}{\ensuremath{\notin}}
\newunicodechar{∧}{\ensuremath{\wedge}}
\newunicodechar{⇔}{\ensuremath{\Leftrightarrow}}
\newunicodechar{⟺}{\ensuremath{\Longleftrightarrow}}
\newunicodechar{⇒}{\ensuremath{\Rightarrow}}
\newunicodechar{⟹}{\ensuremath{\Longrightarrow}}
\newunicodechar{∘}{\ensuremath{\circ}}
\newunicodechar{∪}{\ensuremath{\cup}}
\newunicodechar{∩}{\ensuremath{\cap}}
\newunicodechar{≡}{\ensuremath{\equiv}}
\newunicodechar{⊂}{\ensuremath{\subset}}
\newunicodechar{⊥}{\ensuremath{\perp}}
\newunicodechar{∃}{\ensuremath{\exists}}
\newunicodechar{∀}{\ensuremath{\forall}}
\newunicodechar{∛}{\ensuremath{\sqrt[3]{\,}}}
\newunicodechar{∜}{\ensuremath{\sqrt[4]{\,}}}
\newunicodechar{⃗}{\ensuremath{{}^{\to}}}
\newunicodechar{ⁿ}{\ifmmode^{n}\else\textsuperscript{\ensuremath{n}}\fi}
\newunicodechar{ˣ}{\ifmmode^{x}\else\textsuperscript{\ensuremath{x}}\fi}
\newunicodechar{ᵇ}{\ifmmode^{b}\else\textsuperscript{\ensuremath{b}}\fi}
\newunicodechar{ʳ}{\ifmmode^{r}\else\textsuperscript{\ensuremath{r}}\fi}
\newunicodechar{ᵘ}{\ifmmode^{u}\else\textsuperscript{\ensuremath{u}}\fi}
\newunicodechar{ʸ}{\ifmmode^{y}\else\textsuperscript{\ensuremath{y}}\fi}
\newunicodechar{ᵃ}{\ifmmode^{a}\else\textsuperscript{\ensuremath{a}}\fi}
\newunicodechar{ᵉ}{\ifmmode^{e}\else\textsuperscript{\ensuremath{e}}\fi}
\newunicodechar{ᵏ}{\ifmmode^{k}\else\textsuperscript{\ensuremath{k}}\fi}
\newunicodechar{ᵖ}{\ifmmode^{p}\else\textsuperscript{\ensuremath{p}}\fi}
\newunicodechar{ᴺ}{\ifmmode^{N}\else\textsuperscript{\ensuremath{N}}\fi}
\newunicodechar{ᶜ}{\ifmmode^{c}\else\textsuperscript{\ensuremath{c}}\fi}
\newunicodechar{ʰ}{\ifmmode^{h}\else\textsuperscript{\ensuremath{h}}\fi}
\newunicodechar{ᵠ}{\ifmmode^{q}\else\textsuperscript{\ensuremath{q}}\fi}
\newunicodechar{ₙ}{\ifmmode_{n}\else\textsubscript{\ensuremath{n}}\fi}
\newunicodechar{ₐ}{\ifmmode_{a}\else\textsubscript{\ensuremath{a}}\fi}
\newunicodechar{ᵢ}{\ifmmode_{i}\else\textsubscript{\ensuremath{i}}\fi}
\newunicodechar{ᵣ}{\ifmmode_{r}\else\textsubscript{\ensuremath{r}}\fi}
\newunicodechar{ₖ}{\ifmmode_{k}\else\textsubscript{\ensuremath{k}}\fi}
\newunicodechar{ᵦ}{\ifmmode_{\beta}\else\textsubscript{\ensuremath{\beta}}\fi}
\newunicodechar{ₓ}{\ifmmode_{x}\else\textsubscript{\ensuremath{x}}\fi}
\newfontfamily\popdisplay{ArchivoBlack-Regular.ttf}[Path=../../../fonts/]
\newfontfamily\poplogo{SpaceGrotesk-Bold.ttf}[Path=../../../fonts/]
\newfontfamily\popsemi{PlusJakartaSans-SemiBold.ttf}[Path=../../../fonts/]
\newfontfamily\popxbold{PlusJakartaSans-ExtraBold.ttf}[Path=../../../fonts/]
\newfontfamily\popxboldls{PlusJakartaSans-ExtraBold.ttf}[Path=../../../fonts/,LetterSpace=3.0]

\setlist[enumerate]{leftmargin=1.7em,itemsep=5pt,topsep=4pt}
\setlist[itemize]{leftmargin=1.7em,itemsep=4pt,topsep=4pt,label={\color{poporange}\textbullet}}
\setlength{\parindent}{0pt}
\setlength{\parskip}{5pt}
\renewcommand{\arraystretch}{1.25}

% pgfplots : style Pop par défaut
\pgfplotsset{
  every axis/.append style={axis line style={popink,line width=1pt},tick style={popink},
    grid style={popline,line width=0.5pt},label style={font=\small},tick label style={font=\footnotesize}},
  popcurve/.style={poporange,line width=1.8pt,no markers,smooth},
}
% Même style disponible dans un \draw TikZ simple (hors pgfplots)
\tikzset{popcurve/.style={poporange,line width=1.8pt}}
\tikzset{popgrid/.style={popline,very thin}}
\colorlet{popcurve}{poporange} % le nom du style est parfois utilisé comme couleur

% ============================================================
% Pill / squiggle
% ============================================================
\newcommand{\poppill}[3]{%
  \tikz[baseline=-0.55ex]\node[draw=popink,line width=1.2pt,rounded corners=2.6pt,%
    fill=#2,inner xsep=6.5pt,inner ysep=3pt]{%
    \popxboldls\fontsize{7}{7}\selectfont\color{#3}\MakeUppercase{#1}};%
}
\newcommand{\popsquiggle}[1]{%
  \tikz[baseline=-0.4ex]{
    \draw[#1,line width=2.6pt,line cap=round]
      plot [smooth] coordinates {(0,0.09) (0.34,0.01) (0.68,0.21) (1.02,0.01) (1.36,0.10)};
  }%
}
% Mot-clé surligné (marqueur orange sous le texte)
\newcommand{\cle}[1]{{\popxbold\color{popdark}#1}}

% ============================================================
% En-tête / pied
% ============================================================
\pagestyle{fancy}
\fancyhf{}
\renewcommand{\headrulewidth}{0pt}
\renewcommand{\footrulewidth}{0pt}
\lhead{\raisebox{-0.15ex}{\poplogo\fontsize{13}{13}\selectfont\color{popink}OnBuch\color{poporange}+}}
\rhead{\poppill{\DOCMATIERE\ \textbullet\ \DOCNIVEAU\ \textbullet\ Leçon \DOCLECON}{white}{popink}}
\lfoot{\popxbold\fontsize{7.6}{7.6}\selectfont\color{popmuted}\MakeUppercase{Cours signé OnBuch+}\ \textbullet\ {\popsemi\DOCTITRE}}
\rfoot{\tikz[baseline=-0.6ex]\node[rounded corners=8.5pt,fill=popink,minimum height=17pt,minimum width=17pt,inner xsep=5pt,inner sep=0]{\popxbold\fontsize{8}{8}\selectfont\color{popcream}\thepage\,/\,\pageref*{LastPage}};}

% ============================================================
% Titres
% ============================================================
\newcommand{\chaptitre}[2]{%
  \begin{tcolorbox}[enhanced,colback=poporange,colframe=popink,boxrule=2.6pt,arc=11pt,
      drop shadow=popink,left=15pt,right=15pt,top=12pt,bottom=11pt,before skip=3pt,after skip=18pt]
    \poppill{#2}{popink}{popcream}\par\vspace{9pt}
    {\raggedright\hyphenpenalty=10000\exhyphenpenalty=10000\popdisplay\fontsize{22}{24}\selectfont\color{popcream}\MakeUppercase{#1}\par}
  \end{tcolorbox}
}
% Section numérotée : pastille orange avec le numéro + titre Archivo
\newcounter{popsec}
\newcommand{\coursec}[1]{%
  \stepcounter{popsec}\setcounter{popsub}{0}%
  \par\vspace{16pt}\noindent
  \begin{minipage}{\linewidth}
  \tikz[baseline=-0.7ex]\node[draw=popink,line width=1.6pt,fill=poporange,rounded corners=4pt,minimum size=22pt,inner sep=0]{\popdisplay\fontsize{12}{12}\selectfont\color{popcream}\thepopsec};%
  \hspace{8pt}{\popdisplay\fontsize{15}{17}\selectfont\color{popink}\MakeUppercase{#1}}\par
  \vspace{4pt}\noindent{\color{poporange}\rule{36pt}{3.6pt}}
  \end{minipage}\par\nopagebreak\vspace{8pt}
}
\newcounter{popsub}
\newcommand{\courssub}[1]{%
  \stepcounter{popsub}%
  \par\vspace{9pt}\noindent{\popxbold\fontsize{11.5}{13}\selectfont\color{popink}{\color{poporange}\thepopsec.\thepopsub}\hspace{6pt}#1}\par\nopagebreak\vspace{4pt}
}

% ============================================================
% Boîtes pédagogiques — même recette : fond blanc, bordure épaisse colorée,
% ombre dure encre, badge plein. Titre optionnel [#1].
% ============================================================
\tcbset{popbase/.style={breakable,enhanced,colback=white,boxrule=2.3pt,arc=9pt,
  shadow={3pt}{-3pt}{0pt}{popink},left=14pt,right=14pt,top=11pt,bottom=11pt,before skip=11pt,after skip=15pt,
  fontupper=\popsemi\fontsize{10.6}{14.5}\selectfont\color{popink},
  fontlower=\popsemi\fontsize{10.6}{14.5}\selectfont\color{popink}}}
% Coche dessinée (le glyphe ✓ n'existe pas dans les polices du document)
\newcommand{\popcheck}[1]{\tikz[baseline=-0.5ex]\draw[#1,line width=1.6pt,line cap=round,line join=round](-0.13,0.0)--(-0.04,-0.09)--(0.13,0.1);}
\providecommand{\coche}{\popcheck{popgreen}}
\providecommand{\resultat}[1]{\par\smallskip\noindent\fcolorbox{popgreen}{popgreenL}{\parbox{\dimexpr\linewidth-2\fboxsep-2\fboxrule\relax}{\textbf{Résultat.} #1}}\par\smallskip}
\newcommand{\popboxhead}[3]{\poppill{#1}{#2}{white}\ifx\relax#3\relax\else\hspace{6pt}{\popxbold\fontsize{10.6}{12}\selectfont\color{popink}#3}\fi\par\vspace{7pt}}

\newtcolorbox{aretenir}[1][]{popbase,colframe=popgreen,before upper={\popboxhead{À retenir}{popgreen}{#1}}}
% Texte en anglais (document de lecture, dialogue, extrait) : cadre
% d'étude. Un titre optionnel (source, auteur) s'affiche dans l'en-tête.
\newtcolorbox{textebox}[1][]{popbase,colframe=popblue,before upper={\popboxhead{Text}{popblue}{#1}\begin{otherlanguage}{english}},after upper={\end{otherlanguage}}}
% Marques ✓ / ✗ (forme correcte / fautive) souvent utilisées par les modèles
\providecommand{\cross}{\textcolor{poppink}{\ensuremath{\times}}}
\providecommand{\xmark}{\cross}\providecommand{\crossmark}{\cross}\providecommand{\cmark}{\ensuremath{\checkmark}}
% Ligne de réponse / justification à compléter (inventée par les modèles)
\providecommand{\justification}[1]{\par\noindent\textit{Justification~:} \dotfill\par}
% \textipa (package tipa non chargé) : la police ne couvre pas l'API -> simple passage
\providecommand{\textipa}[1]{#1}
% Soulignement (ulem non chargé) : \uline, \ul
\providecommand{\uline}[1]{\underline{#1}}\providecommand{\ul}[1]{\underline{#1}}
% \glossaire{\item ...} inventée par les modèles : liste de glossaire
\providecommand{\glossaire}[1]{\begin{itemize}#1\end{itemize}}
% \alert (beamer) inventée par les modèles : texte mis en évidence
\providecommand{\alert}[1]{\textbf{\textcolor{poppink}{#1}}}
% variantes ASCII d'ordinaux écrites en macro
\providecommand{\iere}{\textsuperscript{re}}\providecommand{\ieme}{\textsuperscript{e}}\providecommand{\ere}{\textsuperscript{re}}\providecommand{\eme}{\textsuperscript{e}}\providecommand{\er}{\textsuperscript{er}}\providecommand{\ie}{\textsuperscript{e}}
% Ordinaux français écrits en macro par les modèles : 1\ière, 3\ième
\newcommand{\ière}{\textsuperscript{re}}\newcommand{\ième}{\textsuperscript{e}}
% \exemple (inventée par les modèles) : étiquette « Exemple : »
\providecommand{\exemple}{\textbf{Exemple~:}\ }
% \square utilisable aussi hors mode math (cases à cocher)
\AtBeginDocument{\let\squareMath\square\renewcommand{\square}{\ensuremath{\squareMath}}}
% Mot ou phrase en anglais à l'intérieur d'un paragraphe français
\newcommand{\eng}[1]{\ifmmode\text{\textit{#1}}\else\begingroup\selectlanguage{english}\itshape #1\endgroup\fi}
\let\esp\eng % alias toléré
% flèches d'intonation inventées par les modèles
\providecommand{\textsearrow}{}\renewcommand{\textsearrow}{\ensuremath{\searrow}}\providecommand{\textnearrow}{}\renewcommand{\textnearrow}{\ensuremath{\nearrow}}
% \fr{..} : mot français (inventé par les modèles)
\providecommand{\fr}[1]{\textit{#1}}
\newtcolorbox{exemplebox}[1][]{popbase,colframe=poporange,before upper={\popboxhead{Exemple}{poporange}{#1}}}
\newtcolorbox{definition}[1][]{popbase,colframe=popblue,before upper={\popboxhead{Définition}{popblue}{#1}}}
\newtcolorbox{propriete}[1][]{popbase,colframe=poppurple,before upper={\popboxhead{Propriété}{poppurple}{#1}}}
\newtcolorbox{methode}[1][]{popbase,colframe=popgold,before upper={\popboxhead{Méthode}{popgold}{#1}}}
\newtcolorbox{attention}[1][]{popbase,colframe=poppink,before upper={\popboxhead{Attention}{poppink}{#1}}}
\newtcolorbox{experience}[1][]{popbase,colframe=popblue,colback=popblueL,before upper={\popboxhead{Expérience}{popblue}{#1}}}
% « Le savais-tu ? » : carte violette pleine (contexte, culture, Cameroun)
\newtcolorbox{savaistu}[1][]{popbase,colback=poppurple,colframe=popink,
  fontupper=\popsemi\fontsize{10.4}{14.2}\selectfont\color{white},
  before upper={\poppill{Le savais-tu\,?}{white}{poppurple}\ifx\relax#1\relax\else\hspace{6pt}{\popxbold\color{white}#1}\fi\par\vspace{7pt}}}
% Exercice résolu : énoncé en haut, \tcblower puis la solution sur fond crème
\newcounter{popexo}\newcounter{popexr}
\newtcolorbox{exoresolu}[1][]{popbase,colframe=poporange,colbacklower=popcream,
  segmentation style={popink,line width=1.2pt,dashed},
  before upper={\stepcounter{popexr}\popboxhead{Exercice résolu \thepopexr}{poporange}{#1}},
  before lower={\poppill{Solution}{popink}{popcream}\par\vspace{6pt}}}
% Exercice (non résolu en place) — la correction est dans la partie Corrigés
\newtcolorbox{exercice}[1][]{popbase,colframe=popink,
  before upper={\stepcounter{popexo}\popboxhead{Exercice \thepopexo}{popink}{#1}}}
% Corrigé
\newtcolorbox{corrige}[1][]{popbase,colframe=popgreen,colback=popgreenL,
  before upper={\popboxhead{Corrigé}{popgreen}{#1}}}
% Objectifs / prérequis / bilan
\newtcolorbox{objectifs}{popbase,colframe=popink,colback=poporangeL,
  before upper={\popboxhead{Objectifs de la leçon}{popink}{}}}
\newtcolorbox{prerequis}{popbase,colframe=popink,colback=popblueL,
  before upper={\popboxhead{Prérequis}{popblue}{}}}
\newtcolorbox{bilan}{popbase,colframe=popink,colback=white,boxrule=2.8pt,
  before upper={\popboxhead{Fiche bilan}{poporange}{L'essentiel à connaître pour l'examen}}}
% Figure encadrée avec légende
\newtcolorbox{popfigure}[1][]{enhanced,breakable=false,colback=white,colframe=popline,boxrule=1.4pt,arc=8pt,
  left=10pt,right=10pt,top=10pt,bottom=8pt,before skip=10pt,after skip=12pt,halign=center,
  lower separated=false,halign lower=center,
  fontlower=\popsemi\fontsize{9}{11.5}\selectfont\color{popmuted},
  #1}
\newcommand{\legende}[1]{\par\vspace{6pt}{\popsemi\fontsize{9}{11.5}\selectfont\color{popmuted}#1}}

% ============================================================
% Signature OnBuch+
% ============================================================
\newcommand{\poplogoinline}[1]{{\poplogo\fontsize{#1}{#1}\selectfont\color{popink}OnBuch\color{poporange}+}}
\newcommand{\signatureonbuch}{%
  \begin{tcolorbox}[enhanced,colback=popink,colframe=popink,boxrule=2.2pt,arc=10pt,
      drop shadow=poporange,left=14pt,right=14pt,top=10pt,bottom=10pt,before skip=0pt,after skip=0pt]
    \tikz[baseline=-0.7ex]\node[circle,fill=popgreen,draw=popcream,line width=1.3pt,minimum size=22pt,inner sep=0]{\popcheck{white}};%
    \hspace{9pt}\begin{minipage}[c]{0.62\linewidth}
      {\popxbold\fontsize{12}{13}\selectfont\color{popcream}Signé\ \textbullet\ {\poplogo OnBuch\color{poporange}+}}\par\vspace{2pt}
      {\raggedright\popsemi\fontsize{8}{10.5}\selectfont\color{popline}\MakeUppercase{Équipe pédagogique OnBuch+}\\\MakeUppercase{Conforme au programme officiel MINESEC}\par}
    \end{minipage}\hfill
    {\poplogo\fontsize{22}{22}\selectfont\color{popcream}OnBuch\color{poporange}+}
  \end{tcolorbox}%
}

% ============================================================
% Page de garde
% #1 titre · #2 matière · #3 niveau · #4 module · #5 n° leçon · #6 durée ·
% #7 description / objectifs (texte ou itemize)
% ============================================================
\newcommand{\pagegarde}[7]{%
  \thispagestyle{empty}
  \newgeometry{a4paper,margin=1.7cm,top=1.6cm,bottom=1.4cm}%
  \begin{tikzpicture}[remember picture,overlay]
    \fill[poporange!18] ([xshift=-3.2cm,yshift=-3.6cm]current page.north east) circle (4.6cm);
    \fill[poppurple!14] ([xshift=2.4cm,yshift=7.6cm]current page.south west) circle (3.8cm);
  \end{tikzpicture}%
  {\poplogo\fontsize{30}{30}\selectfont\color{popink}OnBuch\color{poporange}+}\hfill
  \raisebox{0.6ex}{\poppill{Cours signé \textbullet\ Premium}{popink}{popcream}}\par
  \vspace{1pt}\popsquiggle{poppurple}\par
  \vspace{8pt}
  \begin{tcolorbox}[enhanced,colback=poporange,colframe=popink,boxrule=2.8pt,arc=13pt,
      drop shadow=popink,left=18pt,right=18pt,top=12pt,bottom=13pt,after skip=10pt]
    \poppill{#2 \textbullet\ #3}{popink}{popcream}\hspace{5pt}\poppill{#4}{white}{popink}\par\vspace{8pt}
    {\popdisplay\fontsize{14}{14}\selectfont\color{popink}LEÇON #5}\par\vspace{4pt}
    {\raggedright\hyphenpenalty=10000\exhyphenpenalty=10000\popdisplay\fontsize{25}{27}\selectfont\color{popcream}\MakeUppercase{#1}\par}
    \vspace{8pt}
    {\popxbold\fontsize{9.5}{9.5}\selectfont\color{popink}\MakeUppercase{Durée conseillée : #6}}
  \end{tcolorbox}
  \begin{tcolorbox}[enhanced,colback=white,colframe=popink,boxrule=2.2pt,arc=10pt,
      drop shadow=popink,left=16pt,right=16pt,top=10pt,bottom=10pt,after skip=8pt]
    \poppill{Dans cette leçon}{poppurple}{white}\par\vspace{5pt}
    \popsemi\fontsize{9.3}{12.6}\selectfont\color{popink}#7
  \end{tcolorbox}
  \noindent\begin{minipage}[t]{0.487\linewidth}
    \begin{tcolorbox}[enhanced,colback=popgreen,colframe=popink,boxrule=2.2pt,arc=10pt,
        drop shadow=popink,left=11pt,right=11pt,top=10pt,bottom=10pt,height=3.1cm,valign=center]
      \tcbox[colback=white,colframe=popink,boxrule=1.3pt,arc=5pt,boxsep=0pt,nobeforeafter,
        left=3pt,right=3pt,top=3pt,bottom=3pt,valign=center]{\qrcode[height=1.9cm]{https://play.google.com/store/apps/details?id=com.luvvix.onbuch&hl=fr}}%
      \hspace{8pt}\begin{minipage}[c]{0.5\linewidth}
        {\popdisplay\fontsize{11}{12.5}\selectfont\color{white}\MakeUppercase{Télécharge OnBuch}}\par\vspace{4pt}
        {\raggedright\popsemi\fontsize{8.4}{11}\selectfont\color{white}Gratuit \textbullet\ Google Play\\Tuteur IA, annales, concours, résultats\par}
      \end{minipage}
    \end{tcolorbox}
  \end{minipage}\hfill
  \begin{minipage}[t]{0.487\linewidth}
    \begin{tcolorbox}[enhanced,colback=poppurple,colframe=popink,boxrule=2.2pt,arc=10pt,
        drop shadow=popink,left=11pt,right=11pt,top=10pt,bottom=10pt,height=3.1cm,valign=center]
      \tikz[baseline=-0.7ex]\node[circle,fill=white,draw=popink,line width=1.3pt,minimum size=1.6cm,inner sep=0]{\popdisplay\fontsize{24}{24}\selectfont\color{poppurple}+};%
      \hspace{8pt}\begin{minipage}[c]{0.6\linewidth}
        {\popdisplay\fontsize{11}{12.5}\selectfont\color{white}\MakeUppercase{Passe à OnBuch+}}\par\vspace{4pt}
        {\raggedright\popsemi\fontsize{8.4}{11}\selectfont\color{white}Tous les cours signés, Tuteur IA illimité, fiches PDF — onglet OnBuch+ de l'app\par}
      \end{minipage}
    \end{tcolorbox}
  \end{minipage}
  \vspace{8pt}
  \popcontenu
  \vfill
  \signatureonbuch
  \restoregeometry
  \newpage
}

% Rangée « Ce cours contient » (page de garde)
\newcommand{\poptile}[3]{% #1 couleur #2 gros texte #3 légende
  \begin{tcolorbox}[enhanced,colback=#1,colframe=popink,boxrule=2pt,arc=9pt,shadow={2.5pt}{-2.5pt}{0pt}{popink},
      left=6pt,right=6pt,top=9pt,bottom=9pt,halign=center,valign=center,height=1.75cm,width=\linewidth,nobeforeafter]
    {\popdisplay\fontsize{12.5}{13}\selectfont\color{white}\MakeUppercase{#2}}\par\vspace{3pt}
    {\centering\popsemi\fontsize{7.8}{9.5}\selectfont\color{white}#3\par}
  \end{tcolorbox}}
\newcommand{\popcontenu}{%
  \par\noindent{\popxboldls\fontsize{7.5}{7.5}\selectfont\color{popmuted}CE COURS SIGNÉ CONTIENT}\par\vspace{7pt}
  \noindent\begin{minipage}[t]{0.235\linewidth}\poptile{popblue}{Cours}{complet, illustré, pas à pas}\end{minipage}\hfill
  \begin{minipage}[t]{0.235\linewidth}\poptile{poporange}{Méthodes}{et exercices résolus}\end{minipage}\hfill
  \begin{minipage}[t]{0.235\linewidth}\poptile{poppink}{Exercices}{type examen + corrigés}\end{minipage}\hfill
  \begin{minipage}[t]{0.235\linewidth}\poptile{popgreen}{Fiche bilan}{l'essentiel à retenir}\end{minipage}\par}

% Sommaire
\newtcolorbox{sommaire}{enhanced,colback=white,colframe=popink,boxrule=2.2pt,arc=10pt,
  drop shadow=popink,left=18pt,right=18pt,top=13pt,bottom=13pt,before skip=6pt,after skip=20pt,
  before upper={\poppill{Sommaire}{poppurple}{white}\par\vspace{9pt}\popsemi\fontsize{10.6}{14.5}\selectfont\color{popink}}}

% Signature de fin de cours — poussée tout en bas de la dernière page
\newcommand{\findecours}{%
  \par\vfill\nopagebreak
  \begin{center}\popsquiggle{poporange}\end{center}\vspace{4pt}
  \signatureonbuch
}
\AtBeginDocument{\def\llbracket{\mathopen{[\mkern-3.5mu[}}\def\rrbracket{\mathclose{]\mkern-3.5mu]}}} % intervalles d'entiers (glyphe ⟦ absent de la police)
% Glyphes absents de Fira Math : redéfinis à partir de glyphes disponibles
\AtBeginDocument{%
  \def\setminus{\mathbin{\backslash}}\let\smallsetminus\setminus
  \def\vdots{\mathord{\vbox{\baselineskip3.2pt\lineskiplimit0pt\kern4pt\hbox{.}\hbox{.}\hbox{.}}}}%
  \def\ddots{\mathinner{\mkern1mu\raise6.5pt\hbox{.}\mkern2mu\raise3.6pt\hbox{.}\mkern2mu\raise0.7pt\hbox{.}\mkern1mu}}%
  \def\triangle{\mathord{\scalebox{0.9}{$\symup{\Delta}$}}}%
  \def\nabla{\mathord{\raisebox{\depth}{\scalebox{1}[-1]{$\symup{\Delta}$}}}}%
  \def\checkmark{\popcheck{popdark}}%
  \def\star{\mathbin{\ast}}\def\bigstar{\mathord{\ast}}%
  \def\diamond{\mathbin{\scalebox{0.6}{$\Diamond$}}}%
  \def\bigcup{\mathop{\mathchoice{\raisebox{-0.3em}{\scalebox{1.7}{$\cup$}}}{\scalebox{1.25}{$\cup$}}{\cup}{\cup}}\displaylimits}%
  \def\bigcap{\mathop{\mathchoice{\raisebox{-0.3em}{\scalebox{1.7}{$\cap$}}}{\scalebox{1.25}{$\cap$}}{\cap}{\cap}}\displaylimits}%
  \def\lesseqqgtr{\mathrel{\lessgtr}}\def\bigoplus{\mathop{\oplus}\displaylimits}%
}
\providecommand{\ssi}{si et seulement si} % abréviation fréquente
\AtBeginDocument{\providecommand{\wideparen}{\overparen}} % arc de cercle

\begin{document}

\pagegarde{Lesson 41 — Grammar: Phrasal verbs}{Anglais}{Tle A}{Chapter 4 : Citizenship and Human Rights}{EN41}{2 h}{Cette leçon de grammaire présente les phrasal verbs (verbes à particule) : leur formation, leurs types (séparables, inséparables, à trois mots) et leurs emplois, en lien avec le thème de la citoyenneté et des droits de l'homme. Elle insiste sur la comparaison avec le français et sur les erreurs fréquentes des apprenants francophones.
\vspace{4pt}
\begin{multicols}{2}\small
\begin{itemize}
\item À la fin de cette leçon, je sais définir ce qu'est un phrasal verb et le distinguer d'un verbe suivi d'une préposition ordinaire.
\item À la fin de cette leçon, je sais identifier le sens d'un phrasal verb courant dans un texte sur la citoyenneté.
\item À la fin de cette leçon, je sais classer un phrasal verb selon son type : inséparable, séparable, ou à trois mots.
\item À la fin de cette leçon, je sais placer correctement le complément (nom ou pronom) avec les phrasal verbs séparables.
\item À la fin de cette leçon, je sais conjuguer les phrasal verbs à tous les temps (present, past, -ing).
\item À la fin de cette leçon, je sais éviter les erreurs typiques des francophones (traduction littérale, mauvaise place du pronom, particule oubliée).
\end{itemize}\end{multicols}}

\begin{sommaire}
\begin{multicols}{2}
\begin{enumerate}
\item Observation et découverte
\item Définition et formation
\item Les trois types : séparables, inséparables, à trois mots
\item Phrasal verbs de la citoyenneté et des droits de l'homme
\item Comparaison français/anglais et pièges des francophones
\item Entraînement guidé et production
\item Activité d'intégration
\item Méthodes et exercices résolus
\item Exercices
\item Corrigés des exercices
\item Fiche bilan
\end{enumerate}
\end{multicols}
\end{sommaire}

\begin{objectifs}
\begin{itemize}
\item À la fin de cette leçon, je sais définir ce qu'est un phrasal verb et le distinguer d'un verbe suivi d'une préposition ordinaire.
\item À la fin de cette leçon, je sais identifier le sens d'un phrasal verb courant dans un texte sur la citoyenneté.
\item À la fin de cette leçon, je sais classer un phrasal verb selon son type : inséparable, séparable, ou à trois mots.
\item À la fin de cette leçon, je sais placer correctement le complément (nom ou pronom) avec les phrasal verbs séparables.
\item À la fin de cette leçon, je sais conjuguer les phrasal verbs à tous les temps (present, past, -ing).
\item À la fin de cette leçon, je sais éviter les erreurs typiques des francophones (traduction littérale, mauvaise place du pronom, particule oubliée).
\item À la fin de cette leçon, je sais utiliser au moins dix phrasal verbs liés aux droits de l'homme dans des phrases et un court paragraphe.
\item À la fin de cette leçon, je sais retrouver le sens d'un phrasal verb inconnu grâce au contexte.
\end{itemize}
\end{objectifs}

\begin{prerequis}
\begin{itemize}
\item Les verbes de base et leur conjugaison au present simple et au past simple (classe de Première).
\item La distinction verbe / préposition / adverbe (notions de grammaire de Seconde).
\item La place du complément d'objet direct dans la phrase anglaise (SVO).
\item Le vocabulaire de base de la citoyenneté : rights, vote, law, freedom (Lessons du Chapter 4).
\end{itemize}
\end{prerequis}

\newpage

\coursec{Observation et découverte}

\courssub{Une situation d'accroche : les élections à Yaoundé}

Imaginez la scène : nous sommes en période électorale à Yaoundé. Un journaliste de la CRTV, face caméra, commente l'ambiance dans un anglais fluide : \eng{“Citizens must stand up for their rights and speak out against injustice. If leaders fail, people vote them out. Activists never give in; they carry on the struggle and look after the victims.”} Un élève de Terminale A, assis dans sa classe du lycée Bilingue d'Etoug-Ebe, lève la main et demande : « Monsieur, pourquoi dit-on \eng{stand up for} et pas simplement \eng{defend} ? Pourquoi le sens change-t-il totalement quand on ajoute une petite particule comme \eng{up}, \eng{out} ou \eng{for} ? » C'est une excellente question. Ces verbes à particule, qu'on appelle \cle{phrasal verbs}, sont omniprésents dans les discours sur la citoyenneté, les droits de l'homme, la politique et les médias anglophones (BBC, CNN, The Guardian, mais aussi The Post ou Cameroon Tribune en version anglaise). Les maîtriser, c'est non seulement comprendre l'actualité internationale, mais aussi enrichir ses propres productions écrites et orales pour le Baccalauréat. Dans cette première section, nous allons observer, repérer et formuler nos premières hypothèses avant d'énoncer la règle rigoureuse.

\courssub{Repérage dans un texte authentique sur la citoyenneté}

Lisez attentivement le court texte ci-dessous, extrait d'un éditorial fictif d'un journal camerounais anglophone. Huit \cle{phrasal verbs} sont mis en \textbf{gras}. Ne cherchez pas encore la règle : observez simplement leur forme et essayez de deviner leur sens grâce au contexte.

\begin{textebox}[Editorial — “Active Citizenship in Cameroon”]
In a healthy democracy, citizens must \textbf{stand up for} their rights and \textbf{speak out} against corruption. When leaders \textbf{put off} necessary reforms, the people should not \textbf{give in} to apathy. Instead, they \textbf{carry on} peaceful protests and \textbf{look after} the most vulnerable. At the polls, voters \textbf{turn up} in large numbers to \textbf{vote out} incompetent officials. Civic duty is not a sprint; it is a marathon.
\end{textebox}

\begin{center}
\begin{tabularx}{\linewidth}{|X|X|X|}
\hline
\rowcolor{popblueL} \textbf{English word / phrase} & \textbf{Nature} & \textbf{Français} \\
\hline
stand up for & phrasal verb (3 mots) & défendre, soutenir activement \\
\hline
speak out & phrasal verb (2 mots) & s'exprimer publiquement, élever la voix \\
\hline
put off & phrasal verb (2 mots) & reporter, remettre à plus tard \\
\hline
give in & phrasal verb (2 mots) & céder, abandonner la résistance \\
\hline
carry on & phrasal verb (2 mots) & continuer, poursuivre \\
\hline
look after & phrasal verb (2 mots) & s'occuper de, prendre soin de \\
\hline
turn up & phrasal verb (2 mots) & se présenter, arriver, venir \\
\hline
vote out & phrasal verb (2 mots) & chasser par le vote, ne pas réélire \\
\hline
apathy & noun (n.) & indifférence, apathie \\
\hline
polls & noun (n. pl.) & urnes, élections \\
\hline
vulnerable & adjective (adj.) & vulnérable, fragile \\
\hline
incompetent & adjective (adj.) & incompétent \\
\hline
\end{tabularx}
\end{center}

\courssub{Analyse guidée : la structure « verbe + particule »}

Regardons de plus près la forme de ces huit verbes. Chacun est composé de deux (parfois trois) éléments :
\begin{enumerate}
    \item Un \textbf{verbe lexical} (verbe d'action ou d'état) : \eng{stand, speak, put, give, carry, look, turn, vote}.
    \item Une \textbf{particule} (souvent un adverbe ou une préposition utilisée comme adverbe) : \eng{up, out, off, in, on, after, for}.
\end{enumerate}

\begin{definition}[Phrasal verb — Structure de base]
Un \textbf{phrasal verb} (verbe à particule) est une unité lexicale formée d'un \textbf{verbe} + une \textbf{particule} (adverbe ou préposition). Le sens de l'ensemble est souvent \textbf{idiomatique} : il ne se déduit pas automatiquement de la simple addition du sens du verbe et du sens de la particule.
\end{definition}

\begin{exemplebox}[Analyse morphologique]
\begin{itemize}
    \item \eng{stand up for} = \eng{stand} (verbe) + \eng{up} (particule) + \eng{for} (particule/préposition) $\rightarrow$ 3 mots.
    \item \eng{speak out} = \eng{speak} (verbe) + \eng{out} (particule) $\rightarrow$ 2 mots.
    \item \eng{look after} = \eng{look} (verbe) + \eng{after} (particule) $\rightarrow$ 2 mots.
\end{itemize}
\end{exemplebox}

\courssub{Le sens global : pas une simple addition (le piège de la traduction mot à mot)}

C'est le point crucial pour un francophone. En français, on utilise souvent un verbe simple unique là où l'anglais utilise un verbe + particule. Le sens du \cle{phrasal verb} est \textbf{global} (unitaire), pas \textbf{compositionnel} (additionnel).

\begin{attention}[Ne pas traduire mot à mot !]
Ne dites jamais : \eng{“give in” = “donner dans”} ou \eng{“give up” = “donner en haut”}.
\begin{itemize}
    \item \eng{Give} = donner.
    \item \eng{Give in} = \textbf{céder} (sous la pression).
    \item \eng{Give up} = \textbf{abandonner} (un effort, une habitude).
    \item \eng{Give away} = \textbf{donner gratuitement} / \textbf{révéler} (un secret).
    \item \eng{Give back} = \textbf{rendre}.
\end{itemize}
La particule change \textbf{radicalement} le sens du verbe de base. C'est un \textbf{nouveau mot} à apprendre comme un vocabulaire à part entière.
\end{attention}

\begin{exemplebox}[Exemples traduits — Sens global vs sens littéral]
\eng{Citizens stand up for their rights.} « Les citoyens défendent leurs droits. » (Pas : « se lèvent vers le haut pour ».)
\eng{They never give in.} « Ils ne cèdent jamais. » (Pas : « donnent dedans ».)
\eng{People vote the government out.} « Les gens chassent le gouvernement par le vote. » (Pas : « votent le gouvernement dehors ».)
\eng{Activists carry on the struggle.} « Les militants poursuivent la lutte. » (Pas : « portent sur ».)
\eng{We look after the victims.} « Nous nous occupons des victimes. » (Pas : « regardons après ».)
\end{exemplebox}

\courssub{Phrasal verb ou verbe + préposition ordinaire ? La distinction cruciale}

Tous les couples \textit{verbe + petit mot} ne sont pas des \cle{phrasal verbs}. Il faut distinguer le \textbf{phrasal verb} (sens idiomatique, nouveau) du \textbf{verbe prépositionnel} (sens littéral, compositionnel).

\begin{propriete}[Distinction : Phrasal verb vs Verbe + Préposition littérale]
\begin{itemize}
    \item \textbf{Phrasal verb (sens nouveau / idiomatique) :} La particule modifie le sens du verbe. L'ensemble forme une unité sémantique.
    \item \textbf{Verbe + préposition (sens littéral / spatial) :} La préposition introduit un complément de lieu, de temps, etc. Le verbe garde son sens de base.
\end{itemize}
\end{propriete}

\begin{exemplebox}[Le test du sens : \eng{look after} vs \eng{look at}]
\begin{itemize}
    \item \eng{“We look \textbf{after} the children.”} $\rightarrow$ \textbf{Phrasal verb}. Sens : « Nous nous occupons des enfants. » (\eng{after} ne marque pas un lieu « derrière », mais change le sens de \eng{look}).
    \item \eng{“Look \textbf{at} that picture!”} $\rightarrow$ \textbf{Verbe + préposition littérale}. Sens : « Regarde ce tableau ! » (\eng{at} indique la cible du regard, \eng{look} garde son sens de « diriger les yeux vers »).
    \item \eng{“He ran \textbf{into} an old friend.”} $\rightarrow$ Phrasal verb (\eng{run into} = rencontrer par hasard).
    \item \eng{“He ran \textbf{into} the house.”} $\rightarrow$ Verbe + préposition littérale (\eng{into} = direction/lieu).
\end{itemize}
\end{exemplebox}

\courssub{Première hypothèse de règle (formulée par les élèves)}

Avant la leçon magistrale, essayez de formuler une règle générale à partir de vos observations. Complétez les phrases ci-dessous :

\begin{experience}[Activité de classe : Formulation d'hypothèses]
En binôme, discutez et écrivez vos réponses au tableau :
\begin{enumerate}
    \item Un phrasal verb est composé de \underline{\hspace{5cm}}.
    \item Le sens d'un phrasal verb est \underline{\hspace{5cm}} (additionnel / global / littéral).
    \item Pour savoir si \eng{look after} est un phrasal verb, on teste \underline{\hspace{5cm}}.
    \item En français, on traduit souvent un phrasal verb anglais par \underline{\hspace{5cm}}.
\end{enumerate}
\textbf{Mise en commun attendue :}
1. Un verbe + une particule (adverbe/préposition). 2. Global (idiomatique). 3. On vérifie si le sens est différent du verbe seul (test de substitution : \eng{look after} $\neq$ \eng{watch} + lieu). 4. Un verbe simple (ex. \eng{defend, continue, surrender, care for}).
\end{experience}

\courssub{Visualisation : Carte mentale de la structure}

La figure ci-dessous résume la composition interne d'un \cle{phrasal verb}. Retenez les deux branches principales : le \textbf{verbe} (porteur de l'action de base) et la \textbf{particule} (modificateur de sens).

\begin{popfigure}
\begin{tikzpicture}[
    mindmap,
    concept color=popblue,
    level 1 concept/.append style={level distance=3.5cm, sibling angle=120},
    level 2 concept/.append style={level distance=2.5cm},
    every node/.style={font=\sffamily\small},
    text=popdark
]
\node[concept, scale=1.2] {PHRASAL VERB}
    child[concept color=poporange] {
        node[concept] {VERBE}
        child { node[rectangle, draw=poporange, fill=poporangeL, rounded corners, inner sep=3pt] {stand} }
        child { node[rectangle, draw=poporange, fill=poporangeL, rounded corners, inner sep=3pt] {give} }
        child { node[rectangle, draw=poporange, fill=poporangeL, rounded corners, inner sep=3pt] {look} }
        child { node[rectangle, draw=poporange, fill=poporangeL, rounded corners, inner sep=3pt] {carry} }
        child { node[rectangle, draw=poporange, fill=poporangeL, rounded corners, inner sep=3pt] {vote} }
    }
    child[concept color=popgreen] {
        node[concept] {PARTICULE}
        child { node[rectangle, draw=popgreen, fill=popgreenL, rounded corners, inner sep=3pt] {up} }
        child { node[rectangle, draw=popgreen, fill=popgreenL, rounded corners, inner sep=3pt] {out} }
        child { node[rectangle, draw=popgreen, fill=popgreenL, rounded corners, inner sep=3pt] {after} }
        child { node[rectangle, draw=popgreen, fill=popgreenL, rounded corners, inner sep=3pt] {on} }
        child { node[rectangle, draw=popgreen, fill=popgreenL, rounded corners, inner sep=3pt] {for} }
    }
    child[concept color=poppurple] {
        node[concept] {SENS GLOBAL}
        child { node[rectangle, draw=poppurple, fill=poppurpleL, rounded corners, inner sep=3pt] {Idiomatique} }
        child { node[rectangle, draw=poppurple, fill=poppurpleL, rounded corners, inner sep=3pt] {Non compositionnel} }
        child { node[rectangle, draw=poppurple, fill=poppurpleL, rounded corners, inner sep=3pt] {= 1 mot en français} }
    };
\end{tikzpicture}
\legende{Structure interne d'un phrasal verb : Verbe + Particule $\rightarrow$ Sens global (idiomatique).}
\end{popfigure}

\courssub{Le saviez-vous ? La fréquence massive des phrasal verbs}

\begin{savaistu}[L'anglais des médias et du quotidien]
L'anglais courant (spoken English), la presse (The Guardian, BBC News) et la littérature contemporaine utilisent des \textbf{dizaines de milliers} de phrasal verbs. Une étude de corpus (COCA - Corpus of Contemporary American English) montre que les \textbf{20 phrasal verbs les plus fréquents} (\eng{go on, pick up, come back, go out, come up, find out, come out, go back, sit down, get up, take off, give up, wake up, put on, turn on, turn off, look for, run out, show up, work out}) couvrent à eux seuls une part énorme des occurrences dans la conversation quotidienne. Au Bac, les correcteurs valorisent l'utilisation naturelle de phrasal verbs adaptés au registre (ex. \eng{point out} au lieu de \eng{show} dans un essai, \eng{bring up} au lieu de \eng{raise} un sujet). Ne les fuyez pas : \textbf{apprenez-les par blocs thématiques} (citoyenneté, environnement, éducation, santé).
\end{savaistu}

\courssub{Entraînement guidé : Identification et traduction}

Appliquons immédiatement notre observation. L'exercice résolu ci-dessous sert de modèle pour votre travail personnel.

\begin{exoresolu}[Repérer, analyser et traduire des phrasal verbs]
\textbf{Énoncé :} Dans les phrases suivantes, issues de discours sur les droits de l'homme :
\begin{enumerate}
    \item Soulignez le phrasal verb.
    \item Séparez le verbe et la/les particule(s) par un trait vertical $|$.
    \item Donnez le sens en français (un verbe simple ou une locution).
    \item Dites si le sens est \textbf{littéral} (L) ou \textbf{idiomatique} (I).
\end{enumerate}
\begin{enumerate}
    \item The government tried to \textbf{cover up} the scandal.
    \item Please \textbf{fill in} this application form.
    \item The children \textbf{looked for} their lost ball in the garden.
    \item We must \textbf{stand up} against discrimination.
    \item The meeting was \textbf{put off} until Monday.
\end{enumerate}

\tcblower

\textbf{Solution détaillée :}
\begin{enumerate}
    \item \textbf{cover $|$ up} — Sens : \textbf{camoufler, dissimuler} (I). \textit{Note :} \eng{cover} = couvrir ; \eng{cover up} = cacher la vérité.
    \item \textbf{fill $|$ in} — Sens : \textbf{remplir (un formulaire)} (I). \textit{Note :} \eng{fill} = remplir (un verre) ; \eng{fill in/out} = compléter un document.
    \item \textbf{looked $|$ for} — Sens : \textbf{chercher} (I). \textit{Attention :} \eng{look at} = regarder (L) ; \eng{look for} = chercher (I). Ici, sens idiomatique.
    \item \textbf{stand $|$ up} — Sens : \textbf{se dresser, s'opposer fermement} (I). \textit{Contexte :} \eng{stand up against} = s'opposer à. Différent de \eng{stand up} (se lever / L).
    \item \textbf{put $|$ off} — Sens : \textbf{reporter, ajourner} (I). \textit{Note :} \eng{put off} $\neq$ mettre hors tension (\eng{turn off}).
\end{enumerate}
\textbf{Bilan :} 5 sur 5 sont des sens idiomatiques (I). Le test de substitution par un verbe simple français confirme le statut de phrasal verb.
\end{exoresolu}

\courssub{Bilan de l'observation : Tableau récapitulatif}

Avant de passer à la définition formelle et aux types grammaticaux (séparables, inséparables), synthétisons nos découvertes dans le tableau comparatif suivant.

\begin{center}
\begin{tabularx}{\linewidth}{|X|X|X|}
\hline
\rowcolor{popblueL} \textbf{Observation} & \textbf{Constat en Anglais} & \textbf{Conséquence pour le Francophone} \\
\hline
\textbf{Forme} & Verbe + Particule (1 ou 2 particules) & Un seul mot anglais $\leftrightarrow$ souvent un seul mot français. \\
\hline
\textbf{Sens} & Global, idiomatique, imprévisible & \textbf{Interdiction de la traduction mot à mot.} Apprendre par cœur + contexte. \\
\hline
\textbf{Test} & Substitution par un verbe simple (\eng{defend, continue, surrender}) & Si remplaçable par 1 verbe français $\rightarrow$ Phrasal verb. \\
\hline
\textbf{Distinction} & \eng{look after} (s'occuper) vs \eng{look at} (regarder) & Vérifier si la particule change le sens du verbe (idiomatique) ou indique un lieu/temps (littéral). \\
\hline
\textbf{Fréquence} & Extrêmement élevé (oral, presse, Bac) & Atout majeur pour la note de \textbf{Expression Écrite} et \textbf{Compréhension}. \\
\hline
\end{tabularx}
\end{center}

\begin{aretenir}[Ce qu'il faut retenir de la Section 1]
\begin{itemize}
    \item Un \cle{phrasal verb} = \textbf{Verbe} + \textbf{Particule(s)}. C'est une unité lexicale unique.
    \item Son sens est \textbf{global (idiomatique)} : \eng{give in} $\neq$ \eng{give} + \eng{in}.
    \item On le distingue du \textbf{verbe + préposition littérale} par le test du sens (substitution par un verbe simple).
    \item En français, la traduction est presque toujours un \textbf{verbe simple} (\eng{stand up for} $\rightarrow$ défendre ; \eng{carry on} $\rightarrow$ continuer).
    \item Les 8 phrasal verbs du thème « Citoyenneté » : \eng{stand up for, speak out, give in, carry on, look after, vote out, put off, turn up}.
\end{itemize}
\end{aretenir}

\coursec{Observation et découverte}

\courssub{Textes supports : discours citoyen et actualité}

\begin{textebox}[Extraits de presse et discours — Cameroun et international]
\textbf{Text A (Editorial, \textit{The Cameroon Tribune})} \\
“Civil society organisations \textbf{call on} the government to \textbf{look into} the rising cost of living. They \textbf{stand up for} the rights of vulnerable populations and refuse to \textbf{give in} to empty promises.”

\textbf{Text B (Speech by a student leader, Buea)} \\
“Fellow students, we cannot \textbf{put up with} injustice any longer. We must \textbf{speak out} against corruption. If the administration \textbf{puts off} the elections again, we will \textbf{carry on} with our peaceful protest. We \textbf{stand for} transparency.”

\textbf{Text C (Radio debate, Yaoundé)} \\
“The Minister \textbf{stepped down} yesterday following the scandal. An independent commission will \textbf{carry out} an inquiry. Citizens hope the new leadership will \textbf{live up to} expectations and not \textbf{let down} the nation.”
\end{textebox}

\courssub{Activité d'observation guidée}

\begin{experience}[Repérage et hypothèses]
Travailler par binômes. À partir des trois textes ci-dessus :
\begin{enumerate}
    \item Surlignez tous les groupes \textbf{verbe + petit mot(s)} qui semblent former un sens unique (ex. : \eng{call on}, \textbf{pas} \eng{call} seul).
    \item Classez-les en deux colonnes : \textit{Sens transparent / déductible} vs \textit{Sens idiomatique / imprévisible}.
    \item Testez la mobilité de la particule : pouvez-vous dire \eng{call the government on} ? \eng{stand transparency for} ? \eng{put the elections off} ? Notez vos constatations.
\end{enumerate}
\textbf{Mise en commun :} Le professeur guide vers la distinction adverbe/préposition et la règle de séparabilité.
\end{experience}

\coursec{Définition et formation}

\courssub{Transition depuis l'observation}

Dans la section précédente, nous avons analysé des extraits de discours citoyens et des articles de presse où des verbes comme \eng{stand up for}, \eng{speak out}, \eng{carry on} ou \eng{put up with} apparaissaient naturellement. Vous avez remarqué que ces verbes sont formés d'un verbe de base suivi d'un petit mot (\eng{up}, \eng{out}, \eng{on}, \eng{for}, \eng{with}…) qui change radicalement le sens initial. Nous allons maintenant formaliser cette observation pour en faire un outil grammatical maîtrisé.

\courssub{Qu'est-ce qu'un \cle{phrasal verb} ?}

\begin{definition}[Définition du phrasal verb]
Un \textbf{phrasal verb} (verbe à particule) est une unité lexicale composée d'un \textbf{verbe} et d'une \textbf{particule} (adverbe ou préposition) — voire de deux particules — dont le sens global est souvent \textbf{différent} de la simple somme des sens du verbe et de la particule pris séparément. Il fonctionne comme un verbe unique à part entière.
\end{definition}

\begin{savaistu}[Origine et fréquence]
Les phrasal verbs sont une caractéristique majeure de l'anglais \textbf{germanique} (old English). Ils sont omniprésents dans l'anglais courant, journalistique et littéraire. L'anglais formel (juridique, académique) leur préfère souvent des verbes d'origine latine/française (\eng{give up} $\rightarrow$ \eng{abandon}, \eng{put off} $\rightarrow$ \eng{postpone}, \eng{look into} $\rightarrow$ \eng{investigate}). Au Cameroun, dans les médias anglophones (\textit{The Cameroon Tribune}, \textit{The Post}) ou les discours officiels, on les rencontre constamment.
\end{savaistu}

\begin{exemplebox}[Verbe seul vs Phrasal verb]
\begin{itemize}
    \item \eng{He \textbf{looks} the document.} $\rightarrow$ Il \textbf{regarde} le document. (sens littéral)
    \item \eng{He \textbf{looks after} the children.} $\rightarrow$ Il \textbf{s'occupe} des enfants. (sens idiomatique)
    \item \eng{The government \textbf{puts} the project \textbf{off}.} $\rightarrow$ Le gouvernement \textbf{reporte} le projet. (sens idiomatique, particule déplacée)
    \item \eng{She \textbf{stands up for} her rights.} $\rightarrow$ Elle \textbf{défend} ses droits. (verbe à trois mots)
\end{itemize}
\end{exemplebox}

\courssub{Formation : les trois structures internes}

On classe les phrasal verbs selon la nature grammaticale de la particule qui suit le verbe. Cette classification est essentielle pour maîtriser leur syntaxe (séparables ou non), traitée en Section 3.

\begin{propriete}[Les trois types de formation]
\begin{center}
\begin{tabularx}{\linewidth}{|X|X|X|}
\hline
\rowcolor{popblueL}
\textbf{Type} & \textbf{Structure} & \textbf{Exemples (Citoyenneté / Droits de l'homme)} \\
\hline
\textbf{1. Verbe + Adverbe} & V + \textit{up, down, in, out, on, off, away...} & \eng{speak out} (s'exprimer), \eng{step down} (démissionner), \eng{carry on} (continuer), \eng{give in} (céder) \\
\hline
\textbf{2. Verbe + Préposition} & V + \textit{for, after, into, on, with, against...} & \eng{stand for} (représenter/tolérer), \eng{look after} (s'occuper de), \eng{look into} (enquêter sur), \eng{stand against} (s'opposer à) \\
\hline
\textbf{3. Verbe + Adverbe + Préposition} & V + Adv + Prep & \eng{put up with} (supporter), \eng{look forward to} (avoir hâte de), \eng{stand up for} (défendre), \eng{get away with} (s'en tirer avec) \\
\hline
\end{tabularx}
\end{center}
\end{propriete}

\courssub{Séparabilité et ordre des mots : règles syntaxiques}

C'est le point crucial pour l'écrit et l'oral (Baccalauréat). La nature de la particule dicte la place de l'objet.

\begin{propriete}[Règles de position de l'objet (nom ou pronom)]
\begin{center}
\begin{tabularx}{\linewidth}{|X|X|X|}
\hline
\rowcolor{popblueL}
\textbf{Type} & \textbf{Objet NOMINAL} & \textbf{Objet PRONOMINAL (le, la, les, lui, y, en...)} \\
\hline
\textbf{Type 1 : V + Adverbe} (Séparable) & \textbf{Deux positions possibles} : \\
\eng{turn off the light} \textbf{OU} \eng{turn the light off} & \textbf{OBLIGATOIREMENT entre V et Adv} : \\
\eng{turn \textbf{it} off} \quad (\textit{jamais} \eng{turn off it}) \\
\hline
\textbf{Type 2 : V + Préposition} (Inséparable) & \textbf{Une seule position} : après la préposition. \\
\eng{look into the matter} \quad (\textit{jamais} \eng{look the matter into}) & \textbf{Après la préposition} : \\
\eng{look into \textbf{it}} \\
\hline
\textbf{Type 3 : V + Adv + Prep} (Inséparable) & \textbf{Une seule position} : après la 2e particule. \\
\eng{put up with the noise} \quad (\textit{jamais} \eng{put the noise up with}) & \textbf{Après la 2e particule} : \\
\eng{put up with \textbf{it}} \\
\hline
\end{tabularx}
\end{center}
\end{propriete}

\begin{attention}[Piège majeur : Objet long $\rightarrow$ particule avant l'objet (Type 1)]
Si l'objet nominal est \textbf{long} (groupe nominal complexe), on préfère ne pas le séparer du verbe pour éviter la lourdeur.
\begin{itemize}
    \item \textbf{Préféré :} \eng{The government \textbf{put off} the meeting scheduled for Monday.} (Particule avant l'objet long).
    \item \textbf{Évité :} \eng{The government \textbf{put} the meeting scheduled for Monday \textbf{off}.} (Trop loin du verbe).
\end{itemize}
\end{attention}

\begin{attention}[Particule vs Préposition : le test de l'objet]
\begin{itemize}
    \item Si la particule est un \textbf{adverbe} (Type 1), elle peut souvent être déplacée après l'objet : \eng{turn off the light} / \eng{turn the light off}.
    \item Si la particule est une \textbf{préposition} (Type 2), elle \textbf{ne peut jamais} être séparée de son objet : \eng{look after the children} (correct) — \eng{*look the children after} (incorrect).
    \item Les verbes à \textbf{trois mots} (Type 3) sont toujours inséparables : \eng{put up with the noise} — \eng{*put the noise up with}.
\end{itemize}
\end{attention}

\courssub{Schéma de structure et règle de conjugaison}

L'erreur n°1 des francophones est d'accorder la particule ou de la faire varier. La règle est absolue : \textbf{seul le verbe porte la morphologie} (temps, personne, nombre, voix, négation).

\begin{popfigure}
\begin{tikzpicture}[
    node distance=1.2cm and 0.8cm,
    box/.style={draw=popblue, thick, rounded corners, fill=popblueL, minimum width=2.5cm, minimum height=1cm, align=center, font=\sffamily\bfseries},
    arrow/.style={-Stealth, thick, popdark},
    labelbox/.style={draw=poporange, thick, rounded corners, fill=poporangeL, minimum width=3cm, minimum height=1.2cm, align=center, font=\sffamily\small}
]
% Nodes
\node[box, align=center] (verb){VERBE\normalsize\\ (give, stand, look)};
\node[box, right=of verb, align=center] (part){PARTICULE(s)\normalsize\\ (up, out, up with)};
\node[labelbox, below=of verb, align=center] (conj){SE CONJUGUE\\ (temps, personne,\\ négation, voix)};
\node[labelbox, below=of part, align=center] (invar){INVARIABLE\\ (jamais de -s,\\ -ed, -ing dessus)};
\node[box, right=of part, xshift=1.5cm, fill=popgreenL, draw=popgreen, align=center] (unit){UNITÉ DE SENS\normalsize\\ (phrasal verb)};

% Arrows
\draw[arrow] (verb) -- (part) node[midway, above, font=\small] {+};
\draw[arrow] (part) -- (unit) node[midway, above, font=\small] {$\rightarrow$};
\draw[arrow, poporange] (verb.south) -- (conj.north);
\draw[arrow, poporange] (part.south) -- (invar.north);
\draw[arrow, dashed, poppurple] (conj.east) -- (invar.west) node[midway, below, font=\tiny, align=center] {Règle :\\ seul V change};
\end{tikzpicture}
\legende{Structure morphologique du phrasal verb : le verbe supporte toute la grammaire, la particule reste figée.}
\end{popfigure}

\courssub{Conjugaison : le verbe seul bouge}

Quel que soit le type (1, 2 ou 3), la conjugaison suit le verbe de base. Prenons l'exemple type de la citoyenneté : \eng{\textbf{look into}} (enquêter sur — Type 2, inséparable), qui permet aisément la voix passive.

\begin{propriete}[Tableau de conjugaison : \eng{look into} (3e pers. sing. \& plur.)]
\begin{center}
\begin{tabular}{|l|l|l|l|}
\hline
\rowcolor{popblueL}
\textbf{Temps / Mode} & \textbf{3e pers. Singulier} & \textbf{3e pers. Pluriel} & \textbf{Traduction} \\
\hline
Present Simple & He \textbf{looks into} & They \textbf{look into} & Il enquête sur / Ils enquêtent sur \\
\hline
Present Continuous & He \textbf{is looking into} & They \textbf{are looking into} & Il est en train d'enquêter sur \\
\hline
Past Simple & He \textbf{looked into} & They \textbf{looked into} & Il a enquêté sur / Ils ont enquêté sur \\
\hline
Present Perfect & He \textbf{has looked into} & They \textbf{have looked into} & Il a enquêté sur (résultat) \\
\hline
Future (will) & He \textbf{will look into} & They \textbf{will look into} & Il enquêtera sur / Ils enquêteront sur \\
\hline
Future (be going to) & He \textbf{is going to look into} & They \textbf{are going to look into} & Il va enquêter sur \\
\hline
Modal (must/can) & He \textbf{must look into} & They \textbf{must look into} & Il doit enquêter sur \\
\hline
Passive Voice & It \textbf{is looked into} & They \textbf{are looked into} & C'est enquêté sur / Ils sont enquêtés sur \\
\hline
Négation & He \textbf{does not look into} & They \textbf{do not look into} & Il n'enquête pas sur \\
\hline
\end{tabular}
\end{center}
\end{propriete}

\begin{exemplebox}[Exemples conjugués en contexte camerounais]
\begin{itemize}
    \item \eng{The Minister \textbf{stepped down} yesterday.} $\rightarrow$ Le Ministre \textbf{a démissionné} hier. (Past simple, verbe irrégulier \eng{step/stepped/stepped}, Type 1).
    \item \eng{Journalists \textbf{are speaking out} against corruption.} $\rightarrow$ Les journalistes \textbf{s'expriment / dénoncent} la corruption. (Present continuous, Type 1).
    \item \eng{We \textbf{will not give in} to intimidation.} $\rightarrow$ Nous \textbf{ne céderons pas} à l'intimidation. (Future + négation, Type 1).
    \item \eng{The committee \textbf{has looked into} the allegations.} $\rightarrow$ Le comité \textbf{a enquêté sur} les allégations. (Present perfect, Type 2 inséparable).
    \item \eng{Citizens must \textbf{put up with} frequent power cuts.} $\rightarrow$ Les citoyens doivent \textbf{supporter} les fréquentes coupures d'électricité. (Modal + Type 3).
    \item \eng{The allegations \textbf{were looked into} by the police.} $\rightarrow$ Les allégations \textbf{ont fait l'objet d'une enquête} de la part de la police. (Passif, Type 2).
\end{itemize}
\end{exemplebox}

\courssub{Transparence vs Idiomaticité : deux degrés de sens}

Tous les phrasal verbs ne sont pas des énigmes. Il existe un continuum entre le sens littéral (compositionnel) et le sens figuré (idiomatique).

\begin{definition}[Phrasal verbs transparents vs idiomatiques]
\begin{itemize}
    \item \textbf{Transparents (littéraux) :} Le sens se déduit de la somme des parties. La particule garde son sens spatial/directionnel.
    \item \textbf{Semi-transparents :} Sens étendu, métaphore accessible (ex. \eng{wake up} = se réveiller, \eng{turn on} = allumer).
    \item \textbf{Idiomatiques (figures) :} Le sens global est imprévisible. Il faut l'apprendre comme un mot de vocabulaire unique.
\end{itemize}
\end{definition}

\begin{center}
\begin{tabularx}{\linewidth}{|X|X|X|}
\hline
\rowcolor{popblueL}
\textbf{Verbe} & \textbf{Degré} & \textbf{Exemple \& Traduction} \\
\hline
\eng{sit down} / \eng{stand up} / \eng{go out} / \eng{come in} & \textbf{Transparent} & \eng{Please \textbf{sit down}.} = Asseyez-vous. \\
\eng{wake up} / \eng{get up} / \eng{turn on/off} & \textbf{Semi-transparent} & \eng{Turn off the generator.} = Éteignez le groupe. \\
\hline
\eng{give up} / \eng{give in} & \textbf{Idiomatique} & \eng{Don't \textbf{give up} the fight.} = N'abandonnez pas le combat. \\
\eng{put off} / \eng{put up with} & \textbf{Idiomatique} & \eng{The election was \textbf{put off}.} = L'élection a été \textbf{reportée}. \\
\eng{stand for} (tolérer/représenter) & \textbf{Idiomatique} & \eng{What does NGO \textbf{stand for}?} = Que \textbf{signifie} ONG ? \\
\eng{carry out} / \eng{look into} & \textbf{Idiomatique} & \eng{Police \textbf{carried out} an investigation.} = La police a \textbf{mené} une enquête. \\
\hline
\end{tabularx}
\end{center}

\begin{aretenir}[Phrasal verbs clés du thème \textit{Citizenship and Human Rights}]
\begin{center}
\begin{tabularx}{\linewidth}{|l|l|X|}
\hline
\rowcolor{popblueL}
\textbf{Phrasal Verb} & \textbf{Type} & \textbf{Sens français (contexte)} \\
\hline
\eng{stand up for} & 3 & défendre (droits, cause) \\
\eng{speak out (against)} & 1 & s'exprimer, dénoncer publiquement \\
\eng{look into} & 2 & enquêter sur \\
\eng{carry out} & 1 & mener, réaliser (enquête, réforme) \\
\eng{call on / upon} & 2 & appeler à, exhorter \\
\eng{step down} & 1 & démissionner \\
\eng{give in (to)} & 1 & céder (pression, exigences) \\
\eng{put up with} & 3 & supporter, tolérer (injustice) \\
\eng{stand for} & 2 & représenter / signifier (sigle) / tolérer (négatif) \\
\eng{live up to} & 3 & être à la hauteur de (attentes) \\
\eng{let down} & 1 & décevoir \\
\eng{come into force} & 2 & entrer en vigueur (loi) \\
\hline
\end{tabularx}
\end{center}
\end{aretenir}

\coursec{Comparaison français/anglais et pièges des francophones}

\courssub{Piège n°1 : Le \textbf{-s} de la 3e personne sur la particule}

C'est l'erreur la plus fréquente chez les francophones qui traitent la particule comme un deuxième verbe ou un adjectif.

\begin{attention}[Le \textbf{-s} ne va JAMAIS sur la particule]
\begin{itemize}
    \item \textbf{Faux :} \eng{*He stand \textbf{ups} for justice.} / \eng{*She look \textbf{afters} the kids.} / \eng{*It \textbf{puts offs} the meeting.}
    \item \textbf{Vrai :} \eng{He \textbf{stands} up for justice.} / \eng{She \textbf{looks} after the kids.} / \eng{It \textbf{puts} off the meeting.}
\end{itemize}
\textbf{Règle mnémotechnique :} \textit{« Le -s colle au verbe, la particule reste nue. »}
\end{attention}

\courssub{Piège n°2 : Confusion \eng{give up} vs \eng{give in}}

\begin{attention}[Nuance sémantique cruciale]
\begin{itemize}
    \item \eng{Give up} = \textbf{Abandonner} une tentative, un objectif, une habitude (souvent réflexif : \eng{give up doing sth}).
    \eng{He \textbf{gave up} smoking.} (Il a arrêté de fumer). \eng{Don't \textbf{give up} the fight.} (N'abandonnez pas le combat).
    \item \eng{Give in} = \textbf{Céder} sous la pression, se rendre (souvent à un adversaire, une demande, une force). Structure : \eng{give in to sb/sth}.
    \eng{The government refused to \textbf{give in} to the strikers' demands.} (Le gouvernement a refusé de céder aux revendications des grévistes).
\end{itemize}
\end{attention}

\courssub{Piège n°3 : \eng{Stand for} — deux sens opposés selon le contexte}

\begin{attention}[Polysémie de \eng{stand for}]
\begin{itemize}
    \item \textbf{Sens 1 (Positif/Neutre) :} Représenter, signifier (sigles, symboles), être candidat.
    \eng{MP \textbf{stands for} Member of Parliament.} (MP \textbf{signifie} Membre du Parlement).
    \eng{He \textbf{stands for} election in the North-West.} (Il \textbf{se présente} aux élections dans le Nord-Ouest).
    \item \textbf{Sens 2 (Négatif, souvent à la forme négative) :} Tolérer, accepter (généralement \eng{not stand for} / \eng{won't stand for}).
    \eng{We \textbf{won't stand for} this injustice.} (Nous ne \textbf{tolérerons pas} cette injustice).
\end{itemize}
\end{attention}

\courssub{Piège n°4 : Traduction littérale des prépositions (\eng{to}, \eng{with}, \eng{against})}

\begin{attention}[Les particules ne sont pas des prépositions françaises]
\begin{itemize}
    \item \eng{Look \textbf{into}} $\neq$ *Regarder \textbf{dans}. $\rightarrow$ \textbf{Enquêter sur} (pas de \eng{on/sur} après).
    \item \eng{Stand \textbf{up for}} $\neq$ *Se lever \textbf{pour}. $\rightarrow$ \textbf{Défendre} (COD direct en anglais, pas de \eng{for} en français).
    \item \eng{Put \textbf{up with}} $\neq$ *Mettre \textbf{en haut avec}. $\rightarrow$ \textbf{Supporter} (COD direct).
    \item \eng{Call \textbf{on}} $\neq$ *Appeler \textbf{sur}. $\rightarrow$ \textbf{Exhorter, demander à qqn de faire qqch} (\eng{call on sb to do sth}).
\end{itemize}
\end{attention}

\coursec{Entraînement guidé et production}

\courssub{Exercice résolu : Identification, conjugaison et syntaxe}

\begin{exoresolu}[Mettre le verbe entre parenthèses à la forme correcte (temps indiqué) et placer l'objet s'il y a lieu]
\textbf{Énoncé :}
\begin{enumerate}
    \item The activists \_\_\_\_\_\_\_\_\_\_ (\eng{stand up for}, Present Simple, 3e pl.) human rights daily.
    \item The meeting \_\_\_\_\_\_\_\_\_\_ (\eng{put off}, Past Simple Passive) until next week.
    \item The leader refuses to \_\_\_\_\_\_\_\_\_\_ (\eng{give in}, Infinitif) to pressure.
    \item Journalists \_\_\_\_\_\_\_\_\_\_ (\eng{speak out}, Present Continuous, 3e pl.) against censorship right now.
    \item We must \_\_\_\_\_\_\_\_\_\_ (\eng{look into}, Infinitif modal) these allegations immediately.
    \item Turn \_\_\_\_\_\_\_\_\_\_ (\eng{turn off}, Impératif) the light, please. \textit{(Test pronominal : Turn \_\_\_\_\_\_\_\_\_\_ it off.)}
    \item The allegations \_\_\_\_\_\_\_\_\_\_ (\eng{look into}, Present Perfect Passive) by the commission.
\end{enumerate}

\tcblower

\textbf{Solution détaillée :}
\begin{enumerate}
    \item \textbf{stand up for} $\rightarrow$ Present Simple, 3e pluriel : \textbf{stand up for}. (Pas de -s au pluriel). \\
    \eng{The activists \textbf{stand up for} human rights daily.}
    \item \textbf{put off} $\rightarrow$ Past Simple Passive : \textbf{was put off}. (Verbe \eng{be} au passé + participe passé \eng{put} — irrégulier : put/put/put). La particule \eng{off} ne bouge pas. \\
    \eng{The meeting \textbf{was put off} until next week.}
    \item \textbf{give in} $\rightarrow$ Infinitif après \eng{refuses to} : \textbf{give in}. (Jamais \eng{gives in} ou \eng{gave in} après \eng{to}). \\
    \eng{The leader refuses to \textbf{give in} to pressure.}
    \item \textbf{speak out} $\rightarrow$ Present Continuous, 3e pluriel : \textbf{are speaking out}. (Auxiliaire \eng{be} conjugué + base verbale \eng{speaking} + particule). \\
    \eng{Journalists \textbf{are speaking out} against censorship right now.}
    \item \textbf{look into} $\rightarrow$ Infinitif après modal \eng{must} : \textbf{look into}. (Modal + base verbale sans \eng{to}). Type 2 inséparable : l'objet \eng{these allegations} suit obligatoirement la préposition \eng{into}. \\
    \eng{We must \textbf{look into} these allegations immediately.}
    \item \textbf{Turn off} (Impératif). Objet pronominal \textbf{obligatoirement} entre V et Adv : \eng{Turn \textbf{it} off}. \\
    \eng{\textbf{Turn off} the light, please.} / \eng{\textbf{Turn it off}, please.}
    \item \textbf{look into} $\rightarrow$ Present Perfect Passive : \textbf{have been looked into}. (Auxiliaire \eng{have been} + participe passé \eng{looked} + particule \eng{into}). \\
    \eng{The allegations \textbf{have been looked into} by the commission.}
\end{enumerate}
\end{exoresolu}

\courssub{Exercice d'application : Conjugaison, syntaxe et traduction}

\begin{exercice}[Complétez avec la forme correcte. Traduisez ensuite.]
\begin{enumerate}
    \item The students \_\_\_\_\_\_\_\_\_\_ (\eng{carry on} / Present Continuous) their protest despite the rain.
    \item The new law \_\_\_\_\_\_\_\_\_\_ (\eng{come into} / Past Simple) force last January.
    \item You should not \_\_\_\_\_\_\_\_\_\_ (\eng{put up with} / Infinitif modal) injustice.
    \item The representative \_\_\_\_\_\_\_\_\_\_ (\eng{stand for} / Present Simple) the North-West region.
    \item The authorities promised to \_\_\_\_\_\_\_\_\_\_ (\eng{look into} / Infinitif) the matter.
    \item Please \_\_\_\_\_\_\_\_\_\_ (\eng{fill in} / Impératif) this form. \textit{(Pronominal : \_\_\_\_\_\_\_\_\_\_ it in.)}
    \item The suspect \_\_\_\_\_\_\_\_\_\_ (\eng{give himself up} / Past Simple) to the police yesterday.
\end{enumerate}
\end{exercice}

\begin{corrige}[Exercice : Conjugaison, syntaxe et traduction]
\begin{enumerate}
    \item \textbf{are carrying on} — Les étudiants \textbf{continuent} leur manifestation malgré la pluie. (Type 1, continu. Objet \eng{their protest} après la particule car court).
    \item \textbf{came into} — La nouvelle loi \textbf{est entrée en} vigueur en janvier dernier. (Type 2, \eng{come into force} = entrer en vigueur, passé simple irrégulier \eng{come/came/come}).
    \item \textbf{put up with} — Tu ne devrais pas \textbf{supporter / tolérer} l'injustice. (Type 3 après modal \eng{should not}).
    \item \textbf{stands for} — Le représentant \textbf{représente} la région du Nord-Ouest. (Type 2, 3e pers. sing. $\rightarrow$ -s sur \eng{stand}).
    \item \textbf{look into} — Les autorités ont promis d'\textbf{enquêter sur} l'affaire. (Type 2 après \eng{promised to}, inséparable).
    \item \textbf{Fill in} (Impératif). Objet pronominal : \eng{Fill \textbf{it} in}. — \textbf{Remplissez} ce formulaire. (Type 1 séparable).
    \item \textbf{gave himself up} — Le suspect \textbf{s'est constitué prisonnier / s'est rendu} à la police hier. (Verbe pronominal \eng{give oneself up}, Type 1, passé simple irrégulier \eng{give/gave/given}).
\end{enumerate}
\end{corrige}

\courssub{Production écrite guidée : Lettre d'opinion / Article de blog}

\begin{methode}[Rédiger un paragraphe militant (100-120 mots)]
\textbf{Consigne :} Vous êtes délégué(e) des élèves de votre lycée. Écrivez une lettre ouverte au Proviseur pour demander l'amélioration des conditions d'études (électricité, manuels, cantine). Utilisez \textbf{au moins 5 phrasal verbs} de la liste .
\textbf{Étapes :}
\begin{enumerate}
    \item \textbf{Planifiez :} Identifiez 3 problèmes + 3 solutions. Choisissez vos 5 phrasal verbs.
    \item \textbf{Rédigez :} Structure formelle (\eng{Dear Principal,} ... \eng{Yours faithfully,}). Un paragraphe par problème.
    \item \textbf{Vérifiez :} Conjugaison (sujet/verbe), syntaxe (séparabilité, pronoms), orthographe des irréguliers.
\end{enumerate}
\textbf{Amorce :} \eng{“Dear Principal, We, the students of Government High School \ldots, \textbf{call on} you to \textbf{look into} the recurring power cuts. We \textbf{will not put up with} studying in the dark anymore. We \textbf{stand for} our right to quality education. If nothing is done, we may \textbf{carry on} with peaceful action. We hope you \textbf{will not let us down}. Yours faithfully, The Students' Council.”}
\end{methode}

\begin{exercice}[Production libre]
Rédigez votre propre lettre (120 mots ± 10\%) en suivant la méthode. Échangez avec un camarade pour correction croisée (check-list : 5 phrasal verbs minimum, 0 faute de -s sur particule, syntaxe objets correcte).
\end{exercice}

\begin{savaistu}[Culture : Le \textit{Pidgin English} camerounais et les phrasal verbs]
Le \textit{Cameroon Pidgin English} (Kamtok) utilise massivement des structures de type verbe + particule, souvent calquées sur l'anglais standard mais avec une créativité lexicale propre.
\begin{itemize}
    \item \eng{Comot} (come out) = sortir, partir, publier. \eng{Di result don comot.} (Les résultats sont sortis).
    \item \eng{Du} (do) + particules : \eng{du away} (abolir, jeter), \eng{du for} (cuisiner pour), \eng{du well} (guérir).
    \item \eng{Stap} (stop/stand up) : \eng{stap kwaet} (se taire), \eng{stap fo} (représenter/soutenir).
\end{itemize}
Cette vitalité montre que la structure \textbf{Verbe + Particule} est au cœur de la pensée linguistique locale, facilitant l'appropriation des phrasal verbs standards par les apprenants camerounais.
\end{savaistu}

\coursec{Les trois types : séparables, inséparables, à trois mots}

Dans la section précédente, nous avons défini le \cle{phrasal verb} comme un verbe suivi d'une ou deux particules qui modifient son sens : \eng{give up} (abandonner), \eng{look after} (s'occuper de). Mais une question se pose immédiatement : où placer le \textbf{complément} ? Peut-on dire indifféremment \eng{turn off the light} et \eng{turn the light off} ? Et pourquoi \eng{turn off it} est-il impossible ? Toute la grammaire des phrasal verbs tient dans cette question de position. Il existe exactement \textbf{trois types} de phrasal verbs, et chacun obéit à une règle de placement différente. C'est ce que nous allons découvrir, observer, puis systématiser.

\courssub{Observation : un même verbe, deux comportements}

\begin{textebox}[A radio debate in Yaoundé]
Yesterday evening, during a radio debate on citizenship, a journalist asked the mayor: “Will the council \textbf{put off} the local elections?” The mayor answered: “No, we will not \textbf{put them off}. Citizens have been \textbf{looking forward to} this vote for months.” A listener then called and said: “Our community must \textbf{look after} its poorest members, and we will never \textbf{put up with} corruption. If leaders fail us, the people will \textbf{vote them out} at the next election.”
\end{textebox}

\textbf{Glossaire :} \emph{council} (nom, conseil municipal) ; \emph{to put off} (verbe, reporter) ; \emph{a vote} (nom, un scrutin) ; \emph{to put up with} (verbe, tolérer) ; \emph{to vote out} (verbe, chasser par le vote).

Observons attentivement les positions des compléments dans ce texte :

\begin{itemize}
\item \eng{put off the local elections} → le complément est \textbf{après} la particule ;
\item \eng{put them off} → le pronom est \textbf{entre} le verbe et la particule ;
\item \eng{look after its poorest members} → le complément est \textbf{toujours après} ;
\item \eng{put up with corruption} → trois mots soudés, complément \textbf{après} ;
\item \eng{vote them out} → à nouveau le pronom \textbf{au milieu}.
\end{itemize}

Pourquoi peut-on dire \eng{put them off} mais jamais \eng{look them after} ? Parce que \eng{put off} et \eng{look after} n'appartiennent pas au même \textbf{type} de phrasal verb. Distinguons-les un par un.

\courssub{Type 1 — Les phrasal verbs inséparables (verbe + préposition)}

\begin{definition}[Phrasal verb inséparable]
Un phrasal verb \cle{inséparable} est formé d'un verbe suivi d'une \textbf{préposition}. Le complément se place \textbf{toujours après la particule}, qu'il s'agisse d'un nom ou d'un pronom. On ne peut jamais insérer quoi que ce soit entre le verbe et la préposition.
\end{definition}

La particule est ici une véritable \textbf{préposition} : elle introduit son complément, exactement comme en français « s'occuper \emph{de} », « compter \emph{sur} ».

\begin{exemplebox}[Inséparables en contexte de citoyenneté]
\eng{We must look after the poor.} « Nous devons prendre soin des pauvres. »\par
\eng{We must look after them.} « Nous devons prendre soin d'eux. »\par
\eng{A good citizen stands for justice.} « Un bon citoyen défend la justice. »\par
\eng{The police are looking into the case of corruption.} « La police enquête sur l'affaire de corruption. »\par
\eng{You can count on your neighbours in difficult times.} « Tu peux compter sur tes voisins dans les moments difficiles. »\par
\eng{Many young people apply for their voter's card online.} « Beaucoup de jeunes demandent leur carte d'électeur en ligne. »
\end{exemplebox}

\begin{propriete}[Règle de position — type inséparable]
Forme : \textbf{verbe + préposition + complément} (nom ou pronom), toujours dans cet ordre.\par
\eng{look after the children} / \eng{look after them} — mais jamais \eng{look the children after}, jamais \eng{look them after}.
\end{propriete}

Quelques inséparables fréquents au programme de Terminale : \eng{look after} (s'occuper de), \eng{look for} (chercher), \eng{look into} (enquêter sur), \eng{stand for} (défendre, représenter), \eng{count on} (compter sur), \eng{apply for} (postuler à, demander), \eng{deal with} (traiter de, s'occuper de), \eng{belong to} (appartenir à), \eng{listen to} (écouter), \eng{wait for} (attendre).

\courssub{Type 2 — Les phrasal verbs séparables (verbe + adverbe)}

\begin{definition}[Phrasal verb séparable]
Un phrasal verb \cle{séparable} est formé d'un verbe suivi d'un \textbf{adverbe de particule} (\eng{off, on, up, down, out, away, back}). Avec un \textbf{complément nominal}, deux positions sont possibles : avant ou après la particule. Avec un \textbf{pronom}, la séparation est \textbf{obligatoire}.
\end{definition}

\begin{exemplebox}[Les deux positions avec un nom]
\eng{Turn off the TV.} = \eng{Turn the TV off.} « Éteins la télévision. »\par
\eng{They put off the election.} = \eng{They put the election off.} « Ils ont reporté l'élection. »\par
\eng{The mayor gave out the forms.} = \eng{The mayor gave the forms out.} « Le maire a distribué les formulaires. »\par
\eng{We must bring up our children well.} = \eng{We must bring our children up well.} « Nous devons bien élever nos enfants. »
\end{exemplebox}

\begin{propriete}[La règle d'or : le pronom coupe toujours]
Avec un complément \textbf{pronom} (\eng{it, them, him, her, us, me}), le phrasal verb séparable est \textbf{toujours coupé} : \textbf{verbe + pronom + particule}.\par
\eng{The people voted them out.} « Le peuple les a chassés par le vote. » — jamais \eng{voted out them}.\par
\eng{They put it off.} « Ils l'ont reportée. » — jamais \eng{put off it}.\par
\eng{Turn it off, please.} « Éteins-la, s'il te plaît. » — jamais \eng{turn off it}.
\end{propriete}

C'est la règle la plus testée de toute la grammaire des phrasal verbs. Mémorisez le réflexe : \textbf{pronom = milieu}.

Quelques séparables fréquents : \eng{put off} (reporter), \eng{turn off / turn on} (éteindre / allumer), \eng{give up} (abandonner), \eng{give out} (distribuer), \eng{bring up} (élever, aborder un sujet), \eng{vote out} (chasser par le vote), \eng{carry out} (mener à bien), \eng{fill in} (remplir un formulaire), \eng{pick up} (ramasser, aller chercher), \eng{throw away} (jeter).

\courssub{Type 3 — Les phrasal verbs à trois mots (verbe + adverbe + préposition)}

\begin{definition}[Phrasal verb à trois mots]
Un phrasal verb \cle{à trois mots} combine un verbe, un adverbe et une préposition : \eng{put up with}, \eng{look forward to}, \eng{come up with}. Ces verbes sont \textbf{toujours inséparables} : rien ne peut s'intercaler entre les trois éléments, et le complément vient après la préposition finale.
\end{definition}

\begin{exemplebox}[Trois mots, zéro séparation]
\eng{We will not put up with corruption.} « Nous ne tolérerons pas la corruption. »\par
\eng{We will not put up with it.} « Nous ne la tolérerons pas. »\par
\eng{Citizens look forward to the elections.} « Les citoyens attendent les élections avec impatience. »\par
\eng{The students came up with a brilliant idea.} « Les élèves ont trouvé une idée géniale. »\par
\eng{Women must stand up for their rights.} « Les femmes doivent défendre leurs droits. »\par
\eng{The government has run out of money for the project.} « Le gouvernement n'a plus d'argent pour le projet. »
\end{exemplebox}

\begin{attention}[Jamais de coupure dans un phrasal verb à trois mots]
Les phrasal verbs à trois mots ne se séparent \textbf{jamais}, même avec un pronom : on dit \eng{look forward to it} « je l'attends avec impatience », et non \eng{look it forward to}. De même : \eng{put up with them} « les tolérer », jamais \eng{put them up with}. Les trois mots forment un bloc soudé.
\end{attention}

\courssub{Tableau récapitulatif des trois types}

\begin{center}
\begin{tabularx}{\linewidth}{|X|X|X|}
\hline
\rowcolor{popblueL} Inséparable (verbe + prép.) & Séparable (verbe + adv.) & Trois mots (verbe + adv. + prép.) \\
\hline
\eng{look after the poor} / \eng{look after them} & \eng{put off the election} ou \eng{put the election off} & \eng{put up with corruption} \\
\hline
« s'occuper des pauvres / d'eux » & « reporter l'élection » ; avec pronom : \eng{put it off} & « tolérer la corruption » ; \eng{put up with it} \\
\hline
\eng{stand for justice} « défendre la justice » & \eng{vote them out} « les chasser par le vote » & \eng{look forward to the elections} « attendre les élections avec impatience » \\
\hline
Complément toujours après & Nom : 2 positions ; pronom : au milieu & Jamais de séparation \\
\hline
\end{tabularx}
\end{center}

\courssub{Comment reconnaître le type ? L'astuce du test}

\begin{methode}[Le test du pronom « it »]
Pour savoir si un phrasal verb est séparable, procédez en trois étapes :
\begin{enumerate}
\item Remplacez le complément nominal par le pronom \eng{it} (ou \eng{them}).
\item Placez ce pronom \textbf{entre} le verbe et la particule : \eng{verbe + it + particule}.
\item Si la phrase obtenue est correcte, le verbe est \textbf{séparable} ; si elle est impossible, il est \textbf{inséparable}.
\end{enumerate}
Exemples : \eng{put off the meeting} → \eng{put it off} est correct → séparable. \eng{look after the children} → \eng{look it after} est impossible → inséparable. Autre indice : si la particule peut se déplacer après le complément nominal, c'est un \textbf{adverbe} (séparable) ; sinon c'est une \textbf{préposition} (inséparable).
\end{methode}

\begin{popfigure}
\begin{tikzpicture}[
node distance=0.9cm,
q/.style={diamond, draw=popblue, fill=popblueL, align=center, inner sep=1.5pt, aspect=2.2, font=\small},
box/.style={rectangle, rounded corners, draw=popdark, fill=popcream, align=center, inner sep=5pt, font=\small},
ok/.style={rectangle, rounded corners, draw=popgreen, fill=popgreenL, align=center, inner sep=5pt, font=\small},
arr/.style={-{Stealth[length=2.5mm]}, thick, popdark}]
\node[q] (q1) {Phrasal verb ?};
\node[q, below=of q1, align=center] (q2){Complément\\= pronom ?};
\node[ok, below left=1cm and -1.2cm of q2, align=center] (yes){OUI : verbe + pronom + particule\\\eng{turn it off}};
\node[box, below right=1cm and -1.2cm of q2, align=center] (no){NON : si séparable, les deux\\positions sont possibles\\\eng{turn off the TV / turn the TV off}};
\draw[arr] (q1) -- (q2);
\draw[arr] (q2) -| (yes) node[midway, above left, font=\small\itshape, popgreen] {oui};
\draw[arr] (q2) -| (no) node[midway, above right, font=\small\itshape, poporange] {non};
\end{tikzpicture}
\legende{Organigramme de décision : où placer le complément d'un phrasal verb ?}
\end{popfigure}

\courssub{Comparaison avec le français et pièges des francophones}

En français, le verbe et sa préposition ne se séparent jamais (« s'occuper \emph{des} enfants »), et le pronom se place \emph{avant} le verbe (« nous \emph{les} éteignons »). L'anglais fait exactement l'inverse avec les séparables : le pronom va \emph{au milieu} du phrasal verb. D'où les erreurs typiques :

\begin{center}
\begin{tabularx}{\linewidth}{|X|X|X|}
\hline
\rowcolor{popblueL} Français & Erreur fréquente & Forme correcte \\
\hline
« Éteins-la. » & \eng{Turn off it.} & \eng{Turn it off.} \\
\hline
« Ils l'ont reportée. » & \eng{They put off it.} & \eng{They put it off.} \\
\hline
« Nous nous occupons d'eux. » & \eng{We look them after.} & \eng{We look after them.} \\
\hline
« Je l'attends avec impatience. » & \eng{I look it forward to.} & \eng{I look forward to it.} \\
\hline
\end{tabularx}
\end{center}

\begin{attention}[Pièges à éviter absolument]
\begin{itemize}
\item Ne placez jamais un pronom après la particule d'un verbe séparable : \eng{turn off it} est une faute grave.
\item Ne coupez jamais un inséparable : \eng{look them after} est impossible.
\item Attention au \eng{to} de \eng{look forward to} : c'est une préposition, donc suivi d'un nom ou d'un \textbf{gérondif} : \eng{I look forward to voting.} « J'ai hâte de voter. », et non \eng{to vote}.
\end{itemize}
\end{attention}

\begin{exoresolu}[Transformation avec un pronom]
Énoncé : réécrivez chaque phrase en remplaçant le complément souligné par un pronom (\eng{it} ou \eng{them}). 1) \eng{The council put off \textbf{the elections}.} 2) \eng{We must look after \textbf{the refugees}.} 3) \eng{Citizens will not put up with \textbf{injustice}.}
\tcblower
Solution détaillée : 1) \eng{put off} est séparable (test : \eng{put it off} est correct) ; le pronom va au milieu → \eng{The council put them off.} « Le conseil les a reportées. » (\eng{elections} est pluriel, donc \eng{them}). 2) \eng{look after} est inséparable ; le pronom reste après la particule → \eng{We must look after them.} « Nous devons nous occuper d'eux. » 3) \eng{put up with} est un verbe à trois mots, jamais coupé → \eng{Citizens will not put up with it.} « Les citoyens ne la toléreront pas. »
\end{exoresolu}

\begin{savaistu}[Cette règle tombe au Bac !]
À l'épreuve d'anglais du Baccalauréat, la consigne classique \eng{Rewrite the sentence replacing the underlined noun by a pronoun} teste exactement la règle de séparation que vous venez d'apprendre. Le correcteur attend \eng{They put it off} et sanctionne \eng{They put off it}. Un point facile à gagner — à condition d'avoir le réflexe « pronom = milieu » pour les séparables, et « pronom = après » pour les inséparables et les verbes à trois mots.
\end{savaistu}

\begin{exercice}[À vous de jouer]
Classez ces phrasal verbs en trois colonnes (inséparable / séparable / trois mots), puis réécrivez les phrases 1 à 3 avec un pronom : \eng{look into, bring up, come up with, fill in, count on, stand up for}. 1) \eng{Fill in the form.} 2) \eng{They look into the complaint.} 3) \eng{She came up with a solution.}
\end{exercice}

\begin{corrige}[Exercice — À vous de jouer]
Inséparables : \eng{look into, count on}. Séparables : \eng{bring up, fill in}. Trois mots : \eng{come up with, stand up for}. 1) \eng{Fill it in.} (séparable, pronom au milieu). 2) \eng{They look into it.} (inséparable, pronom après). 3) \eng{She came up with it.} (trois mots, jamais coupé).
\end{corrige}

\coursec{Phrasal verbs de la citoyenneté et des droits de l'homme}

Dans la section précédente, nous avons classé les phrasal verbs en trois types : séparables (\eng{to turn the radio off} / \eng{to turn it off}), inséparables (\eng{to look after the children}) et à trois mots (\eng{to put up with injustice}). Nous allons maintenant mettre cette grammaire au service du chapitre : la citoyenneté et les droits de l'homme. En effet, les articles de presse, les discours politiques et les rapports des ONG anglophones regorgent de phrasal verbs. Les maîtriser, c'est pouvoir lire un article de la BBC ou du \emph{Guardian} sur une élection, un procès ou une manifestation — et c'est aussi enrichir considérablement vos propres productions écrites au Bac, où l'emploi correct d'un phrasal verb idiomatique est toujours valorisé.

\courssub{Observation : un texte pour découvrir le vocabulaire}

\begin{textebox}[After the crisis — an original press article]
When the law broke down in the northern region, ordinary citizens refused to give in. Local NGOs helped out the victims, gave out food and blankets, and took in hundreds of refugees who had lost everything. Volunteers looked after the wounded in temporary camps, while journalists looked into the scandal that had caused the violence. Their reports showed that some officials had taken away basic rights from minority groups.

Civil society leaders called on people to speak out against injustice. “We must stand up for our rights and carry on the struggle peacefully,” declared Amina Bello, a human rights activist from Bamenda. “Nobody should put off this fight until tomorrow.”

The government finally set up an independent commission to investigate the abuses. The commission's findings brought about real change: a new law was passed, and at the next election, citizens turned up in huge numbers to vote the corrupt leaders out. A young teacher even decided to stand for parliament. “Voting is not enough,” she said. “We must also help out in our communities every day.”
\end{textebox}

\textbf{Glossaire}

\begin{center}
\begin{tabularx}{\linewidth}{|l|l|X|}
\hline
\rowcolor{popblueL} \textbf{Mot anglais} & \textbf{Nature} & \textbf{Traduction} \\
\hline
to break down & phrasal verb & s'effondrer, cesser de fonctionner \\
\hline
to give in & phrasal verb & céder, abandonner \\
\hline
to help out & phrasal verb & dépanner, donner un coup de main \\
\hline
to give out & phrasal verb & distribuer \\
\hline
to take in & phrasal verb & accueillir, héberger \\
\hline
to look after & phrasal verb & s'occuper de \\
\hline
to look into & phrasal verb & enquêter sur, examiner \\
\hline
to take away & phrasal verb & retirer, enlever (un droit) \\
\hline
to speak out & phrasal verb & s'exprimer publiquement, dénoncer \\
\hline
to stand up for & phrasal verb (3 mots) & défendre (une cause, un droit) \\
\hline
to carry on & phrasal verb & continuer, poursuivre \\
\hline
to put off & phrasal verb & reporter, remettre à plus tard \\
\hline
to set up & phrasal verb & créer, mettre en place (une institution) \\
\hline
to bring about & phrasal verb & provoquer, entraîner (un changement) \\
\hline
to turn up & phrasal verb & venir, se présenter (en grand nombre) \\
\hline
to vote out & phrasal verb & chasser par le vote, battre aux élections \\
\hline
to stand for & phrasal verb & se porter candidat à ; représenter, signifier \\
\hline
abuse & nom & abus \\
\hline
findings & nom (pluriel) & conclusions (d'une enquête) \\
\hline
\end{tabularx}
\end{center}

\textbf{Questions de compréhension}
\begin{enumerate}
\item What did the NGOs do for the victims? Mention three actions.
\item Why did the government set up an independent commission?
\item What change did the commission's findings bring about?
\item What did citizens do at the next election?
\end{enumerate}

\begin{corrige}[Questions de compréhension]
\begin{enumerate}
\item \eng{They helped out the victims, gave out food and blankets, and took in hundreds of refugees.} (Trois actions : dépanner, distribuer, accueillir.)
\item \eng{The government set up an independent commission to investigate the abuses.} (Pour enquêter sur les abus.)
\item \eng{The findings brought about real change: a new law was passed.} (Un vrai changement : une nouvelle loi a été votée.)
\item \eng{Citizens turned up in huge numbers and voted the corrupt leaders out.} (Ils sont venus voter en masse et ont chassé les dirigeants corrompus.)
\end{enumerate}
\end{corrige}

\courssub{Champ lexical 1 : défendre et lutter}

\begin{definition}[Les verbes du combat citoyen]
\begin{itemize}
\item \cle{stand up for} + nom = défendre (un droit, une personne, une cause). Inséparable, trois mots. \eng{We must stand up for freedom of speech.} « Nous devons défendre la liberté d'expression. »
\item \cle{speak out} (against/about) = s'exprimer publiquement, dénoncer. Inséparable. \eng{Students spoke out against corruption.} « Les étudiants ont dénoncé publiquement la corruption. »
\item \cle{fight for} + nom = lutter pour. Inséparable. \eng{They fought for equal rights.} « Ils ont lutté pour l'égalité des droits. »
\item \cle{carry on} (+ with / -ing) = continuer, poursuivre. Inséparable. \eng{She carried on campaigning despite the threats.} « Elle a continué sa campagne malgré les menaces. »
\item \cle{give in} = céder, abandonner. Inséparable. \eng{The protesters refused to give in.} « Les manifestants ont refusé de céder. »
\end{itemize}
\end{definition}

\begin{exemplebox}[En contexte camerounais et international]
\eng{In Bamenda, teachers stood up for the rights of their students.} « À Bamenda, les enseignants ont défendu les droits de leurs élèves. »

\eng{Journalists who speak out sometimes face threats.} « Les journalistes qui dénoncent publiquement font parfois face à des menaces. »

\eng{Nelson Mandela fought for justice and never gave in.} « Nelson Mandela a lutté pour la justice et n'a jamais cédé. »

\eng{After the meeting, the activists carried on with their project.} « Après la réunion, les militants ont poursuivi leur projet. »
\end{exemplebox}

\begin{attention}[Conjugaison : verbes irréguliers à connaître]
Plusieurs de ces phrasal verbs reposent sur des verbes irréguliers. Au passé et au participe passé : \eng{stand -- stood -- stood} (\eng{They stood up for their rights}), \eng{fight -- fought -- fought} (\eng{She fought for justice}), \eng{give -- gave -- given} (\eng{He never gave in}), \eng{speak -- spoke -- spoken} (\eng{They have spoken out}). Ne dites jamais \emph{standed}, \emph{fighted} ou \emph{gived} !
\end{attention}

\courssub{Champ lexical 2 : vie civique et élections}

\begin{definition}[Les verbes de la démocratie]
\begin{itemize}
\item \cle{vote out} = chasser (quelqu'un) par le vote. \textbf{Séparable} : \eng{to vote the leaders out} / \eng{to vote them out}. \eng{The people voted the mayor out.} « Le peuple a chassé le maire par le vote. »
\item \cle{stand for} = se porter candidat à (une élection) ; représenter, signifier. Inséparable. \eng{She stood for parliament in 2020.} « Elle s'est portée candidate au parlement en 2020. »
\item \cle{turn up} = venir, se présenter, arriver. Inséparable. \eng{Thousands of voters turned up at the polling stations.} « Des milliers d'électeurs se sont présentés aux bureaux de vote. »
\item \cle{bring about} = provoquer, entraîner (un changement). Séparable : \eng{to bring about change} / \eng{to bring it about}. \eng{The election brought about a peaceful change.} « L'élection a provoqué un changement pacifique. »
\end{itemize}
\end{definition}

\begin{attention}[to stand for : ne pas confondre avec to stand]
\eng{To stand} seul signifie « être debout » ou « supporter » (\eng{I can't stand noise} = « Je ne supporte pas le bruit »). \eng{To stand for} est un phrasal verb à part entière avec deux sens : « se porter candidat » et « signifier ». Ne traduisez jamais \eng{He stood for the election} par « Il s'est levé pour l'élection » !
\end{attention}

\courssub{Champ lexical 3 : justice et loi}

\begin{definition}[Les verbes du système judiciaire]
\begin{itemize}
\item \cle{break down} = s'effondrer, cesser de fonctionner (loi, dialogue, système). Inséparable. \eng{Talks between the two sides broke down.} « Les pourparlers entre les deux parties ont échoué. »
\item \cle{put off} = reporter, remettre à plus tard. Séparable : \eng{to put the trial off} / \eng{to put it off}. Suivi d'un verbe, il exige la forme en \emph{-ing} : \eng{They put off holding the trial.} \eng{The judge put off the hearing.} « Le juge a reporté l'audience. »
\item \cle{look into} = enquêter sur, examiner. Inséparable. \eng{The police are looking into the case.} « La police enquête sur l'affaire. »
\item \cle{take away} = retirer, enlever (un droit, une liberté). Séparable. \eng{No one can take away your basic rights.} « Personne ne peut vous retirer vos droits fondamentaux. »
\item \cle{set up} = créer, mettre en place (une commission, une institution). Séparable. \eng{They set up a committee to investigate.} « Ils ont créé un comité pour enquêter. »
\end{itemize}
\end{definition}

\begin{exemplebox}[Exemples fondateurs à retenir]
\eng{The commission was set up to protect human rights.} « La commission a été créée pour protéger les droits de l'homme. » (Notez l'emploi du passif, très fréquent dans les textes officiels.)

\eng{The court looked into the case.} « Le tribunal a enquêté sur l'affaire. »

\eng{When the law breaks down, citizens suffer.} « Quand la loi s'effondre, les citoyens souffrent. »

\eng{The authorities took away his passport.} « Les autorités lui ont retiré son passeport. »
\end{exemplebox}

\courssub{Champ lexical 4 : solidarité}

\begin{definition}[Les verbes de l'entraide]
\begin{itemize}
\item \cle{look after} = s'occuper de, prendre soin de. Inséparable. \eng{She looks after elderly neighbours.} « Elle s'occupe de ses voisins âgés. »
\item \cle{help out} = dépanner, donner un coup de main. Séparable : \eng{to help somebody out}. \eng{Volunteers helped out the flood victims in Douala.} « Des bénévoles ont dépanné les victimes des inondations à Douala. »
\item \cle{give out} = distribuer. Séparable. \eng{They gave out food and clothes at the market.} « Ils ont distribué de la nourriture et des vêtements au marché. »
\item \cle{take in} = accueillir, héberger. Séparable. \eng{The village took in fifty refugee families.} « Le village a accueilli cinquante familles de réfugiés. »
\end{itemize}
\end{definition}

\begin{attention}[Faux amis et pièges de traduction]
\begin{itemize}
\item \eng{To take in} ne signifie pas « prendre dans » : c'est « accueillir » (ou « tromper » : \eng{I was taken in by his lies} = « J'ai été trompé par ses mensonges »).
\item \eng{To give out} n'est pas « donner dehors » mais « distribuer » ; attention, il signifie aussi « tomber en panne, s'épuiser » (\eng{My voice gave out} = « Ma voix m'a lâché »).
\item \eng{To look after} ne veut pas dire « regarder après » : c'est « s'occuper de ». Le français « regarder » se dit \eng{to look at}.
\item \eng{To put off} ne signifie pas « éteindre » (qui se dit \eng{to turn off} ou \eng{to switch off}) : c'est « reporter ».
\end{itemize}
\end{attention}

\courssub{Tableau récapitulatif du vocabulaire}

\begin{center}
\begin{tabularx}{\linewidth}{|l|l|X|}
\hline
\rowcolor{popblueL} \textbf{Phrasal verb} & \textbf{Traduction} & \textbf{Exemple en contexte (citizenship)} \\
\hline
stand up for & défendre & \eng{Stand up for your rights!} \\
\hline
speak out & dénoncer publiquement & \eng{Citizens spoke out against injustice.} \\
\hline
fight for & lutter pour & \eng{They fought for democracy.} \\
\hline
carry on & continuer & \eng{Carry on the peaceful struggle.} \\
\hline
give in & céder & \eng{Never give in to intimidation.} \\
\hline
vote out & chasser par le vote & \eng{Voters voted the corrupt mayor out.} \\
\hline
stand for & se porter candidat ; signifier & \eng{She stood for election in Yaoundé.} \\
\hline
turn up & venir, se présenter & \eng{Many young people turned up to vote.} \\
\hline
bring about & provoquer (un changement) & \eng{The reform brought about progress.} \\
\hline
break down & s'effondrer & \eng{The dialogue broke down.} \\
\hline
put off & reporter & \eng{The trial was put off until May.} \\
\hline
look into & enquêter sur & \eng{The court looked into the scandal.} \\
\hline
take away & retirer (un droit) & \eng{They took away his freedom.} \\
\hline
set up & créer (une institution) & \eng{A commission was set up.} \\
\hline
look after & s'occuper de & \eng{She looks after orphans.} \\
\hline
help out & dépanner & \eng{Neighbours helped out the family.} \\
\hline
give out & distribuer & \eng{NGOs gave out blankets.} \\
\hline
take in & accueillir & \eng{The town took in refugees.} \\
\hline
\end{tabularx}
\end{center}

\begin{popfigure}
\begin{tikzpicture}[mindmap, grow cyclic, every node/.style={concept}, concept color=popblue!40, text=popdark, root concept/.append style={concept color=popblue!70, text=white, font=\bfseries}, level 1/.append style={level distance=3.2cm, sibling angle=90}, level 2/.append style={level distance=2.2cm, sibling angle=45, concept color=popgreenL, font=\small}]
\node[root concept]{Phrasal verbs of citizenship}
  child[concept color=poporangeL]{ node {Defend and fight}
    child{ node {stand up for} }
    child{ node {speak out} }
    child{ node {fight for} }
    child{ node {carry on} }
    child{ node {give in} }
  }
  child[concept color=popgoldL]{ node {Elections}
    child{ node {vote out} }
    child{ node {stand for} }
    child{ node {turn up} }
    child{ node {bring about} }
  }
  child[concept color=poppinkL]{ node {Justice and law}
    child{ node {look into} }
    child{ node {set up} }
    child{ node {take away} }
    child{ node {put off} }
    child{ node {break down} }
  }
  child[concept color=poppurpleL]{ node {Solidarity}
    child{ node {look after} }
    child{ node {help out} }
    child{ node {give out} }
    child{ node {take in} }
  };
\end{tikzpicture}
\legende{Carte mentale : les quatre champs lexicaux des phrasal verbs de la citoyenneté.}
\end{popfigure}

\courssub{Comparaison avec le français}

\begin{propriete}[Un verbe français = souvent un phrasal verb anglais]
Le français exprime ces idées avec un verbe unique et soutenu (« distribuer », « enquêter », « reporter »), alors que l'anglais courant et journalistique préfère le phrasal verb (\eng{give out}, \eng{look into}, \eng{put off}). En production écrite, utiliser le phrasal verb donne un style naturel et idiomatique, très apprécié au Bac. À l'inverse, traduire mot à mot un phrasal verb en français produit presque toujours une absurdité : la particule ne se traduit jamais seule.
\end{propriete}

\begin{center}
\begin{tabularx}{\linewidth}{|X|X|X|}
\hline
\rowcolor{popblueL} \textbf{Français} & \textbf{English (phrasal verb)} & \textbf{Remarque} \\
\hline
enquêter sur & to look into & inséparable \\
\hline
reporter & to put off & séparable ; ≠ « rapporter » \\
\hline
distribuer & to give out & séparable \\
\hline
s'occuper de & to look after & inséparable \\
\hline
défendre & to stand up for & trois mots, inséparable \\
\hline
céder & to give in & ≠ « donner dans » \\
\hline
créer (une institution) & to set up & séparable \\
\hline
retirer (un droit) & to take away & séparable \\
\hline
\end{tabularx}
\end{center}

\courssub{Méthode et exercices d'application}

\begin{methode}[Comment mémoriser durablement les phrasal verbs]
\begin{enumerate}
\item N'apprenez \textbf{jamais} un phrasal verb dans une liste isolée : la particule (\eng{up}, \eng{out}, \eng{in}) n'a aucun sens seule.
\item Apprenez chaque verbe \textbf{dans une phrase complète liée au thème} du chapitre : \eng{Citizens turned up to vote} est bien plus facile à retenir que « turn up = venir ».
\item Classez-les par \textbf{champ lexical} (lutte, élections, justice, solidarité), comme dans la carte mentale ci-dessus.
\item Notez toujours si le verbe est \textbf{séparable ou inséparable} : c'est ce qui détermine la place du pronom (\eng{to vote them out}, jamais \emph{to vote out them}).
\item Réutilisez-les dans \textbf{vos propres phrases} sur le Cameroun ou l'actualité : c'est la meilleure préparation à l'essai du Bac.
\end{enumerate}
\end{methode}

\begin{exoresolu}[Complétez avec le phrasal verb correct]
Complete the sentences with: \eng{set up, looked into, gave out, stand for, turned up, take away}.
\begin{enumerate}
\item The government \ldots\ a commission to investigate corruption.
\item The police \ldots\ the disappearance of the funds.
\item Volunteers \ldots\ food to the displaced families.
\item What does the letter “H” \ldots\ in “HRW”?
\item Despite the rain, voters \ldots\ in large numbers.
\item No law can \ldots\ your dignity.
\end{enumerate}
\tcblower
\textbf{Solution détaillée}
\begin{enumerate}
\item \eng{set up} — créer une institution (passé : \eng{set} est irrégulier : set -- set -- set).
\item \eng{looked into} — enquêter sur une affaire ; inséparable, passé régulier en \emph{-ed}.
\item \eng{gave out} — distribuer ; verbe irrégulier give -- gave -- given.
\item \eng{stand for} — ici au sens de « signifier » (présent, après l'auxiliaire \eng{does}, base verbale).
\item \eng{turned up} — venir en grand nombre ; passé régulier.
\item \eng{take away} — retirer un droit ; après \eng{can}, base verbale.
\end{enumerate}
\end{exoresolu}

\begin{exercice}[À vous de jouer]
\begin{enumerate}
\item Traduisez en anglais avec un phrasal verb : « Les citoyens doivent défendre leurs droits et ne jamais céder. »
\item Traduisez en anglais : « Le gouvernement a créé une commission pour enquêter sur les abus. »
\item Remplacez le verbe souligné par un phrasal verb équivalent : \eng{The judge \emph{postponed} the hearing.} / \eng{The association \emph{distributes} food every Saturday.}
\item Remettez le pronom à la bonne place : \eng{They voted out them.} / \eng{She looked after them.}
\item Rédigez trois phrases sur la vie citoyenne au Cameroun en utilisant \eng{turn up}, \eng{speak out} et \eng{help out}.
\end{enumerate}
\end{exercice}

\begin{corrige}[Exercice « À vous de jouer »]
\begin{enumerate}
\item \eng{Citizens must stand up for their rights and never give in.}
\item \eng{The government set up a commission to look into the abuses.}
\item \eng{The judge put off the hearing.} (ou \eng{put the hearing off}) ; \eng{The association gives out food every Saturday.}
\item \eng{They voted them out.} (verbe séparable : le pronom se place entre le verbe et la particule) ; \eng{She looked after them.} (verbe inséparable : le pronom suit la particule).
\item Exemples de réponses : \eng{Many young people turned up to vote in Yaoundé.} / \eng{Students spoke out against injustice on campus.} / \eng{Neighbours helped out the victims of the floods in Douala.}
\end{enumerate}
\end{corrige}

\begin{experience}[Débat de classe : “Youth and citizenship”]
Par groupes de quatre, préparez un mini-débat sur la question : \eng{Should voting be compulsory for all citizens?} Chaque élève doit utiliser au moins \textbf{quatre phrasal verbs} de la leçon dans ses prises de parole (\eng{stand up for}, \eng{speak out}, \eng{turn up}, \eng{vote out}, \eng{carry on}\ldots). Un rapporteur par groupe résume ensuite les arguments au tableau, en réutilisant les phrasal verbs au passé : \eng{Our group said that young people should turn up in large numbers\ldots}
\end{experience}

\begin{savaistu}[Les deux sens de “stand for” au Bac]
\eng{To stand for} a deux sens fréquents dans les épreuves : « représenter, signifier » — \eng{What does UNESCO stand for?} (« Que signifie UNESCO ? ») — et « se porter candidat » — \eng{to stand for election} (« se présenter aux élections »). Au Royaume-Uni, on dit \eng{to stand for parliament}, alors que les Américains préfèrent \eng{to run for office} : même idée, deux phrasal verbs différents selon le continent !
\end{savaistu}

\begin{aretenir}[L'essentiel de la leçon]
\begin{itemize}
\item Les phrasal verbs du chapitre se classent en quatre champs : \textbf{lutter} (\eng{stand up for}, \eng{speak out}, \eng{give in}), \textbf{voter} (\eng{vote out}, \eng{stand for}, \eng{turn up}), \textbf{la justice} (\eng{look into}, \eng{set up}, \eng{take away}, \eng{put off}) et \textbf{la solidarité} (\eng{look after}, \eng{help out}, \eng{give out}, \eng{take in}).
\item Verbes \textbf{séparables} : le pronom se place obligatoirement au milieu (\eng{to vote them out}). Verbes \textbf{inséparables} : le pronom suit la particule (\eng{to look after them}).
\item Attention aux verbes irréguliers : stand -- stood -- stood, fight -- fought -- fought, give -- gave -- given, take -- took -- taken, set -- set -- set, speak -- spoke -- spoken.
\item Un phrasal verb ne se traduit jamais mot à mot : apprenez-le toujours dans une phrase complète.
\end{itemize}
\end{aretenir}

\coursec{Comparaison français/anglais et pièges des francophones}

\courssub{Transition : du lexique thématique à la contrastive}

Dans la section précédente, nous avons répertorié les \cle{phrasal verbs} incontournables du thème « Citoyenneté et Droits de l'homme » (\emph{stand up for}, \emph{speak out against}, \emph{crack down on}, etc.). Nous savons désormais \emph{quels} couples verbe + particule utiliser. Reste à maîtriser \emph{comment} les manipuler sans calquer le français. C'est l'objet de cette section : observer les divergences structurelles entre nos deux langues pour neutraliser les erreurs systémiques des apprenants francophones.

\begin{textebox}{Observation : deux langues, deux logiques}
Read the following pairs of sentences. The English versions use phrasal verbs; the French versions use single verbs (often of Latin origin). Notice how the particle adds a nuance of completion, direction or aspect that French expresses by changing the verb entirely.

\eng{1. The government \textbf{put off} the election.}
« Le gouvernement a \textbf{reporté} l'élection. »

\eng{2. Citizens must \textbf{stand up for} their rights. }
« Les citoyens doivent \textbf{défendre} leurs droits. »

\eng{3. The police \textbf{cracked down on} illegal voting. }
« La police a \textbf{réprimé} le vote illégal. »

\eng{4. She \textbf{gave up} her seat for the elderly man. }
« Elle a \textbf{cédé} sa place à l'homme âgé. »

\eng{5. We need to \textbf{look into} these allegations. }
« Nous devons \textbf{enquêter sur} ces allégations. »
\end{textebox}

\begin{propriete}{Différence majeure : verbe unique (FR) vs verbe + particule (EN)}
\emph{Forme} : En français, le sens lexical est porté quasi exclusivement par le verbe (souvent d'origine latine : \emph{reporter, défendre, réprimer, céder, enquêter}). En anglais courant, le sens est porté par l'association \textbf{verbe de base (germanique) + particule/adverbe} (\emph{put off, stand up for, crack down on, give up, look into}).

\emph{Conséquence immédiate} : 
\begin{itemize}
    \item Le phrasal verb est \textbf{plus concret, plus imagé} (métaphore spatiale : \emph{up, down, off, out, in}).
    \item Il est \textbf{neutre à familier} à l'oral ; le verbe latinisé (\emph{postpone, defend, repress, surrender, investigate}) est \textbf{formel, écrit, administratif}.
    \item \textbf{Au Bac et à l'oral}, l'examinateur attend le phrasal verb naturel (\eng{put off}) plutôt que le verbe savant (\eng{postpone}), sauf registre très formel explicite.
\end{itemize}
\end{propriete}

\begin{center}
\begin{tabularx}{\linewidth}{|X|X|X|}
\hline
\rowcolor{popblueL} \textbf{Français (verbe unique)} & \textbf{English (Phrasal Verb)} & \textbf{Registre / Nuance} \\
\hline
reporter & \eng{put off} & Courant, oral, neutre \\
ajourner & \eng{put back} & Légèrement plus formel \\
défendre (une cause) & \eng{stand up for} & Engagé, concret (image : se lever) \\
soutenir & \eng{back up} & Concret (appui dans le dos) \\
réprimer / sévir & \eng{crack down on} & Très imagé (craquer le fouet) \\
céder / abandonner & \eng{give up} / \eng{give in} & \eng{give up} = arrêter ; \eng{give in} = céder sous pression \\
enquêter sur & \eng{look into} / \eng{probe into} & \eng{look into} = courant ; \eng{probe} = plus technique \\
tolérer / supporter & \eng{put up with} & Image : « mettre en haut avec » = supporter le poids \\
\hline
\end{tabularx}
\end{center}

\courssub{Les cinq pièges majeurs des francophones}

\begin{attention}[Piège n°1 : Oublier la particule — le sens s'effondre]
La particule n'est pas un « petit mot décoratif » : elle \textbf{change radicalement le sens} du verbe. L'omettre revient à dire un autre verbe.
\begin{itemize}
    \item \eng{He \textbf{gave up} smoking.} → « Il a \textbf{arrêté} de fumer. » (succès, achèvement)
    \item \eng{He \textbf{gave} smoking.} → *Incorrect / Sens impossible* (on ne « donne » pas la cigarette).
    \item \eng{The meeting \textbf{broke up} at 6 p.m.} → « La réunion s'est \textbf{terminée/dispersée} à 18 h. »
    \item \eng{The meeting \textbf{broke} at 6 p.m.} → « La réunion s'est \textbf{cassée/rompue} » (absurde).
\end{itemize}
\textbf{Règle d'or} : Apprenez le phrasal verb comme un \textbf{bloc lexical indissociable} (une seule entrée dans votre carnet : \eng{give up = abandonner}, pas \eng{give} + \eng{up}).
\end{attention}

\begin{attention}[Piège n°2 : Mauvaise place du pronom objet — le calque « éteins-le »]
En français, le pronom objet précède le verbe : « \emph{Éteins-\textbf{le}} ». En anglais, si l'objet est un \textbf{pronom} (\eng{it, them, me, him, her, us}), il \textbf{doit s'intercaler} entre le verbe et la particule (séparation obligatoire pour les phrasal verbs \textbf{séparables}).

\textbf{Règle de placement du pronom}
\begin{center}
\begin{tabular}{|l|l|l|}
\hline
\rowcolor{poporangeL} \textbf{Phrase correcte} & \textbf{Traduction} & \textbf{Structure} \\
\hline
\eng{Turn \textbf{it} off, please.} & Éteins-\textbf{le}, s'il te plaît. & V + \textbf{pronom} + Particule \\
\eng{The government \textbf{put it off}.} & Le gouvernement l'a \textbf{reporté}. & V + \textbf{pronom} + Particule \\
\eng{She \textbf{looked them up}.} & Elle les a \textbf{cherchés (dans l'annuaire)}. & V + \textbf{pronom} + Particule \\
\hline
\rowcolor{poporangeL} \textbf{Erreur fréquente} & \textbf{Pourquoi c'est faux} & \textbf{Correction} \\
\hline
\eng{*Turn off it.} & Calque français « éteins-le » & \eng{Turn it off.} \\
\eng{*Put off it.} & Calque « reporte-le » & \eng{Put it off.} \\
\eng{*Look up them.} & Calque « cherche-les » & \eng{Look them up.} \\
\hline
\end{tabular}
\end{center}

\textbf{Exception} : Les phrasal verbs \textbf{inséparables} (ex. \eng{look after, look into, stand for}) ne séparent \textbf{jamais}, même avec un pronom : \eng{look after \textbf{her}} (s'occuper d'\textbf{elle}), \eng{look into \textbf{it}} (enquêter sur \textbf{ça}).
\end{attention}

\begin{attention}[Piège n°3 : Traduction littérale de la particule — \eng{up} $\neq$ « en haut »]
Les particules (\eng{up, down, off, on, out, in, away, back}) ont des \textbf{sens aspectuels} (achèvement, début, intensité, retrait) et non seulement spatiaux. Ne les traduisez jamais mot à mot.
\begin{center}
\begin{tabularx}{\linewidth}{|X|X|X|}
\hline
\rowcolor{poppurpleL} \textbf{Phrasal Verb} & \textbf{Sens réel} & \textbf{Piège littéral à éviter} \\
\hline
\eng{eat up} & manger entièrement / finir son assiette & * « manger en haut » \\
\eng{speak up} & parler plus fort / donner son avis & * « parler en haut » \\
\eng{clean up} & nettoyer à fond / tout ranger & * « nettoyer en haut » \\
\eng{give up} & abandonner / cesser & * « donner en haut » \\
\eng{put off} & reporter / repousser & * « mettre hors de » \\
\eng{call off} & annuler & * « appeler hors de » \\
\eng{find out} & découvrir (par recherche) & * « trouver dehors » \\
\hline
\end{tabularx}
\end{center}
\textbf{Astuce} : Associez la particule à sa \textbf{valeur aspectuelle} principale :
\begin{itemize}
    \item \eng{up} = achèvement, totalité (\eng{eat up, clean up, use up}).
    \item \eng{off} = séparation, départ, annulation (\eng{put off, call off, set off}).
    \item \eng{out} = extériorisation, découverte, élimination (\eng{find out, speak out, wipe out}).
    \item \eng{down} = réduction, écriture, soumission (\eng{write down, calm down, crack down}).
\end{itemize}
\end{attention}

\begin{attention}[Piège n°4 : Préférer le verbe latinisé (registre trop formel) à l'oral]
Les anglophones natifs utilisent massivement les phrasal verbs dans la conversation quotidienne, les médias, les discours politiques télévisés. L'apprenant francophone, par sécurité, choisit le verbe latinisé (cognat) qui lui semble « plus correct ». C'est une \textbf{erreur de registre} : cela sonne pédant, administratif ou médical.

\textbf{Registre oral / courant vs Registre formel / écrit}
\begin{center}
\begin{tabular}{|l|l|l|}
\hline
\rowcolor{popgreenL} \textbf{Contexte oral / Bac Speaking / Essay} & \textbf{Contexte très formel (rapport ONU, loi)} & \textbf{Nuance} \\
\hline
\eng{The president \textbf{called off} the visit.} & The president \textbf{cancelled} the visit. & \eng{Call off} = naturel, standard. \\
\eng{They \textbf{put off} the deadline.} & They \textbf{postponed} the deadline. & \eng{Postpone} = administratif. \\
\eng{We need to \textbf{find out} the truth.} & We need to \textbf{ascertain} the truth. & \eng{Ascertain} = juridique/technique. \\
\eng{The law \textbf{cracks down on} corruption.} & The law \textbf{represses} corruption. & \eng{Repress} = fort, autoritaire. \\
\eng{He \textbf{turned down} the offer.} & He \textbf{declined} the offer. & \eng{Decline} = poli, écrit. \\
\hline
\end{tabular}
\end{center}
\textbf{Conseil} : Dans votre \emph{Essay} (épreuve écrite) et à l'oral, \textbf{privilégiez le phrasal verb} si vous le maîtrisez. C'est un marqueur de niveau \textbf{B2/C1} (compétence idiomatique).
\end{attention}

\begin{attention}[Piège n°5 : Inventer des phrasal verbs calqués sur le français — Faux amis structurels]
Le français utilise souvent \emph{à} + infinitif ou \emph{de} + infinitif là où l'anglais change de structure. Ne créez pas de phrasal verbs qui n'existent pas.
\begin{center}
\begin{tabularx}{\linewidth}{|X|X|X|}
\hline
\rowcolor{poppinkL} \textbf{Français (calque fautif)} & \textbf{Erreur inventée} & \textbf{Anglais correct} \\
\hline
assister à une réunion & \eng{*assist to a meeting} / \eng{*attend to} & \eng{attend a meeting} / \eng{turn up at a meeting} / \eng{show up at a meeting} \\
participer à & \eng{*participate to} & \eng{take part in} / \eng{join in} \\
s'occuper de (soins) & \eng{*look for} (chercher) & \eng{look after} / \eng{take care of} \\
chercher & \eng{*look after} (s'occuper) & \eng{look for} \\
enquêter sur & \eng{*look for} / \eng{*search} & \eng{look into} / \eng{investigate} \\
\hline
\end{tabularx}
\end{center}
\textbf{Rappel vital} : \eng{look for} = chercher (actif) ; \eng{look after} = s'occuper de (responsabilité) ; \eng{look into} = enquêter sur (processus). \textbf{La particule change tout le sens.}
\end{attention}

\courssub{Stratégie de choix : la méthode en 3 étapes}

\begin{methode}[Comment choisir le bon phrasal verb sans calquer le français]
\textbf{Étape 1 — Identifier le sens précis en français (contexte + registre).}
\emph{Exemple} : « Le ministre a annulé sa visite. » $\rightarrow$ sens = \emph{annuler}, registre = information courante.

\textbf{Étape 2 — Récupérer le couple [Verbe + Particule] appris en classe (mémoire lexicale), JAMAIS l'inverse.}
Ne partez pas de \eng{call} ou \eng{put} pour deviner la particule. Partez du sens « annuler » $\rightarrow$ \eng{call off} (vu en classe). Si vous hésitez entre \eng{call off} et \eng{put off}, rappelez-vous : \eng{put off} = reporter (date future) ; \eng{call off} = annuler (plus d'événement).

\textbf{Étape 3 — Vérifier la syntaxe (séparable ? inséparable ? pronom ?).}
\eng{The minister called \textbf{it} off.} (séparable, pronom intercalé).
\eng{The minister called off \textbf{the visit}.} (nom après la particule possible).

\textbf{Étape 4 — Contrôler le registre.}
Oral / Essay standard $\rightarrow$ \eng{call off} ✓. Traité international $\rightarrow$ \eng{cancel} / \eng{terminate} possible.
\end{methode}

\begin{savaistu}[Le phrasal verb, marqueur de niveau au Baccalauréat]
Les correcteurs du Baccalauréat (épreuve d'Expression Écrite — \emph{Essay} ou \emph{Letter}) et de l'Oral (\emph{Conversation}) sont formés pour repérer la \textbf{compétence idiomatique}. Un candidat qui écrit naturellement :
\eng{“ The government \textbf{called off} the demonstration instead of \textbf{cracking down on} protesters. ”}
montre une maîtrise de la \textbf{phrase verbale anglaise} (niveau B2+). Un candidat qui écrit :
\eng{“ The government \textbf{cancelled} the demonstration instead of \textbf{repressing} protesters. ”}
a un anglais correct mais \textbf{scolaire, latinisé, niveau B1}. Le phrasal verb bien employé vaut des points de \textbf{Vocabulaire} et de \textbf{Grammaire} (structure complexe). Entraînez-vous à en glisser 2 ou 3 par production écrite.
\end{savaistu}

\courssub{Tableau récapitulatif des erreurs corrigées (Illustration TikZ)}

\begin{popfigure}
\begin{tikzpicture}[
    node distance=0.2cm and 0.5cm,
    err/.style={rectangle, draw=poppink, thick, fill=poppink!10, rounded corners=3pt, minimum height=1.1cm, text width=7.5cm, align=left, font=\small},
    cor/.style={rectangle, draw=popgreen, thick, fill=popgreen!10, rounded corners=3pt, minimum height=1.1cm, text width=7.5cm, align=left, font=\small},
    hdr/.style={rectangle, draw=popdark, thick, fill=popblue, text=white, rounded corners=3pt, minimum height=0.9cm, text width=7.5cm, align=center, font=\bfseries\small},
    arr/.style={-Stealth, thick, poporange, shorten <=2pt, shorten >=2pt}
]
% En-têtes
\node[hdr] (h1) {\textbf{Erreur fréquente (Calque francophone)}};
\node[hdr, right=of h1] (h2) {\textbf{Correction + Traduction}};

% Ligne 1
\node[err, below=of h1, align=center] (e1){\eng{*He gave up it.} \\ (Il a abandonné \textbf{là-haut}.)};
\node[cor, below=of h2, align=center] (c1){\eng{He gave \textbf{it} up.} \\ « Il l'a abandonné. » (Pronom intercalé)};
\draw[arr] (e1.east) -- (c1.west) node[midway, above, font=\tiny, poporange] {Pronom $\rightarrow$ milieu};

% Ligne 2
\node[err, below=of e1, align=center] (e2){\eng{*She looks after her keys.} \\ (Elle s'occupe de ses clés.)};
\node[cor, below=of c1, align=center] (c2){\eng{She \textbf{looks for} her keys.} \\ « Elle \textbf{cherche} ses clés. » (Particule \eng{for})};
\draw[arr] (e2.east) -- (c2.west) node[midway, above, font=\tiny, poporange] {Particule : \eng{after} $\rightarrow$ \eng{for}};

% Ligne 3
\node[err, below=of e2, align=center] (e3){\eng{*The meeting broke.} \\ (La réunion s'est cassée.)};
\node[cor, below=of c2, align=center] (c3){\eng{The meeting \textbf{broke up}.} \\ « La réunion s'est \textbf{terminée/dispersée}. »};
\draw[arr] (e3.east) -- (c3.west) node[midway, above, font=\tiny, poporange] {Particule \eng{up} = achèvement};

% Ligne 4
\node[err, below=of e3, align=center] (e4){\eng{*We must assist to the debate.} \\ (Calque « assister à »).};
\node[cor, below=of c3, align=center] (c4){\eng{We must \textbf{attend} the debate.} \\ ou \eng{\textbf{turn up at} the debate.} \\ « Nous devons assister au débat. »};
\draw[arr] (e4.east) -- (c4.west) node[midway, above, font=\tiny, poporange] {Verbe simple / Phrasal \eng{turn up}};

% Ligne 5
\node[err, below=of e4, align=center] (e5){\eng{*They put off the lights.} \\ (Ils ont reporté les lumières.)};
\node[cor, below=of c4, align=center] (c5){\eng{They \textbf{turned off} the lights.} \\ « Ils ont \textbf{éteint} les lumières. »};
\draw[arr] (e5.east) -- (c5.west) node[midway, above, font=\tiny, poporange] {Verbe : \eng{put} $\rightarrow$ \eng{turn}};

% Cadre global
\node[fit=(h1)(h2)(e5)(c5), draw=popline, thick, rounded corners=5pt, inner sep=4pt] {};
\end{tikzpicture}
\legende{Tableau des 5 erreurs types : pronom mal placé, confusion de particules, particule omise, calque structurel « assister à », mauvais verbe de base.}
\end{popfigure}

\courssub{Entraînement guidé : correction d'erreurs authentiques}

\begin{exoresolu}[Correction de 5 phrases fautives — Niveau Bac]
\textbf{Consigne} : Corrigez chaque phrase. Justifiez la correction (règle de syntaxe, choix lexical, registre).

\begin{enumerate}
    \item \eng{*The NGO looked after the causes of the conflict.}
    \item \eng{*The president put off it until next month.}
    \item \eng{*We need to assist to the conference on human rights.}
    \item \eng{*The police cracked the protest down.}
    \item \eng{*She gave up to smoke last year.}
\end{enumerate}

\tcblower
\textbf{Solution détaillée}

\begin{enumerate}
    \item \textbf{Correction} : \eng{The NGO \textbf{looked into} the causes of the conflict.}
    \textbf{Justification} : \eng{Look after} = s'occuper de (soins). \eng{Look into} = enquêter sur. Contexte « causes » $\rightarrow$ enquête. \textbf{Piège n°3/5}.
    
    \item \textbf{Correction} : \eng{The president \textbf{put it off} until next month.}
    \textbf{Justification} : \eng{Put off} est séparable. Objet pronom \eng{it} $\rightarrow$ placement \textbf{obligatoire} entre verbe et particule. \textbf{Piège n°2}.
    
    \item \textbf{Correction} : \eng{We need to \textbf{attend} the conference on human rights.} (ou \eng{turn up at} / \eng{show up at}).
    \textbf{Justification} : \eng{Assist to} n'existe pas (faux ami \emph{assister à}). \eng{Attend} est transitif direct. \textbf{Piège n°5}.
    
    \item \textbf{Correction} : \eng{The police \textbf{cracked down on} the protest.}
    \textbf{Justification} : \eng{Crack down on} est \textbf{inséparable} (trois mots). La particule \eng{down} et la préposition \eng{on} ne se séparent jamais. Objet nominal \eng{the protest} $\rightarrow$ après \eng{on}. \textbf{Piège n°1/2}.
    
    \item \textbf{Correction} : \eng{She \textbf{gave up smoking} last year.} (ou \eng{quit smoking}).
    \textbf{Justification} : \eng{Give up} est suivi du \textbf{gérondif} (\eng{-ing}), pas de l'infinitif \eng{to}. \eng{Give up to} n'existe pas. \textbf{Grammaire du verbe + -ing}.
\end{enumerate}
\end{exoresolu}

\begin{exercice}[Entraînement personnel — À rendre au professeur]
\textbf{Traduisez en anglais en utilisant un phrasal verb du thème « Citoyenneté / Droits de l'homme » (liste Section 4). Attention aux pronoms, à l'ordre des mots et au gérondif.}

\begin{enumerate}
    \item Le gouvernement a reporté le référendum. (put off)
    \item Les militants défendent la liberté d'expression. (stand up for)
    \item La police a réprimé la manifestation. (crack down on)
    \item Nous devons enquêter sur ces violations. (look into)
    \item Elle a abandonné l'idée de se présenter. (give up)
    \item Éteins-la (la télévision) ! (turn off + pronom)
    \item Ils ont annulé la grève. (call off)
    \item Je ne supporterai plus cette injustice. (put up with + pronom)
    \item Le député s'est exprimé contre la loi. (speak out against)
    \item Les observateurs ont remarqué des irrégularités. (note down / observe $\rightarrow$ \eng{spot} / \eng{pick up on} ?) \textit{Indice : verbe de perception + particule.}
\end{enumerate}
\end{exercice}

\begin{corrige}[Exercice n°1]
\begin{enumerate}
    \item The government \textbf{put off} the referendum. (ou \textbf{put it off}).
    \item Activists \textbf{stand up for} freedom of speech.
    \item The police \textbf{cracked down on} the demonstration.
    \item We must \textbf{look into} these violations.
    \item She \textbf{gave up} running for office. (gérondif \eng{running}).
    \item \textbf{Turn it off}! (Pronom \eng{it} intercalé).
    \item They \textbf{called off} the strike. (ou \textbf{called it off}).
    \item I won't \textbf{put up with} this injustice \textbf{any longer}. (ou \eng{put up with it}).
    \item The MP \textbf{spoke out against} the bill.
    \item The observers \textbf{picked up on} irregularities. (\eng{Pick up on} = remarquer, détecter). \textit{Autre possibilité} : \eng{noted down} (noté par écrit).
\end{enumerate}
\end{corrige}

\coursec{Observation et découverte}

Pour entrer dans la leçon, nous partons d'une situation de communication authentique : une interview radiophonique sur l'engagement des jeunes au Cameroun. L'objectif est d'entendre et de voir les \cle{phrasal verbs} en action avant d'en analyser la structure.

\begin{textebox}[Radio Buea — “Youth Voice” : Interview with a student leader]
\textbf{Presenter:} Good morning, listeners. Today we talk about civic engagement. With me is Ako, the head boy of Government High School Buea. Ako, your class recently \textbf{carried out} a clean-up campaign. Why?

\textbf{Ako:} Good morning. Well, we noticed the rubbish heap near the market. People had \textbf{put up with} the smell for years. We decided to \textbf{stand up for} our right to a healthy environment. We \textbf{turned up} early Saturday, \textbf{looked into} the problem, and \textbf{sorted out} a plan.

\textbf{Presenter:} Did the authorities help?

\textbf{Ako:} The mayor promised to \textbf{look into} the illegal dumping. But we won't \textbf{give up} if they \textbf{put off} their decision. We want them to \textbf{set up} a recycling centre. We must \textbf{speak out} and \textbf{look after} our community. We can \textbf{bring about} change.
\end{textebox}

\begin{definition}[Glossaire de l'interview]
\begin{itemize}
\item \eng{civic engagement} (n.) : engagement civique, citoyen
\item \eng{rubbish heap} (n.) : tas d'ordures
\item \eng{illegal dumping} (n.) : dépôt sauvage (d'ordures)
\item \eng{recycling centre} (n.) : centre de recyclage
\item \eng{head boy} (n.) : élève major, délégué principal
\end{itemize}
\end{definition}

\begin{experience}[Activité d'écoute / lecture guidée]
\begin{enumerate}
\item \textbf{Écoutez / Lisez} le texte une première fois pour la compréhension globale : quel est le sujet ? Qui parle ? Quel est le ton ?
\item \textbf{Repérez} les mots en \textbf{gras} dans le texte. Combien y en a-t-il ? (Réponse : 11).
\item \textbf{Inférez} le sens de chaque expression en gras grâce au contexte. Exemple : \eng{carry out a campaign} $\neq$ \eng{carry} (porter) + \eng{out} (dehors) ; ici, cela signifie « mener, réaliser ».
\item \textbf{Classez}-les mentalement en deux groupes : ceux où le verbe et la particule semblent « collés » (\eng{look into}) et ceux où l'on pourrait insérer un mot entre les deux (\eng{carried out a campaign} $\to$ \eng{carried it out}).
\end{enumerate}
\end{experience}

\begin{aretenir}[Ce qu'on observe déjà]
\begin{itemize}
\item Un \cle{phrasal verb} = \textbf{verbe de base} + \textbf{particule} (adverbe ou préposition : \eng{up, out, off, into, with, for, after}…).
\item Le sens global est souvent \textbf{idiomatique} (imprévisible à partir des mots isolés) : \eng{give up} $\neq$ « donner en haut » mais « abandonner ».
\item Certains acceptent une séparation (objet nominal entre les deux), d'autres non.
\end{itemize}
\end{aretenir}

\coursec{Définition et formation}

\begin{definition}[Phrasal verb (verbe à particule)]
Un \textbf{phrasal verb} est une unité lexicale formée d'un \textbf{verbe} (généralement d'origine germanique, monosyllabique : \eng{go, put, take, look, give, set, turn, carry, bring, call, speak, stand}) suivi d'une \textbf{particule} (adverbe ou préposition) qui en modifie radicalement le sens.

\begin{center}
\begin{tabularx}{\linewidth}{|X|X|X|}\hline
\rowcolor{popblueL} \textbf{Verbe de base} & \textbf{Particule} & \textbf{Phrasal verb (sens idiomatique)} \\ \hline
\eng{look} (regarder) & \eng{after} & \eng{look after} = s'occuper de \\ \hline
\eng{put} (mettre) & \eng{off} & \eng{put off} = reporter \\ \hline
\eng{give} (donner) & \eng{up} & \eng{give up} = abandonner \\ \hline
\eng{bring} (apporter) & \eng{about} & \eng{bring about} = provoquer, entraîner \\ \hline
\end{tabularx}
\end{center}
\end{definition}

\begin{propriete}[Formation et grammaticalité]
\begin{itemize}
\item \textbf{Verbe} : se conjugue normalement (temps, mode, personne, voix). \emph{Ex. : \eng{carries out, carried out, will carry out, has carried out, is carrying out}.}
\item \textbf{Particule} : \textbf{invariable} (jamais de \eng{-s}, \eng{-ed}, \eng{-ing}). Elle reste identiquement \eng{up, out, off, into, with, for, after}…
\item \textbf{Transitivité} : le phrasal verb peut être intransitif (sans objet : \eng{The plane took off}) ou transitif (avec objet : \eng{They put off the meeting}).
\end{itemize}
\end{propriete}

\begin{attention}[Piège n°1 : traduction mot à mot]
Ne traduisez \textbf{jamais} le verbe et la particule séparément.
\begin{itemize}
\item \eng{He \textbf{turned up} late.} $\to$ « Il est arrivé en retard. » (Pas : « Il a tourné en haut ».)
\item \eng{She \textbf{looks after} her brother.} $\to$ « Elle s'occupe de son frère. » (Pas : « Elle regarde après son frère ».)
\end{itemize}
Apprenez le phrasal verb comme un \textbf{seul mot de vocabulaire} (unité de sens).
\end{attention}

\coursec{Les trois types de phrasal verbs transitifs}

C'est la distinction cruciale pour la grammaire (place de l'objet, pronom, passif). Nous distinguons trois catégories.

\begin{propriete}[Type 1 : Intransitifs — Sans objet]
Le phrasal verb ne prend pas de complément d'objet. Verbe et particule sont inséparables.
\begin{center}
\begin{tabular}{|l|l|l|}\hline
\rowcolor{popblueL} \textbf{Phrasal verb} & \textbf{Sens} & \textbf{Exemple} \\ \hline
\eng{turn up} & arriver, se présenter & \eng{Many voters \textbf{turned up} early.} \\ \hline
\eng{speak out} & s'exprimer publiquement & \eng{We must \textbf{speak out} against corruption.} \\ \hline
\eng{show up} & apparaître, arriver & \eng{He didn't \textbf{show up} at the meeting.} \\ \hline
\eng{come back} & revenir & \eng{When will she \textbf{come back}?} \\ \hline
\end{tabular}
\end{center}
\end{propriete}

\begin{propriete}[Type 2 : Transitifs séparables — Verbe + Particule (adverbe)]
L'objet (nominal) peut se placer \textbf{avant} ou \textbf{après} la particule. \textbf{Mais si l'objet est un pronom (\eng{it, them, me, us}…), il se place OBLIGATOIREMENT ENTRE le verbe et la particule.}
\begin{center}
\begin{tabular}{|l|l|l|}\hline
\rowcolor{popgreenL} \textbf{Phrasal verb} & \textbf{Sens} & \textbf{Exemples (Nom / Pronom)} \\ \hline
\eng{carry out} & mener, réaliser & \eng{carry out the plan / carry \textbf{it} out} \\ \hline
\eng{put off} & reporter & \eng{put off the meeting / put \textbf{it} off} \\ \hline
\eng{call off} & annuler & \eng{call off the strike / call \textbf{it} off} \\ \hline
\eng{set up} & créer, établir & \eng{set up a centre / set \textbf{it} up} \\ \hline
\eng{sort out} & résoudre, régler & \eng{sort out the problem / sort \textbf{it} out} \\ \hline
\eng{give up} & abandonner & \eng{give up smoking / give \textbf{it} up} \\ \hline
\eng{bring about} & provoquer & \eng{bring about change / bring \textbf{it} about} \\ \hline
\end{tabular}
\end{center}
\end{propriete}

\begin{propriete}[Type 3 : Transitifs inséparables — Verbe + Particule (préposition) ou Verbes à trois mots]
L'objet (nominal \textbf{ou} pronominal) se place \textbf{TOUJOURS APRÈS} la particule (ou la dernière particule pour les verbes à trois mots). Aucune séparation n'est possible.
\begin{center}
\begin{tabular}{|l|l|l|}\hline
\rowcolor{poppurpleL} \textbf{Phrasal verb} & \textbf{Sens} & \textbf{Exemples (Nom / Pronom)} \\ \hline
\eng{look after} & s'occuper de & \eng{look after the children / look after \textbf{them}} \\ \hline
\eng{look into} & enquêter sur & \eng{look into the case / look into \textbf{it}} \\ \hline
\eng{stand up for} & défendre (cause) & \eng{stand up for rights / stand up for \textbf{them}} \\ \hline
\eng{put up with} & tolérer, supporter & \eng{put up with noise / put up with \textbf{it}} \\ \hline
\eng{cope with} & faire face à & \eng{cope with the pressure / cope with \textbf{it}} \\ \hline
\eng{go for} & choisir, opter pour & \eng{go for the option / go for \textbf{it}} \\ \hline
\end{tabular}
\end{center}
\end{propriete}

\begin{methode}[Comment savoir si un phrasal verb est séparable ?]
\begin{enumerate}
\item \textbf{Test du pronom} : remplacez l'objet par \eng{it} (ou \eng{them}). Si la phrase naturelle est \eng{V + it + Part} $\to$ \textbf{Séparable}. Si c'est \eng{V + Part + it} $\to$ \textbf{Inséparable}.
\item \textbf{Dictionnaire} : les bons dictionnaires indiquent \eng{[sep]} ou \eng{[inseparable]} / \eng{[no obj]}.
\item \textbf{Règle pratique} : si la particule est une \textbf{préposition} introduisant un complément (\eng{after, into, with, for, on, at}…), le verbe est \textbf{inséparable}. Si la particule est un \textbf{adverbe} (\eng{up, out, off, on, down, away}…), le verbe est \textbf{séparable} (sauf exceptions comme \eng{go for}).
\end{enumerate}
\end{methode}

\begin{aretenir}[Résumé visuel des trois types]
\begin{center}
\begin{tikzpicture}[
  tbxone/.style={rectangle, rounded corners=4pt, draw=popblue, fill=popblueL, thick, align=center, minimum width=4.5cm, minimum height=2.8cm, font=\small},
  tbxtwo/.style={rectangle, rounded corners=4pt, draw=popgreen, fill=popgreenL, thick, align=center, minimum width=4.5cm, minimum height=2.8cm, font=\small},
  tbxthree/.style={rectangle, rounded corners=4pt, draw=poppurple, fill=poppurpleL, thick, align=center, minimum width=4.5cm, minimum height=2.8cm, font=\small},
  arr/.style={-{Stealth[length=3mm]}, thick, poporange}
]
\node[tbxone, align=center] (t1) at (0,0){\textbf{Type 1 : Intransitif}\\\eng{turn up, speak out, show up}\\\textit{Pas d'objet.}\\\textit{Jamais de séparation.}};
\node[tbxtwo, align=center] (t2) at (5.5,0){\textbf{Type 2 : Séparable}\\\eng{put off, carry out, set up}\\\textit{Objet nominal : V O Part \textbf{ou} V Part O.}\\\textit{Objet pronom : V \textbf{it} Part.}};
\node[tbxthree, align=center] (t3) at (11,0){\textbf{Type 3 : Inséparable}\\\eng{look after, look into, stand up for}\\\textit{Objet (nom/pronom) :}\\\textit{toujours APRÈS la particule.}};
\end{tikzpicture}
\end{center}
\end{aretenir}

\coursec{Phrasal verbs de la citoyenneté et des droits de l'homme (Chapitre 4)}

Le programme lie ce point de grammaire au thème « Citizenship and Human Rights ». Voici le lexique prioritaire à maîtriser pour la compréhension et la production (orale et écrite).

\begin{definition}[Lexique thématique — Phral verbs \& Citoyenneté]
\begin{center}
\begin{tabularx}{\linewidth}{|l|X|l|X|}\hline
\rowcolor{popblueL} \textbf{Phrasal verb} & \textbf{Sens (FR)} & \textbf{Type} & \textbf{Exemple contextualisé (Cameroun / Anglophone)} \\ \hline
\eng{stand up for} & défendre (droits, cause) & 3 mots & \eng{NGOs \textbf{stand up for} human rights in conflict zones.} \\ \hline
\eng{speak out} & s'exprimer, dénoncer & Insep. & \eng{Journalists \textbf{speak out} against censorship.} \\ \hline
\eng{look into} & enquêter sur & Insep. & \eng{The commission will \textbf{look into} the allegations.} \\ \hline
\eng{bring about} & provoquer, entraîner & Séparable & \eng{The new law \textbf{brought about} major changes.} \\ \hline
\eng{carry out} & mener, exécuter & Séparable & \eng{The youth \textbf{carried out} a voter registration drive.} \\ \hline
\eng{set up} & créer, mettre en place & Séparable & \eng{They \textbf{set up} a legal aid clinic in Yaoundé.} \\ \hline
\eng{put up with} & tolérer, supporter & 3 mots & \eng{Citizens should not \textbf{put up with} police brutality.} \\ \hline
\eng{give up} & abandonner & Séparable & \eng{Activists never \textbf{give up} their fight for justice.} \\ \hline
\eng{look after} & prendre soin de, gérer & Insep. & \eng{Community leaders \textbf{look after} vulnerable families.} \\ \hline
\eng{put off} & reporter, remettre à plus tard & Séparable & \eng{The election was \textbf{put off} due to the crisis.} \\ \hline
\eng{call off} & annuler & Séparable & \eng{The protest was \textbf{called off} after negotiations.} \\ \hline
\eng{sort out} & régler, résoudre & Séparable & \eng{We must \textbf{sort out} the land dispute peacefully.} \\ \hline
\eng{crack down on} & sévir contre, réprimer & 3 mots & \eng{The government \textbf{cracked down on} illegal logging.} \\ \hline
\eng{step down} & démissionner, se retirer & Intrans. & \eng{The minister \textbf{stepped down} after the scandal.} \\ \hline
\eng{back up} & soutenir, appuyer & Séparable & \eng{The UN \textbf{backed up} the peace process.} \\ \hline
\end{tabularx}
\end{center}
\end{definition}

\begin{exemplebox}[Contexte camerounais : phrases modèles]
\begin{enumerate}
\item \eng{The students \textbf{stood up for} the rights of internally displaced persons in Buea.}
\item \eng{The Human Rights Commission promised to \textbf{look into} the arbitrary arrests in Douala.}
\item \eng{We cannot \textbf{put up with} corruption in the public service any longer.}
\item \eng{The association \textbf{set up} a legal aid desk at the Court of First Instance in Yaoundé.}
\item \eng{If the authorities \textbf{put off} the dialogue, tensions will rise.}
\end{enumerate}
\end{exemplebox}

\coursec{Comparaison français/anglais et pièges des francophones}

\begin{propriete}[Verbe savant (latin) vs Phrasal verb (germanique) — Registre de langue]
L'anglais possède souvent deux verbes pour une même idée : un verbe « savant » (origine latine/française, formel, écrit) et un phrasal verb (origine germanique, neutre/standard, oral et écrit courant). Au Bac, on vous demande souvent de passer de l'un à l'autre.

\begin{center}
\begin{tabularx}{\linewidth}{|X|X|X|}\hline
\rowcolor{popblueL} \textbf{Verbe savant (Formel / Écrit soutenu)} & \textbf{Phrasal verb (Standard / Courant)} & \textbf{Français} \\ \hline
\eng{postpone} & \eng{put off} & reporter \\ \hline
\eng{tolerate / endure} & \eng{put up with} & tolérer, supporter \\ \hline
\eng{investigate / examine} & \eng{look into} & enquêter sur \\ \hline
\eng{cancel} & \eng{call off} & annuler \\ \hline
\eng{establish / create / found} & \eng{set up} & créer, établir \\ \hline
\eng{continue} & \eng{carry on} & continuer \\ \hline
\eng{abandon / renounce} & \eng{give up} & abandonner \\ \hline
\eng{cause / provoke / generate} & \eng{bring about} & provoquer, entraîner \\ \hline
\eng{resolve / solve} & \eng{sort out} & résoudre, régler \\ \hline
\eng{defend / support} & \eng{stand up for} & défendre \\ \hline
\eng{maintain / sustain} & \eng{keep up} & maintenir \\ \hline
\eng{discover / find out} & \eng{find out} & découvrir \\ \hline
\end{tabularx}
\end{center}
\emph{Note} : Les phrasal verbs ne sont pas « familiers » ou « argotiques ». Ils sont la norme en anglais standard (médias, administration, littérature, Bac). Les verbes savants sont réservés au style très formel (juridique, académique pur).
\end{propriete}

\begin{attention}[Piège n°2 : La place du pronom — L'erreur n°1 des francophones]
En français, le pronom objet \textbf{précède} le verbe : \emph{« Je l'ai reportée. »} \\
En anglais (verbe séparable), le pronom objet \textbf{s'intercale} : \eng{I \textbf{put it off}.} \\
\begin{center}
\begin{tabularx}{\linewidth}{|X|X|}\hline
\rowcolor{poporangeL} \textbf{Incorrect (Calque français)} & \textbf{Correct (Anglais)} \\ \hline
\eng{*They put off it.} & \eng{They \textbf{put it off}.} \\ \hline
\eng{*She carried out it.} & \eng{She \textbf{carried it out}.} \\ \hline
\eng{*We set up it yesterday.} & \eng{We \textbf{set it up} yesterday.} \\ \hline
\end{tabularx}
\end{center}
\emph{Astuce mnémotechnique} : « Le pronom est petit, il se glisse \textbf{au milieu} du verbe et de sa particule. »
\end{attention}

\begin{attention}[Piège n°3 : Verbes inséparables — Ne pas séparer !]
Avec \eng{look after, look into, stand up for, put up with, cope with}, le pronom reste \textbf{après} la particule.
\begin{center}
\begin{tabularx}{\linewidth}{|X|X|}\hline
\rowcolor{poporangeL} \textbf{Incorrect} & \textbf{Correct} \\ \hline
\eng{*She looks them after.} & \eng{She \textbf{looks after them}.} \\ \hline
\eng{*The police looked it into.} & \eng{The police \textbf{looked into it}.} \\ \hline
\eng{*We put it up with.} & \eng{We \textbf{put up with it}.} \\ \hline
\end{tabularx}
\end{center}
\end{attention}

\begin{attention}[Piège n°4 : Le passif — La particule reste collée au participe]
Au passif, le phrasal verb (séparable ou non) se comporte comme un bloc : \eng{be + participe passé + particule}. On ne sépare \textbf{jamais} la particule du participe.
\begin{center}
\begin{tabularx}{\linewidth}{|X|X|}\hline
\rowcolor{popgreenL} \textbf{Actif} & \textbf{Passif} \\ \hline
\eng{They \textbf{put off} the meeting.} & \eng{The meeting \textbf{was put off}.} \\ \hline
\eng{They \textbf{carried out} the plan.} & \eng{The plan \textbf{was carried out}.} \\ \hline
\eng{They \textbf{look after} the children.} & \eng{The children \textbf{are looked after}.} \\ \hline
\eng{They \textbf{looked into} the case.} & \eng{The case \textbf{was looked into}.} \\ \hline
\end{tabularx}
\end{center}
\end{attention}

\begin{attention}[Piège n°5 : Orthographe et particules confondues]
\begin{itemize}
\item \eng{put \textbf{off}} (reporter) $\neq$ \eng{put \textbf{of}} (n'existe pas). \textbf{Off} = deux \eng{f}.
\item \eng{give \textbf{up}} $\neq$ \eng{give \textbf{of}}.
\item \eng{look \textbf{after}} $\neq$ \eng{look \textbf{afterward}} (adverbe = par la suite).
\item Verbes irréguliers : \eng{set -- set -- set} (\eng{set up}), \eng{put -- put -- put} (\eng{put off, put up with}), \eng{bring -- brought -- brought} (\eng{bring about}), \eng{stand -- stood -- stood} (\eng{stand up for}), \eng{speak -- spoke -- spoken} (\eng{speak out}).
\end{itemize}
\end{attention}

\coursec{Entraînement guidé et production}

Après avoir comparé les phrasal verbs avec leurs équivalents français et repéré les pièges qui guettent les francophones, il est temps de passer à la pratique. Cette dernière partie de la leçon vous fait progresser en quatre étapes : \cle{repérer}, \cle{compléter}, \cle{transformer}, \cle{produire}. C'est exactement le cheminement exigé au Baccalauréat : on commence par reconnaître la structure dans un texte, on finit par l'utiliser soi-même dans une production écrite.

\begin{popfigure}
\begin{center}
\begin{tikzpicture}[
  etape/.style={rectangle, rounded corners=4pt, draw=popblue, fill=popblueL, thick, align=center, minimum width=2.6cm, minimum height=1.1cm, font=\small},
  fleche/.style={-{Stealth[length=3mm]}, thick, poporange}
]
\node[etape, align=center] (e1) at (0,0){\textbf{1. Repérer}\\ identifier dans un texte};
\node[etape, fill=popgreenL, draw=popgreen, align=center] (e2) at (4.2,0){\textbf{2. Compléter}\\ choisir la particule};
\node[etape, fill=poppurpleL, draw=poppurple, align=center] (e3) at (8.4,0){\textbf{3. Transformer}\\ verbe savant $\to$ phrasal};
\node[etape, fill=poporangeL, draw=poporange, align=center] (e4) at (12.6,0){\textbf{4. Produire}\\ rédiger un paragraphe};
\draw[fleche] (e1) -- (e2);
\draw[fleche] (e2) -- (e3);
\draw[fleche] (e3) -- (e4);
\end{tikzpicture}
\end{center}
\legende{La progression de l'entraînement : du repérage passif à la production active.}
\end{popfigure}

\courssub{Exercice 1 : repérer et traduire les phrasal verbs dans un texte}

Lisez d'abord le texte suivant, qui servira de support à l'exercice.

\begin{textebox}[Youth and civic engagement in Buea]
Last Saturday, the students of Government High School Buea decided to \textbf{carry out} a clean-up campaign in their neighbourhood. Early in the morning, they \textbf{turned up} in front of the town hall, wearing gloves and carrying bags. The mayor had promised to \textbf{look into} the problem of illegal dumping, and the students wanted to show that young people can \textbf{stand up for} a clean environment.

During the campaign, an old woman came to thank them. She explained that her family had always \textbf{put up with} the bad smell from the rubbish heap near her house, because nobody had ever tried to \textbf{sort out} the situation. The students promised they would not \textbf{give up} until the heap disappeared. They also asked the town council to \textbf{set up} a recycling centre.

At the end of the day, the head boy declared: “We must \textbf{speak out} against injustice and \textbf{look after} our community. If the authorities \textbf{put off} their decisions, citizens must act.” The local radio reporter noted that the students had managed to \textbf{bring about} a real change in people's attitudes.
\end{textebox}

\begin{definition}[Glossaire du texte]
\begin{itemize}
\item \eng{clean-up campaign} (n.) : campagne de nettoyage
\item \eng{dumping} (n.) : dépôt sauvage d'ordures
\item \eng{rubbish heap} (n.) : tas d'ordures
\item \eng{town council} (n.) : conseil municipal
\item \eng{head boy} (n.) : élève major / délégué principal
\item \eng{attitude} (n.) : attitude, état d'esprit
\end{itemize}
\end{definition}

\begin{exercice}[Exercice 1 — Repérage et traduction]
Soulignez les phrasal verbs du texte ci-dessus, puis traduisez chacun d'eux en français, en vous aidant du contexte.
\end{exercice}

\begin{corrige}[Exercice 1]
\begin{center}
\begin{tabularx}{\linewidth}{|l|X|X|}\hline
\rowcolor{popblueL} \textbf{Phrasal verb} & \textbf{Traduction} & \textbf{Type} \\ \hline
carry out & mener, réaliser & séparable \\ \hline
turn up & arriver, se présenter & intransitif / inséparable \\ \hline
look into & examiner, enquêter sur & inséparable \\ \hline
stand up for & défendre (une cause) & à trois mots \\ \hline
put up with & tolérer, supporter & à trois mots \\ \hline
sort out & régler, résoudre & séparable \\ \hline
give up & abandonner & séparable \\ \hline
set up & créer, mettre en place & séparable \\ \hline
speak out & s'exprimer publiquement & intransitif / inséparable \\ \hline
look after & s'occuper de & inséparable \\ \hline
put off & reporter & séparable \\ \hline
bring about & provoquer, entraîner & séparable \\ \hline
\end{tabularx}
\end{center}
\emph{Astuce de repérage} : un phrasal verb se reconnaît à sa particule (\eng{up, out, off, into, with, for, after}) qui modifie le sens du verbe de base. Comparez \eng{to look} (regarder) et \eng{to look into} (enquêter sur) : le sens change complètement.
\end{corrige}

\courssub{Exercice 2 : compléter avec la particule correcte}

\begin{methode}[Comment choisir la bonne particule]
\begin{enumerate}
\item Lisez la phrase entière et devinez le sens général.
\item Demandez-vous quel sens le verbe seul ne peut pas exprimer.
\item Testez mentalement chaque particule : c'est le \textbf{sens global} qui décide, pas une traduction mot à mot.
\item Vérifiez la conjugaison du verbe (temps, personne) : la particule, elle, ne change jamais.
\end{enumerate}
\end{methode}

\begin{exercice}[Exercice 2 — La bonne particule]
Complétez avec \eng{up}, \eng{out}, \eng{off}, \eng{after}, \eng{into} ou \eng{with}.
\begin{enumerate}
\item The police promised to look \ldots{} the disappearance of the documents.
\item We cannot put \ldots{} with corruption any longer.
\item The demonstration was called \ldots{} because of the heavy rain.
\item Who will look \ldots{} the orphans after the crisis?
\item The association carried \ldots{} a survey on children's rights.
\item Many citizens did not turn \ldots{} to vote on election day.
\end{enumerate}
\end{exercice}

\begin{corrige}[Exercice 2]
\begin{enumerate}
\item \eng{look \textbf{into}} = enquêter sur. « La police a promis d'enquêter sur la disparition des documents. »
\item \eng{put \textbf{up} with} = tolérer. « Nous ne pouvons plus tolérer la corruption. »
\item \eng{called \textbf{off}} = annulée. « La manifestation a été annulée à cause de la forte pluie. »
\item \eng{look \textbf{after}} = s'occuper de. « Qui s'occupera des orphelins après la crise ? »
\item \eng{carried \textbf{out}} = menée. « L'association a mené une enquête sur les droits des enfants. »
\item \eng{turn \textbf{up}} = se présenter. « Beaucoup de citoyens ne se sont pas présentés pour voter. »
\end{enumerate}
\end{corrige}

\courssub{Exercice 3 : remplacer le complément nominal par un pronom}

Rappel de la règle étudiée dans la section 3 : avec un phrasal verb \cle{séparable}, un complément \emph{nominal} peut se placer avant ou après la particule, mais un complément \emph{pronom} se place \textbf{obligatoirement entre le verbe et la particule}.

\begin{exemplebox}[La règle du pronom en images]
\eng{The authorities postponed the meeting.} $\to$ \eng{The authorities put the meeting off.} $\to$ \eng{The authorities put \textbf{it} off.}\\
« Les autorités ont reporté la réunion. » $\to$ « Les autorités l'ont reportée. »\\[0.3em]
Jamais : \eng{put off it}. C'est l'inverse du français, où le pronom précède le verbe (\emph{ils l'ont reportée}).
\end{exemplebox}

\begin{exoresolu}[Exercice 3 — Le pronom à sa place]
Énoncé : réécrivez chaque phrase en remplaçant le complément souligné par un pronom (\eng{it}, \eng{them}, \eng{her}).
\begin{enumerate}
\item The students carried out \underline{the campaign}.
\item The mayor promised to sort out \underline{the problem}.
\item We must give up \underline{our bad habits}.
\item The council set up \underline{a recycling centre}.
\end{enumerate}
\tcblower
Solution détaillée :
\begin{enumerate}
\item \eng{The students carried \textbf{it} out.} — \eng{carry out} est séparable ; le pronom \eng{it} (la campagne) se place entre verbe et particule. « Les élèves l'ont menée. »
\item \eng{The mayor promised to sort \textbf{it} out.} — le pronom suit l'infinitif \eng{to sort} et précède \eng{out}. « Le maire a promis de le régler. »
\item \eng{We must give \textbf{them} up.} — \eng{habits} est pluriel, donc \eng{them}. « Nous devons y renoncer. »
\item \eng{The council set \textbf{it} up.} — attention au passé : \eng{set} est irrégulier (set--set--set). « Le conseil l'a créé. »
\end{enumerate}
\end{exoresolu}

\begin{attention}[Piège : les verbes inséparables]
Avec un phrasal verb \cle{inséparable} (\eng{look after}, \eng{look into}, \eng{cope with}), le pronom se place \textbf{après la particule}, jamais entre : \eng{She looks after \textbf{them}.} « Elle s'occupe d'eux. » — et non \eng{looks them after}. Avant de placer le pronom, demandez-vous donc toujours : ce phrasal verb est-il séparable ou inséparable ?
\end{attention}

\courssub{Exercice 4 : transformation — du verbe savant au phrasal verb}

C'est l'exercice type du Baccalauréat : on vous donne une phrase avec un verbe « savant » (souvent d'origine latine, proche du français) et vous devez le remplacer par le phrasal verb équivalent.

\begin{methode}[Réussir les exercices de transformation au Bac]
\begin{enumerate}
\item \textbf{Repérez} le verbe à remplacer et identifiez son temps et sa voix (active ou passive).
\item \textbf{Choisissez} le phrasal verb équivalent : \eng{postpone} $\to$ \eng{put off} ; \eng{tolerate} $\to$ \eng{put up with} ; \eng{investigate} $\to$ \eng{look into} ; \eng{cancel} $\to$ \eng{call off} ; \eng{establish} $\to$ \eng{set up} ; \eng{continue} $\to$ \eng{carry on}.
\item \textbf{Conjuguez} le verbe du phrasal verb au même temps et à la même personne que le verbe d'origine (la particule reste invariable).
\item \textbf{Vérifiez} la place du complément : nom avant ou après la particule, pronom obligatoirement au milieu (verbe séparable) ou après la particule (verbe inséparable).
\end{enumerate}
\end{methode}

\begin{center}
\begin{tabularx}{\linewidth}{|X|X|X|}\hline
\rowcolor{popblueL} \textbf{Verbe savant} & \textbf{Phrasal verb} & \textbf{Français} \\ \hline
postpone & put off & reporter \\ \hline
tolerate & put up with & tolérer \\ \hline
investigate & look into & enquêter sur \\ \hline
cancel & call off & annuler \\ \hline
establish / create & set up & créer \\ \hline
continue & carry on & continuer \\ \hline
abandon & give up & abandonner \\ \hline
cause / provoke & bring about & provoquer \\ \hline
solve / resolve & sort out & résoudre \\ \hline
defend & stand up for & défendre \\ \hline
\end{tabularx}
\end{center}

\begin{exoresolu}[Exercice 4 — Transformation guidée]
Énoncé : réécrivez les phrases en remplaçant le verbe souligné par un phrasal verb, sans changer le sens.
\begin{enumerate}
\item The judge \underline{postponed} the trial.
\item The commission will \underline{investigate} the case.
\item The villagers \underline{tolerated} the injustice for years.
\end{enumerate}
\tcblower
Solution détaillée :
\begin{enumerate}
\item \eng{The judge \textbf{put off} the trial} ou \eng{The judge \textbf{put the trial off}.} — \eng{postponed} est au prétérit ; \eng{put} est irrégulier (put--put--put), donc \eng{put} au prétérit. « Le juge a reporté le procès. »
\item \eng{The commission will \textbf{look into} the case.} — futur : \eng{will} + base verbale ; \eng{look into} est inséparable, le complément reste après. « La commission enquêtera sur l'affaire. »
\item \eng{The villagers \textbf{put up with} the injustice for years.} — prétérit : \eng{put} ; verbe à trois mots, jamais de séparation. « Les villageois ont toléré l'injustice pendant des années. »
\end{enumerate}
\end{exoresolu}

\begin{propriete}[Le phrasal verb au passif]
Au passif, le phrasal verb garde sa particule \textbf{accolée au verbe} ; on ne la sépare jamais de \eng{be} + participe passé :
\begin{itemize}
\item \eng{The meeting was \textbf{put off}.} « La réunion a été reportée. »
\item \eng{The victims are being \textbf{looked after}.} « Les victimes sont prises en charge. »
\item \eng{The case will be \textbf{looked into}.} « L'affaire sera examinée. »
\end{itemize}
Remarquez que même un verbe inséparable (\eng{look after}) voit sa particule se coller au participe passé : \eng{are being looked after}.
\end{propriete}

\courssub{Exercice 5 : production écrite — « An active citizen in my community »}

\begin{attention}[Deux conditions pour chaque phrasal verb]
Dans une production écrite, chaque phrasal verb doit être \textbf{correctement conjugué} (temps, personne, forme irrégulière) \emph{et} \textbf{correctement placé} (particule à sa place, pronom au milieu si le verbe est séparable). Une particule oubliée ou mal placée est comptée comme une faute, même si le verbe est juste. Écrire \eng{He put with the noise} au lieu de \eng{He put up with the noise} coûte des points.
\end{attention}

\begin{exercice}[Exercice 5 — Production guidée]
Rédigez un paragraphe de 5 à 6 lignes intitulé \eng{“An active citizen in my community”}. Utilisez \textbf{au moins cinq phrasal verbs} de la leçon, par exemple : \eng{stand up for}, \eng{carry out}, \eng{look after}, \eng{speak out}, \eng{give up}, \eng{set up}, \eng{put off}, \eng{bring about}. Soulignez-les dans votre copie.
\end{exercice}

\begin{exemplebox}[Modèle de production (corrigé type)]
\eng{In my community in Bamenda, Mrs Ngong is a \textbf{truly} active citizen. She \textbf{stands up for} the rights of widows and orphans, and she never \textbf{gives up}, even when people discourage her. Last year, she \textbf{carried out} a campaign to \textbf{set up} a small library for the children of her quarter. When the landlord tried to \textbf{put off} the project, she \textbf{spoke out} at the town meeting and convinced everybody. Today she \textbf{looks after} the library with the help of volunteers, and she has \textbf{brought about} a real change: children now love reading.}\\[0.3em]
« Dans ma communauté à Bamenda, Mme Ngong est une vraie citoyenne active. Elle défend les droits des veuves et des orphelins et n'abandonne jamais, même quand on la décourage. L'an dernier, elle a mené une campagne pour créer une petite bibliothèque pour les enfants de son quartier. Quand le propriétaire a essayé de reporter le projet, elle a pris la parole à la réunion municipale et a convaincu tout le monde. Aujourd'hui, elle s'occupe de la bibliothèque avec l'aide de bénévoles et elle a provoqué un vrai changement : les enfants aiment désormais lire. »
\end{exemplebox}

\courssub{Correction collective : justifier chaque réponse par la règle}

La correction en classe ne consiste pas seulement à donner la bonne réponse : chaque réponse doit être \cle{justifiée} par la règle correspondante. Voici la démarche à adopter oralement.

\begin{methode}[Justifier une réponse en trois phrases]
\begin{enumerate}
\item \textbf{Nommer la règle} : « \eng{put off} is a separable phrasal verb. »
\item \textbf{Appliquer la règle} : « The object is a pronoun, so it must come between the verb and the particle. »
\item \textbf{Produire la phrase correcte} : « The authorities put it off. »
\end{enumerate}
\end{methode}

\begin{experience}[Correction collective interactive]
L'enseignant écrit au tableau cinq réponses d'élèves (certaines volontairement fautives : \eng{put off it}, \eng{look them after}, \eng{he putted off}). Par groupes de trois, les élèves :
\begin{enumerate}
\item identifient la faute éventuelle ;
\item citent la règle violée (séparation, conjugaison, place du pronom) ;
\item proposent la correction à voix haute.
\end{enumerate}
Chaque groupe gagne un point par faute corrigée \emph{avec} sa justification. L'activité se termine par la relecture collective des paragraphes de l'exercice 5 : les camarades vérifient que les cinq phrasal verbs sont bien conjugués et bien placés.
\end{experience}

\begin{aretenir}[L'essentiel de l'entraînement]
\begin{itemize}
\item Repérer un phrasal verb = verbe + particule dont le sens diffère du verbe seul.
\item Pronom complément : \textbf{au milieu} si le verbe est séparable (\eng{put it off}), \textbf{après} s'il est inséparable (\eng{look after them}).
\item Transformation au Bac : même temps, même voix, particule invariable, complément bien placé.
\item Au passif, la particule reste collée au participe : \eng{The meeting was put off.}
\item En production : chaque phrasal verb doit être conjugué \emph{et} placé correctement — une particule oubliée est une faute.
\end{itemize}
\end{aretenir}

\newpage

\coursec{Activité d'intégration}

\courssub{Objectifs de l'activité}

À l'issue de cette activité d'intégration, l'élève sera capable de :

\begin{itemize}
\item utiliser spontanément et correctement au moins \cle{six phrasal verbs} de la leçon dans un échange oral authentique ;
\item respecter la syntaxe des phrasal verbs : place du pronom (\eng{vote them out}, jamais \eng{vote out them}), inséparabilité de certains verbes (\eng{look into the matter}), constructions à trois mots (\eng{put up with injustice}) ;
\item mobiliser le lexique de la citoyenneté et des droits de l'homme (\eng{rights, corruption, election, community, justice}) ;
\item tenir un rôle défini (animateur, activiste, chef traditionnel, jeune électeur) avec le registre de langue approprié ;
\item écouter attentivement ses camarades et repérer les phrasal verbs sur une grille d'écoute ;
\item rédiger, en variante écrite, un court éditorial de clôture réutilisant les phrasal verbs dans un texte cohérent.
\end{itemize}

\courssub{Situation}

La CRTV, la radio nationale camerounaise, organise une émission hebdomadaire intitulée \emph{“Stand up for your rights!”}. Cette semaine, l'émission est enregistrée en public dans votre lycée, à Yaoundé. Le sujet du jour : un problème local qui touche directement la communauté — la corruption dans l'attribution des marchés publics, les droits des femmes, ou l'accès à l'eau potable dans le quartier. Trois invités sont attendus en studio : un activiste des droits de l'homme, un chef traditionnel et un jeune électeur qui votera pour la première fois. Votre groupe de quatre élèves incarne l'animateur et les trois invités.

\begin{textebox}[CRTV programme announcement — “Stand up for your rights!”]
This Saturday at 6 p.m., tune in to CRTV for a special edition of “Stand up for your rights!”, recorded live at a secondary school in Yaoundé. Our presenter will speak out with three guests about a problem that concerns us all: access to clean water in our neighbourhoods. A human rights activist will explain why citizens must never give in to injustice. A traditional chief will tell us how his community looks after its only borehole. And a young first-time voter will ask why corrupt officials are not voted out at election time. Should we put up with broken promises, or should we carry on the struggle for change? The authorities have promised to look into the matter — but will they? Don't give in! Carry on the struggle! Join the debate this Saturday on CRTV.
\end{textebox}

\begin{definition}[Glossaire du document]
\begin{itemize}
\item \eng{tune in} (phrasal verb) : se brancher sur une station, écouter ;
\item \eng{borehole} (nom) : forage, puits équipé d'une pompe ;
\item \eng{first-time voter} (nom) : électeur qui vote pour la première fois ;
\item \eng{broken promises} (nom pluriel) : promesses non tenues ;
\item \eng{the authorities} (nom pluriel) : les autorités.
\end{itemize}
\end{definition}

\courssub{Consignes}

\begin{enumerate}
\item \textbf{Compréhension du document.} Lis l'annonce de la CRTV. Souligne tous les phrasal verbs que tu reconnais et classe-les en trois colonnes : séparables, inséparables, à trois mots. Vérifie ton classement avec ton groupe.
\item \textbf{Choix du sujet.} En groupe de quatre, choisissez un seul problème local pour votre émission : \eng{corruption in public contracts}, \eng{women's rights in the village} ou \eng{access to clean water}. Distribuez les rôles : presenter, human rights activist, traditional chief, young voter.
\item \textbf{Préparation du lexique et de la grammaire.} Chaque intervenant choisit au moins \cle{six phrasal verbs} de la leçon et rédige ses répliques. Attention à la syntaxe : avec un pronom, dis \eng{we must vote them out} ; dis \eng{the council must look into the problem} (inséparable) ; dis \eng{we cannot put up with this situation} (trois mots).
\item \textbf{Répétition.} Répétez l'émission une première fois. L'animateur veille à poser des questions ouvertes (\eng{What do you think of...?}, \eng{How does your community deal with...?}) et termine obligatoirement par l'appel aux auditeurs : \eng{Don't give in! Carry on the struggle!}
\item \textbf{Performance.} Jouez l'émission devant la classe (5 à 7 minutes). Parlez sans lire intégralement votre texte : des notes sont autorisées, pas un script lu mot à mot.
\item \textbf{Écoute active.} Pendant la prestation des autres groupes, coche sur ta grille d'écoute chaque phrasal verb correctement utilisé (verbe + particule + place correcte du complément). Note aussi toute erreur entendue pour en débattre ensuite.
\item \textbf{Variante écrite.} Rédige individuellement l'éditorial de clôture de l'émission en 8 à 10 lignes : résume le débat, donne ton opinion et termine par un appel à l'action, en employant au moins cinq phrasal verbs.
\end{enumerate}

\courssub{Ressources à mobiliser}

\begin{aretenir}[Phrasal verbs utiles pour le débat]
\begin{center}
\begin{tabularx}{\linewidth}{|l|X|l|}\hline
\rowcolor{popblueL} \textbf{Phrasal verb} & \textbf{Signification} & \textbf{Type} \\ \hline
\eng{stand up for (rights)} & défendre (ses droits) & trois mots, inséparable \\ \hline
\eng{speak out (against)} & dénoncer publiquement & inséparable \\ \hline
\eng{give in} & céder, abandonner & inséparable \\ \hline
\eng{carry on (the struggle)} & continuer (la lutte) & séparable \\ \hline
\eng{look after (the borehole)} & s'occuper de, entretenir & inséparable \\ \hline
\eng{put up with (injustice)} & tolérer, supporter & trois mots, inséparable \\ \hline
\eng{vote (them) out} & renvoyer par le vote & séparable \\ \hline
\eng{look into (the matter)} & enquêter sur, examiner & inséparable \\ \hline
\eng{deal with (a problem)} & s'occuper d'un problème & inséparable \\ \hline
\eng{bring about (change)} & provoquer (le changement) & séparable \\ \hline
\end{tabularx}
\end{center}
\end{aretenir}

\begin{propriete}[Rappel des règles essentielles]
\begin{itemize}
\item \textbf{Verbe séparable + pronom} : le pronom se place obligatoirement entre le verbe et la particule : \eng{we will vote them out} — jamais \eng{vote out them}.
\item \textbf{Verbe inséparable} : le complément suit toujours la particule : \eng{they promised to look into the case}.
\item \textbf{Verbe à trois mots} : jamais de coupure : \eng{we must stand up for our rights}, \eng{we cannot put up with corruption}.
\item \textbf{Impératif négatif} : \eng{Don't give in!} ; \textbf{impératif affirmatif} : \eng{Carry on the struggle!} ou \eng{Carry it on!}
\end{itemize}
\end{propriete}

\begin{attention}[Pièges des francophones]
Ne traduis pas mot à mot : \eng{stand up for} ne signifie pas « se lever » mais « défendre » ; \eng{give in} signifie « céder » et non « donner dedans ». N'ajoute pas de préposition française : on dit \eng{look into the matter} (et non \eng{look into on the matter}). Enfin, la particule change tout le sens : compare \eng{look after} (s'occuper de), \eng{look for} (chercher) et \eng{look into} (enquêter sur).
\end{attention}

\courssub{Proposition de production et corrigé commenté}

\begin{exoresolu}[Modèle d'extrait d'émission et d'éditorial de clôture]
Énoncé : rédige l'éditorial de clôture de l'émission (8 à 10 lignes) en employant au moins cinq phrasal verbs de la leçon.
\tcblower
\textbf{Extrait du débat (modèle oral)} :

\eng{Presenter: Good evening, dear listeners. Tonight we are talking about clean water in our neighbourhoods. Mr Ndi, as a human rights activist, what is your message?}

\eng{Activist: My message is simple: we must stand up for our rights. For too long, we have put up with empty promises. We will not give in.}

\eng{Chief: In my village, we look after our borehole together. When it breaks down, we do not wait for the government; we deal with the problem ourselves.}

\eng{Young voter: And when officials steal public money, we must vote them out at the next election. The authorities say they will look into the matter, but we want action, not words.}

\eng{Presenter: Thank you all. Dear listeners, don't give in! Carry on the struggle!}

\textbf{Éditorial de clôture (production écrite modèle)} :

\eng{Tonight's debate has shown that our community refuses to put up with injustice any longer. Our guests agreed on one point: we must stand up for our rights and speak out against corruption. The traditional chief reminded us that we can look after our common resources ourselves, while the young voter insisted that dishonest leaders must be voted out. The authorities have promised to look into the water crisis, but promises are not enough. Real change will only come about if we carry on the struggle together. So, dear listeners, don't give in: stand up, speak out, and carry on!}

\textbf{Commentaire des choix de langue} :
\begin{itemize}
\item Sept phrasal verbs sont employés : \eng{put up with}, \eng{stand up for}, \eng{speak out}, \eng{look after}, \eng{vote out}, \eng{look into}, \eng{carry on}, \eng{come about}.
\item Syntaxe respectée : \eng{must be voted out} (forme passive correcte du verbe séparable), \eng{look into the water crisis} (inséparable, complément après la particule), \eng{put up with injustice} (trois mots, jamais coupé).
\item Registre adapté à un éditorial : phrases courtes, impératifs finaux (\eng{stand up, speak out, carry on}) qui créent un appel à l'action rythmé.
\item Lexique du chapitre réinvesti : \eng{rights, corruption, community, leaders, authorities, change}.
\end{itemize}
\end{exoresolu}

\courssub{Bilan de l'activité}

Cette activité a permis de mobiliser l'ensemble des compétences de la leçon dans une tâche finale authentique. D'un point de vue grammatical, l'élève a manipulé les trois types de phrasal verbs — séparables (\eng{vote them out}), inséparables (\eng{look into the matter}) et à trois mots (\eng{stand up for our rights}) — dans un contexte de communication réel, ce qui favorise la mémorisation durable. D'un point de vue lexical, le vocabulaire de la citoyenneté et des droits de l'homme a été réemployé dans un registre oral formel (débat radiophonique) puis écrit (éditorial). La grille d'écoute a développé la compréhension orale et l'attention à la forme : repérer un phrasal verb correctement construit chez autrui renforce sa propre compétence. Enfin, le contexte camerounais (CRTV, chef traditionnel, forage, jeune électeur) ancre l'apprentissage dans la réalité des élèves et illustre concrètement le thème du chapitre : un citoyen informé est un citoyen qui sait \eng{stand up for his rights} — et qui, surtout, ne \eng{gives in} jamais.

\newpage

\coursec{Méthodes et exercices résolus}

\courssub{Savoir-faire 1 : Identifier un phrasal verb dans un texte}

Un \cle{phrasal verb} (verbe à particule) est un groupe formé d'un \textbf{verbe} et d'une ou deux \textbf{particules} dont le sens global est souvent très différent du sens du verbe seul. C'est l'une des structures les plus fréquentes de l'anglais parlé et écrit, et l'une des plus redoutées des francophones, car le français n'a pas d'équivalent : il faut apprendre chaque phrasal verb comme un mot de vocabulaire à part entière.

\begin{methode}[Comment repérer et comprendre un phrasal verb]
\begin{enumerate}
\item Cherche dans le texte un \cle{verbe} suivi d'une \cle{particule} (adverbe : \eng{up, down, on, off, out, in, away, back} ; ou préposition : \eng{for, after, into, with}).
\item Demande-toi si le sens du groupe est \textbf{différent} du sens du verbe seul. \eng{The prisoner gave up.} ne veut pas dire « donner vers le haut » mais « abandonner » : c'est un phrasal verb.
\item Teste avec un pronom : si l'on peut dire \eng{gave it up}, le verbe est \textbf{séparable} ; si l'on doit dire \eng{look after her} (jamais \eng{look her after}), il est \textbf{inséparable}.
\item Retrouve l'équivalent français (souvent un verbe simple : \eng{give up} = abandonner, \eng{look after} = s'occuper de) ou un verbe anglais plus soutenu (\eng{put off} = \eng{postpone}).
\item Note le phrasal verb dans ton carnet avec la particule, un exemple et la traduction.
\end{enumerate}
\end{methode}

\begin{exoresolu}[Repérage dans un texte]
\textbf{Énoncé.} Read the sentence and answer the questions : \eng{After the election, the new mayor promised to look into the complaints of the citizens and to do away with corruption in the council.}\\
a) Underline the two phrasal verbs.\\
b) Give their meaning in French.\\
c) Replace them with more formal single verbs.
\tcblower
\textbf{Solution détaillée.}\\
a) Les deux groupes verbe + particule dont le sens dépasse le verbe simple sont : \eng{look into} (verbe + préposition \eng{into}) et \eng{do away with} (phrasal verb à trois mots).\\
b) \eng{look into the complaints} = examiner, enquêter sur les plaintes ; \eng{do away with corruption} = supprimer, éliminer la corruption.\\
c) Équivalents soutenus : \eng{look into} = \eng{investigate} ; \eng{do away with} = \eng{abolish} ou \eng{eliminate}.\\
\textbf{Réponse rédigée :} \eng{The two phrasal verbs are “to look into”, which means to investigate, and “to do away with”, which means to abolish. The mayor promised to investigate the citizens' complaints and to abolish corruption.}\\
\textbf{Conclusion :} un phrasal verb se reconnaît à son sens global, souvent très différent du verbe de base ; il a presque toujours un équivalent formel en un seul mot.
\end{exoresolu}

\courssub{Savoir-faire 2 : Construire correctement la phrase (séparable ou inséparable)}

La grande difficulté des phrasal verbs est la \textbf{place du complément}, surtout quand c'est un pronom. Il faut d'abord classer le verbe dans l'une des trois catégories, puis appliquer la règle de position correspondante.

\begin{propriete}[Les trois types de phrasal verbs]
\begin{enumerate}
\item \textbf{Séparables} (\eng{give up, put off, turn down, bring up}) : avec un \textbf{nom}, deux positions sont possibles : \eng{She gave up smoking} ou \eng{She gave smoking up}. Avec un \textbf{pronom}, la séparation est \textbf{obligatoire} : \eng{She gave it up.}
\item \textbf{Inséparables} (\eng{look after, look into, deal with, stand for}) : la particule est une préposition et reste toujours collée au verbe : \eng{They look after the children} / \eng{They look after them.}
\item \textbf{À trois mots} (\eng{put up with, look forward to, do away with}) : jamais de coupure ; le complément suit tout le groupe : \eng{We look forward to it.}
\end{enumerate}
\end{propriete}

\begin{center}
\begin{tabularx}{\linewidth}{|X|X|X|}
\hline
\rowcolor{popblueL} Type & Avec un nom & Avec un pronom \\
\hline
Séparable : \eng{turn down} (rejeter) & \eng{He turned down the offer.} / \eng{He turned the offer down.} & \eng{He turned it down.} (jamais \eng{turned down it}) \\
\hline
Inséparable : \eng{look after} (s'occuper de) & \eng{She looks after the orphans.} & \eng{She looks after them.} (jamais \eng{looks them after}) \\
\hline
Trois mots : \eng{put up with} (tolérer) & \eng{We cannot put up with injustice.} & \eng{We cannot put up with it.} \\
\hline
\end{tabularx}
\end{center}

\begin{methode}[Placer le complément et le pronom au bon endroit]
\begin{enumerate}
\item \textbf{Verbe séparable} : deux positions possibles avec un nom : \eng{She gave up smoking} ou \eng{She gave smoking up}.
\item Mais avec un \cle{pronom} (\eng{it, them, him}), la séparation est \textbf{obligatoire} : \eng{She gave it up.} — jamais \eng{gave up it}.
\item \textbf{Verbe inséparable} : la particule reste collée au verbe : \eng{They look after the children} / \eng{They look after them.}
\item \textbf{Verbe à trois mots} : jamais de coupure : \eng{We look forward to it.}
\item Réflexe d'examen : si le complément est \eng{it} ou \eng{them}, demande-toi d'abord si le verbe est séparable ; si oui, place le pronom au milieu.
\end{enumerate}
\end{methode}

\begin{exoresolu}[Transformation avec un pronom]
\textbf{Énoncé.} Rewrite each sentence, replacing the underlined noun with a pronoun.\\
1. \eng{The government must do away with \underline{discrimination}.}\\
2. \eng{The judge turned down \underline{the appeal}.}\\
3. \eng{Our association looks after \underline{street children} in Douala.}
\tcblower
\textbf{Solution détaillée.}\\
1. \eng{do away with} est un verbe à trois mots, inséparable : le pronom se place après tout le groupe. \eng{The government must do away with it.} (discrimination = singulier neutre $\rightarrow$ \eng{it}).\\
2. \eng{turn down} (rejeter) est séparable : avec un pronom, la coupure est obligatoire. \eng{The judge turned it down.} (the appeal = singulier $\rightarrow$ \eng{it}).\\
3. \eng{look after} est inséparable : \eng{Our association looks after them in Douala.} (street children = pluriel $\rightarrow$ \eng{them}).\\
\textbf{Conclusion :} la position du pronom est le test le plus fiable : milieu pour les séparables, fin pour les inséparables et les verbes à trois mots.
\end{exoresolu}

\courssub{Savoir-faire 3 : Employer les phrasal verbs du thème « citoyenneté et droits »}

Dans le chapitre \emph{Citizenship and Human Rights}, certains phrasal verbs reviennent constamment dans les textes et les sujets de composition. Les voici réunis avec leur construction et leur traduction.

\begin{definition}[Phrasal verbs de la citoyenneté et des droits de l'homme]
\begin{itemize}
\item \eng{stand up for + nom} : défendre (une cause, des droits). \eng{We must stand up for women's rights.} « Nous devons défendre les droits des femmes. »
\item \eng{speak out against + nom} : dénoncer, s'élever contre. \eng{She spoke out against corruption.} « Elle a dénoncé la corruption. »
\item \eng{look into + nom} : enquêter sur, examiner. \eng{The commission is looking into the case.} « La commission enquête sur l'affaire. »
\item \eng{do away with + nom} : supprimer, abolir. \eng{They want to do away with child labour.} « Ils veulent abolir le travail des enfants. »
\item \eng{put up with + nom} : tolérer, supporter. \eng{We will not put up with injustice.} « Nous ne tolérerons pas l'injustice. »
\item \eng{look after + nom} : s'occuper de. \eng{The centre looks after street children.} « Le centre s'occupe des enfants des rues. »
\item \eng{abide by + nom} : respecter (une loi, une règle). \eng{Every citizen must abide by the law.} « Tout citoyen doit respecter la loi. »
\item \eng{take part in + nom} : participer à. \eng{Many students took part in the march.} « Beaucoup d'élèves ont participé à la marche. »
\end{itemize}
\end{definition}

\begin{methode}[Choisir le bon phrasal verb selon le contexte]
\begin{enumerate}
\item Identifie l'idée exprimée en français : se mobiliser, défendre = \eng{stand up for} ; dénoncer = \eng{speak out against} ; enquêter = \eng{look into} ; supprimer = \eng{do away with} ; tolérer = \eng{put up with} ; s'occuper de = \eng{look after} ; respecter une loi = \eng{abide by} ; participer = \eng{take part in}.
\item Vérifie la construction : \eng{stand up for + nom} (défendre), \eng{speak out against + nom} (s'élever contre) — attention aux prépositions contraires (\eng{for} / \eng{against}).
\item Conjugue le verbe normalement (la particule ne change jamais) : \eng{stood up for} (prétérit de \eng{stand} : stand -- stood -- stood), \eng{spoke out against} (prétérit de \eng{speak} : speak -- spoke -- spoken).
\item Au Bac, dans une composition sur les droits de l'homme, utilise 2 ou 3 phrasal verbs bien maîtrisés plutôt que beaucoup de verbes approximatifs.
\end{enumerate}
\end{methode}

\begin{exoresolu}[Complétion à trous]
\textbf{Énoncé.} Complete the sentences with the correct phrasal verb in the correct tense : \eng{stand up for, look into, put up with, take part in, speak out against}.\\
1. Last year, thousands of students \dots\ the peaceful march for free education.\\
2. Journalists should \dots\ injustice without fear.\\
3. The police promised to \dots\ the case of police brutality.\\
4. Every citizen must \dots\ the rights of minorities.\\
5. We cannot \dots\ corruption any longer.
\tcblower
\textbf{Solution détaillée.}\\
1. \eng{took part in} — « participer à » ; \eng{last year} impose le prétérit ; \eng{take} est irrégulier : take -- took -- taken.\\
2. \eng{speak out against} — « dénoncer » ; après le modal \eng{should}, base verbale.\\
3. \eng{look into} — « enquêter sur » ; après \eng{promised to}, infinitif.\\
4. \eng{stand up for} — « défendre » ; après \eng{must}, base verbale.\\
5. \eng{put up with} — « tolérer » ; après \eng{cannot}, base verbale.\\
\textbf{Conclusion :} seule la base verbale se conjugue ; la ou les particules restent invariables et la préposition finale fait partie du verbe (\eng{against}, \eng{for}, \eng{with}, \eng{into}).
\end{exoresolu}

\courssub{Savoir-faire 4 : Traduire sans calquer le français}

Le français exprime par un verbe simple ce que l'anglais courant exprime par un phrasal verb. Traduire, c'est donc souvent \textbf{changer de structure}, pas chercher un mot à mot.

\begin{center}
\begin{tabularx}{\linewidth}{|X|X|X|}
\hline
\rowcolor{popblueL} Français & English (courant) & English (soutenu) \\
\hline
abandonner & \eng{give up} & \eng{abandon} \\
\hline
remettre à plus tard & \eng{put off} & \eng{postpone} \\
\hline
tolérer & \eng{put up with} & \eng{tolerate} \\
\hline
enquêter sur & \eng{look into} & \eng{investigate} \\
\hline
supprimer, abolir & \eng{do away with} & \eng{abolish} \\
\hline
rejeter & \eng{turn down} & \eng{reject} \\
\hline
s'occuper de & \eng{look after} & \eng{take care of} \\
\hline
chercher (dans un dictionnaire) & \eng{look up} & \eng{consult} \\
\hline
\end{tabularx}
\end{center}

\begin{methode}[Éviter les calques et choisir le bon registre]
\begin{enumerate}
\item Ne traduis jamais un phrasal verb mot à mot : \eng{give up} n'a rien à voir avec « donner » + « haut ».
\item À l'écrit formel (lettre au ministre, article), préfère le verbe soutenu : \eng{postpone} plutôt que \eng{put off}, \eng{tolerate} plutôt que \eng{put up with}.
\item À l'oral et dans les dialogues, les phrasal verbs sont naturels : \eng{Come on, speak up!} « Allez, parle plus fort ! »
\item Attention aux faux amis : \eng{demand} = exiger (et non « demander » = \eng{ask for}) ; \eng{assist} = aider (et non « assister à » = \eng{attend}).
\item Vérifie la particule dans un dictionnaire : c'est elle qui porte le sens (\eng{look up} = chercher dans un dictionnaire, \eng{look up to} = admirer, \eng{look after} = s'occuper de).
\end{enumerate}
\end{methode}

\begin{exoresolu}[Traduction guidée]
\textbf{Énoncé.} Translate into English : « Les citoyens ont dénoncé la fraude électorale et ont demandé au tribunal d'examiner les plaintes. Ils ne toléreront plus l'injustice. »
\tcblower
\textbf{Solution détaillée.}\\
« dénoncer » $\rightarrow$ \eng{speak out against} ; prétérit : \eng{spoke out against} (speak -- spoke -- spoken).\\
« demander » $\rightarrow$ attention au faux ami : \eng{asked}, pas \eng{demanded} (qui signifie « exiger »). Structure : \eng{ask somebody to do something}.\\
« examiner » $\rightarrow$ \eng{look into} (après \eng{to}, base verbale).\\
« tolérer » $\rightarrow$ \eng{put up with} ; futur négatif : \eng{will not put up with} ; la particule \eng{with} reste en fin de groupe.\\
\textbf{Réponse rédigée :} \eng{The citizens spoke out against electoral fraud and asked the court to look into the complaints. They will no longer put up with injustice.}\\
\textbf{Conclusion :} une bonne traduction repose sur le choix du phrasal verb juste, la conjugaison correcte du verbe de base et le refus des calques du français.
\end{exoresolu}

\courssub{Savoir-faire 5 : Rédiger un court texte avec des phrasal verbs}

\begin{methode}[Produire un paragraphe argumentatif au Bac]
\begin{enumerate}
\item Prépare une mini-banque de 5 phrasal verbs du thème : \eng{stand up for, speak out against, look into, do away with, abide by}.
\item Structure ton paragraphe : phrase d'introduction (constat) $\rightarrow$ arguments $\rightarrow$ conclusion (appel à l'action).
\item Place chaque phrasal verb dans une phrase simple et correcte ; vérifie l'accord du verbe (3\up{e} personne du singulier : \eng{The government looks into\dots}).
\item Varie les temps : présent pour les vérités générales, prétérit pour les faits passés, futur ou modal pour les recommandations.
\item Relis en traquant : particule oubliée, pronom mal placé, verbe irrégulier mal conjugué.
\end{enumerate}
\end{methode}

\begin{exoresolu}[Composition courte]
\textbf{Énoncé.} Write a paragraph of about 80 words on : \eng{How can young Cameroonians defend human rights?} Use at least three phrasal verbs.
\tcblower
\textbf{Solution détaillée.}\\
Plan : 1) constat (les jeunes ont un rôle à jouer) ; 2) actions concrètes avec phrasal verbs ; 3) conclusion.\\
Vérifications : \eng{stand up for} + nom ; \eng{speak out against} + nom ; \eng{abide by} = respecter (la loi) ; \eng{take part in} + activité. Tous les verbes au présent de vérité générale, avec -s à la 3\up{e} personne du singulier si nécessaire.\\
\textbf{Réponse rédigée :} \eng{Young Cameroonians have a key role to play in the defence of human rights. First, they should stand up for the victims of injustice instead of remaining silent. They can also speak out against corruption and tribalism on social media and in their communities. Moreover, they must abide by the law themselves if they want others to respect it. Finally, they should take part in associations that look after vulnerable people, such as street children. If every young citizen acts, our society will gradually do away with injustice.}\\
\textbf{Conclusion :} six phrasal verbs correctement construits (\eng{stand up for, speak out against, abide by, take part in, look after, do away with}) donnent un texte naturel et valorisé par le correcteur.
\end{exoresolu}

\begin{attention}[Les 5 erreurs qui coûtent des points au Bac]
\begin{enumerate}
\item \textbf{Pronom mal placé :} écrire \eng{She gave up it} au lieu de \eng{She gave it up.} Avec un verbe séparable, le pronom va obligatoirement au milieu.
\item \textbf{Particule fausse ou oubliée :} confondre \eng{look for} (chercher), \eng{look after} (s'occuper de), \eng{look into} (enquêter sur), \eng{look up to} (admirer). La particule change tout le sens.
\item \textbf{Traduction mot à mot :} traduire \eng{give up} par « donner vers le haut » ou chercher un calque du français ; apprends chaque phrasal verb comme un mot nouveau avec sa traduction globale.
\item \textbf{Verbe de base mal conjugué :} \eng{He stand up for} au lieu de \eng{He stands up for} ; \eng{they speaked out} au lieu de \eng{they spoke out} (speak -- spoke -- spoken). Seul le verbe se conjugue, la particule ne bouge pas.
\item \textbf{Faux amis et registre :} utiliser \eng{demand} pour « demander » (= \eng{ask for}), \eng{assist} pour « assister à » (= \eng{attend}) ; ou multiplier les phrasal verbs familiers (\eng{put up with}) dans une lettre officielle où l'on attend \eng{tolerate}.
\end{enumerate}
\end{attention}

\newpage

\coursec{Exercices}

\courssub{Je vérifie mes connaissances}

\begin{exercice}[QCM — La bonne particule]
Entoure la particule qui complète correctement chaque phrasal verb.
\begin{enumerate}
\item The government must look \textbf{after / up / off} the rights of every citizen.
\item The students carried \textbf{on / out / in} a survey on youth participation in elections.
\item We cannot put \textbf{off / up / out} with corruption any longer.
\item The judge decided to look \textbf{into / after / over} the case of the arrested journalist.
\item During the crisis, many families ran \textbf{out of / up with / into} food and water.
\end{enumerate}
\end{exercice}

\begin{exercice}[Vrai ou faux — justifie ta réponse]
Dis si chaque phrase est correcte (True) ou incorrecte (False). Corrige les phrases fausses.
\begin{enumerate}
\item \eng{He gave up it last year.}
\item \eng{She looks after her grandmother every weekend.}
\item \eng{They put the meeting off until Friday.}
\item \eng{I am looking forward to vote in the next elections.}
\item \eng{The police are looking into the matter.}
\end{enumerate}
\end{exercice}

\begin{exercice}[Vocabulaire — Trouve le phrasal verb]
Retrouve le phrasal verb qui correspond à chaque définition. Choisis dans la liste : \eng{bring about, stand up for, give in, take over, set up, break down}.
\begin{enumerate}
\item to cause something to happen (change, progress) : \ldots
\item to defend a person or an idea : \ldots
\item to stop resisting and accept defeat : \ldots
\item to take control of something (a country, a company) : \ldots
\item to create or establish (an organisation, a committee) : \ldots
\item to stop working (a system, negotiations) : \ldots
\end{enumerate}
\end{exercice}

\begin{exercice}[Conjugaison et forme — Complète]
Complète les phrases avec le phrasal verb entre parenthèses, à la forme qui convient (temps, participe, gérondif).
\begin{enumerate}
\item Last year, our community \ldots\ a health centre. (to set up — past simple)
\item The minister promised that he would never \ldots\ to pressure. (to give in — base verbale)
\item The activists are \ldots\ the rights of street children. (to stand up for — present continuous)
\item \ldots\ early is a good habit for every student. (to get up — gérondif sujet)
\item The old law has been \ldots\ by the parliament. (to do away with — participe passé)
\end{enumerate}
\end{exercice}

\courssub{Je m'entraîne}

\begin{exercice}[Traduction — Du français vers l'anglais]
Traduis les phrases suivantes en anglais en utilisant un phrasal verb.
\begin{enumerate}
\item Les citoyens ont rejeté les mauvais dirigeants lors des élections.
\item Nous devons défendre les droits des femmes dans notre village.
\item La réunion a été reportée à la semaine prochaine.
\item Je ne peux pas tolérer l'injustice.
\item Les jeunes volontaires ont nettoyé le marché de Mokolo.
\end{enumerate}
\end{exercice}

\begin{exercice}[Transformation — Remplace le complément par un pronom]
Réécris chaque phrase en remplaçant le complément nominal souligné par le pronom approprié (\eng{it, them, him, her}). Attention à la place du pronom !
\begin{enumerate}
\item The people voted \textbf{the corrupt leaders} out.
\item The students put \textbf{the demonstration} off.
\item She looks after \textbf{her little brother}.
\item Please turn \textbf{the radio} on; the news is starting.
\item The committee turned \textbf{his application} down.
\end{enumerate}
\end{exercice}

\begin{exercice}[Texte à trous — Une élection au Cameroun]
Complète le texte avec les particules suivantes : \eng{up, out, off, after, into, with, for, on}. Chaque particule n'est utilisée qu'une fois.
\begin{textebox}[Election day in Yaoundé]
On election day, voters woke (1) \ldots\ early and queued (2) \ldots\ patiently in front of the polling stations. The electoral commission had set (3) \ldots\ transparent ballot boxes in every neighbourhood. Journalists carried (4) \ldots\ interviews with voters, while observers looked (5) \ldots\ the whole process carefully. When a problem was reported, officials promised to look (6) \ldots\ it immediately. Some candidates tried to put (7) \ldots\ the counting, but the population refused to put up (8) \ldots\ any fraud. In the end, the citizens voted the dishonest leaders (9) \ldots\ and brought (10) \ldots\ a new era of hope.
\end{textebox}
\end{exercice}

\begin{exercice}[Appariement — Phrasal verbs et synonymes]
Associe chaque phrasal verb de la colonne A à son équivalent de la colonne B.
\begin{center}
\begin{tabularx}{\linewidth}{|X|X|}
\hline
\rowcolor{popblueL} \textbf{Column A — Phrasal verb} & \textbf{Column B — Equivalent} \\
\hline
1. to put off & a. to investigate \\
\hline
2. to look into & b. to postpone \\
\hline
3. to put up with & c. to establish \\
\hline
4. to set up & d. to tolerate \\
\hline
5. to bring about & e. to cause \\
\hline
6. to do away with & f. to abolish \\
\hline
\end{tabularx}
\end{center}
\end{exercice}

\courssub{Je me prépare au Bac}

\begin{exercice}[Type Bac — Compréhension et langue]
Lis le texte suivant, puis réponds aux questions en anglais, avec tes propres mots.
\begin{textebox}[Youth and citizenship in Cameroon]
In many towns of Cameroon, young people are standing up for their rights more than ever before. In Douala, a group of students set up an association called “Young Voices” to bring about change in their community. Every Saturday, they carry out cleaning campaigns in the markets and visit schools to talk about civic education.

At first, the local authorities did not take them seriously. Some officials even tried to put off their meetings. But the young activists refused to give in. “We cannot put up with dirty streets and abandoned schools,” says Amina, the seventeen-year-old president of the association. “If the government looks into our problems, things will change.”

Little by little, the population joined them. Shopkeepers gave out free drinks during the campaigns, and parents came along to help. Today, “Young Voices” has taken over an old building and turned it into a library. Their example shows that when citizens work together, they can bring about real progress.
\end{textebox}
\textbf{Comprehension questions}
\begin{enumerate}
\item What is the name of the association created by the students, and what is its goal? (2 marks)
\item Give two activities that the members of the association carry out. (2 marks)
\item How did the local authorities react at the beginning? (2 marks)
\item According to Amina, what will happen if the government examines their problems? (2 marks)
\item What lesson can we learn from this text about citizenship? (2 marks)
\end{enumerate}
\textbf{Language work}
\begin{enumerate}
\item Find in the text three phrasal verbs and give their meaning in English or in French. (3 marks)
\item Rewrite this sentence, replacing the underlined words with a pronoun: \eng{The students put off \textbf{their meetings}.} (1 mark)
\item Replace the underlined verb with the correct phrasal verb: \eng{The activists \textbf{tolerated} the situation for two years.} (1 mark)
\end{enumerate}
\end{exercice}

\begin{exercice}[Type Bac — Traduction et correction d'erreurs]
\textbf{A. Translation (5 marks).} Traduis le passage suivant en anglais.

\begin{textebox}[Texte à traduire]
Les jeunes Camerounais veulent participer à la vie politique de leur pays. Ils ont créé des associations pour défendre leurs droits. Ils ne peuvent plus tolérer la corruption et l'injustice. Quand les citoyens travaillent ensemble, ils peuvent provoquer de grands changements.
\end{textebox}

\textbf{B. Error correction (5 marks).} Chaque phrase contient une erreur de phrasal verb. Souligne l'erreur et réécris la phrase correctement.
\begin{enumerate}
\item \eng{He gave up it before the examination.}
\item \eng{We are looking forward to meet the minister.}
\item \eng{She put off with the noise for three hours.}
\item \eng{The police looked the case into yesterday.}
\item \eng{They carried on the survey out last month.}
\end{enumerate}
\end{exercice}

\begin{exercice}[Type Bac — Essay practice]
\textbf{Topic:} \eng{How young Cameroonians can bring about change.}

Rédige un paragraphe argumentatif de \textbf{60 à 80 mots} en anglais sur ce sujet. Tu dois :
\begin{enumerate}
\item utiliser \textbf{au moins cinq phrasal verbs} différents (souligne-les dans ta copie) ;
\item présenter une opinion claire et au moins deux arguments ;
\item respecter la grille d'évaluation du Bac : pertinence des idées (content), organisation logique (organisation), correction grammaticale (accuracy) et richesse du vocabulaire (vocabulary).
\end{enumerate}
\textbf{Useful phrasal verbs:} \eng{stand up for, bring about, set up, carry out, give in, put up with, look after, take part in, do away with, work out}.
\end{exercice}

\newpage

\coursec{Corrigés des exercices}

\begin{corrige}[Exercice 1 — QCM : La bonne particule]
\begin{enumerate}
\item \textbf{after} — \eng{look after} = take care of / s'occuper de, veiller sur. \eng{The government must look after the rights of every citizen.} (Le gouvernement doit veiller aux droits de chaque citoyen.)
\item \textbf{out} — \eng{carry out} = conduct, perform / mener, réaliser. \eng{The students carried out a survey…} (Les étudiants ont réalisé une enquête…)
\item \textbf{up with} — \eng{put up with} = tolerate / tolérer, supporter. \eng{We cannot put up with corruption any longer.} (Nous ne pouvons plus tolérer la corruption.)
\item \textbf{into} — \eng{look into} = investigate / enquêter sur, examiner. \eng{The judge decided to look into the case…} (Le juge a décidé d'examiner l'affaire…)
\item \textbf{out of} — \eng{run out of} = have no more left / manquer de, être à court de. \eng{Many families ran out of food and water.} (De nombreuses familles ont manqué de nourriture et d'eau.)
\end{enumerate}
\end{corrige}

\begin{corrige}[Exercice 2 — Vrai ou faux : justifie ta réponse]
\begin{enumerate}
\item \textbf{False.} \eng{He gave up it last year.} $\rightarrow$ \textbf{Correction :} \eng{He gave \textbf{it} up last year.} \\
\emph{Justification :} \eng{give up} est un \textbf{phrasal verb séparable}. Lorsque le complément est un pronom (\eng{it, them, him, her}), il \textbf{doit} se placer \textbf{entre} le verbe et la particule.
\item \textbf{True.} \eng{She looks after her grandmother every weekend.} \\
\emph{Justification :} \eng{look after} est \textbf{inséparable} ; le complément (nom ou pronom) se place toujours \textbf{après} la particule.
\item \textbf{True.} \eng{They put the meeting off until Friday.} \\
\emph{Justification :} \eng{put off} est \textbf{séparable}. Avec un complément nominal (\eng{the meeting}), on peut le placer après la particule (\eng{put off the meeting}) ou entre les deux (\eng{put the meeting off}). Les deux sont corrects.
\item \textbf{False.} \eng{I am looking forward to vote in the next elections.} $\rightarrow$ \textbf{Correction :} \eng{I am looking forward to \textbf{voting} in the next elections.} \\
\emph{Justification :} Dans \eng{look forward to}, \eng{to} est une \textbf{préposition} (et non la marque de l'infinitif). Elle est donc suivie du \textbf{gérondif} (\eng{-ing}).
\item \textbf{True.} \eng{The police are looking into the matter.} \\
\emph{Justification :} \eng{look into} est \textbf{inséparable} (verbe + particule + préposition \eng{into}) ; le complément suit l'ensemble.
\end{enumerate}
\end{corrige}

\begin{corrige}[Exercice 3 — Vocabulaire : Trouve le phrasal verb]
\begin{enumerate}
\item \textbf{bring about} — \eng{to cause something to happen (change, progress)} (provoquer, entraîner un changement).
\item \textbf{stand up for} — \eng{to defend a person or an idea} (défendre une personne, une idée).
\item \textbf{give in} — \eng{to stop resisting and accept defeat} (céder, abandonner la résistance).
\item \textbf{take over} — \eng{to take control of something (a country, a company)} (prendre le contrôle de, reprendre).
\item \textbf{set up} — \eng{to create or establish (an organisation, a committee)} (créer, établir, mettre sur pied).
\item \textbf{break down} — \eng{to stop working (a system, negotiations)} (tomber en panne, échouer [pour négociations/système]).
\end{enumerate}
\end{corrige}

\begin{corrige}[Exercice 4 — Conjugaison et forme : Complète]
\begin{enumerate}
\item Last year, our community \textbf{set up} a health centre. (\eng{to set up} — \emph{past simple} : forme irrégulière \eng{set - set - set}).
\item The minister promised that he would never \textbf{give in} to pressure. (\eng{to give in} — \emph{base verbale / infinitif sans to} après le modal \eng{would}).
\item The activists \textbf{are standing up for} the rights of street children. (\eng{to stand up for} — \emph{present continuous} : \eng{be + V-ing} ; \eng{standing} pour \eng{stand}).
\item \textbf{Getting} up early is a good habit for every student. (\eng{to get up} — \emph{gérondif sujet} : \eng{V-ing} en début de phrase).
\item The old law has been \textbf{done away with} by the parliament. (\eng{to do away with} — \emph{participe passé} à la voix passive : \eng{been + V-en/ed} ; forme irrégulière \eng{do - did - done} $\rightarrow$ \eng{done away with}).
\end{enumerate}
\end{corrige}

\begin{corrige}[Exercice 5 — Traduction : Du français vers l'anglais]
\begin{enumerate}
\item \eng{The citizens \textbf{voted the bad leaders out} (or \textbf{voted out the bad leaders}) during the elections.} (\eng{vote out} = révoquer, chasser par le vote ; séparable).
\item \eng{We must \textbf{stand up for} women's rights in our village.} (\eng{stand up for} = défendre ; inséparable).
\item \eng{The meeting \textbf{was put off} until next week.} (\eng{put off} = reporter ; voix passive, donc le sujet \eng{the meeting} remonte).
\item \eng{I \textbf{cannot put up with} injustice.} (\eng{put up with} = tolérer ; inséparable, suivi de l'objet \eng{injustice}).
\item \eng{The young volunteers \textbf{cleaned up} the Mokolo market.} (or \eng{\textbf{cleaned} the Mokolo market \textbf{up}}.) (\eng{clean up} = nettoyer, assainir ; séparable).
\end{enumerate}
\end{corrige}

\begin{corrige}[Exercice 6 — Transformation : Remplace le complément par un pronom]
\textbf{Règle :} Avec un \textbf{phrasal verb séparable}, le pronom objet (\eng{it, them, him, her}) \textbf{doit} se placer \textbf{entre} le verbe et la particule. Avec un \textbf{phrasal verb inséparable}, le pronom se place \textbf{après} la particule (ou la préposition finale).
\begin{enumerate}
\item \eng{The people voted \textbf{them} out.} (\eng{vote out} séparable $\rightarrow$ pronom au milieu).
\item \eng{The students put \textbf{it} off.} (\eng{put off} séparable $\rightarrow$ pronom au milieu).
\item \eng{She looks after \textbf{him}.} (\eng{look after} inséparable $\rightarrow$ pronom après la particule).
\item \eng{Please turn \textbf{it} on; the news is starting.} (\eng{turn on} séparable $\rightarrow$ pronom au milieu).
\item \eng{The committee turned \textbf{it} down.} (\eng{turn down} = refuser, séparable $\rightarrow$ pronom au milieu).
\end{enumerate}
\end{corrige}

\begin{corrige}[Exercice 7 — Texte à trous : Une élection au Cameroun]
\textbf{Attention :} La liste fournie (8 particules pour 10 trous) ne permet pas de respecter l'instruction « chaque particule une seule fois » car certaines particules (\eng{up, out}) sont nécessaires deux fois et la particule \eng{about} (pour \eng{bring about}) manque à la liste. La correction linguistique exacte est donnée ci-dessous.
\begin{enumerate}
\item \textbf{up} — \eng{woke \textbf{up}} early (se réveiller).
\item \textbf{up} — \eng{queued \textbf{up}} patiently (faire la queue ; \eng{queue up} courant).
\item \textbf{up} — \eng{had set \textbf{up}} transparent ballot boxes (installer, mettre en place).
\item \textbf{out} — \eng{carried \textbf{out}} interviews (mener, réaliser).
\item \textbf{on} — \eng{observers looked \textbf{on}} the whole process carefully (regarder, observer sans intervenir ; \eng{look on}).
\item \textbf{into} — \eng{promised to look \textbf{into}} it immediately (enquêter sur).
\item \textbf{off} — \eng{tried to put \textbf{off}} the counting (reporter).
\item \textbf{up with} — \eng{refused to put \textbf{up with}} any fraud (tolérer ; \eng{put up with} inséparable à 3 mots).
\item \textbf{out} — \eng{voted the dishonest leaders \textbf{out}} (chasser par le vote).
\item \textbf{about} — \eng{brought \textbf{about}} a new era of hope (provoquer, amener ; \eng{bring about} inséparable). \textit{Particule manquante dans la liste fournie.}
\end{enumerate}
\end{corrige}

\begin{corrige}[Exercice 8 — Appariement : Phrasal verbs et synonymes]
\begin{center}
\begin{tabularx}{\linewidth}{|X|X|X|}
\hline
\rowcolor{popblueL} \textbf{N°} & \textbf{Phrasal verb (Colonne A)} & \textbf{Equivalent (Colonne B)} \\
\hline
1 & to put off & \textbf{b. to postpone} \\
\hline
2 & to look into & \textbf{a. to investigate} \\
\hline
3 & to put up with & \textbf{d. to tolerate} \\
\hline
4 & to set up & \textbf{c. to establish} \\
\hline
5 & to bring about & \textbf{e. to cause} \\
\hline
6 & to do away with & \textbf{f. to abolish} \\
\hline
\end{tabularx}
\end{center}
\end{corrige}

\begin{corrige}[Exercice 9 — Type Bac : Compréhension et langue]
\textbf{Barème indicatif : 10 marks (Comprehension) + 5 marks (Language) = 15 marks}

\end{corrige}

\courssub{Comprehension questions (10 marks)}
\begin{enumerate}
\item \textbf{Name:} “Young Voices”. \textbf{Goal:} To \eng{bring about change} in their community / to bring about change in Douala. (2 marks)
\item \textbf{Activities (any two):} \eng{carry out cleaning campaigns in the markets} ; \eng{visit schools to talk about civic education}. (2 marks)
\item \textbf{Reaction:} They \eng{did not take them seriously} and \eng{tried to put off their meetings}. (2 marks)
\item \textbf{Amina's view:} If the government \eng{looks into their problems}, \eng{things will change}. (2 marks)
\item \textbf{Lesson:} When citizens \eng{work together}, they can \eng{bring about real progress} / Active citizenship and collective action lead to positive change. (2 marks)
\end{enumerate}

\courssub{Language work (5 marks)}
\begin{enumerate}
\item \textbf{Three phrasal verbs from the text (1 mark each):}
\begin{itemize}
\item \eng{stand up for} = defend / défendre
\item \eng{set up} = create, establish / créer, mettre sur pied
\item \eng{bring about} = cause, provoke / provoquer, entraîner
\item \eng{carry out} = conduct, do / mener, réaliser
\item \eng{put off} = postpone / reporter
\item \eng{give in} = surrender, yield / céder
\item \eng{put up with} = tolerate / tolérer
\item \eng{look into} = investigate / enquêter sur
\item \eng{take over} = take control of / prendre le contrôle de
\item \eng{turn into} = transform into / transformer en
\end{itemize}
\item \eng{The students put \textbf{them} off.} (1 mark) — \emph{Règle : pronom objet placé entre verbe et particule pour phrasal verb séparable.}
\item \eng{The activists \textbf{put up with} the situation for two years.} (1 mark) — \eng{tolerate} $\rightarrow$ \eng{put up with} (inséparable, 3 mots).
\end{enumerate}

\courssub{Points de vigilance (Bac)}
\begin{itemize}
\item Répondre \textbf{en anglais} et \textbf{avec ses propres mots} (pas de copier-coller brut).
\item Pour \textit{Language work} : bien distinguer séparables / inséparables (place du pronom) et la forme \eng{-ing} après \eng{look forward to}.
\item Vérifier l'orthographe des verbes irréguliers (\eng{set - set - set}, \eng{bring - brought - brought}).
\end{itemize}


\begin{corrige}[Exercice 10 — Type Bac : Traduction et correction d'erreurs]
\end{corrige}

\courssub{A. Translation (5 marks)}
\eng{Young Cameroonians want to \textbf{take part in} the political life of their country. They have \textbf{set up} associations to \textbf{stand up for} their rights. They can no longer \textbf{put up with} corruption and injustice. When citizens \textbf{work together}, they can \textbf{bring about} great changes.}

\emph{Notes de traduction :}
\begin{itemize}
\item \eng{participer à} $\rightarrow$ \eng{take part in} (phrasal verb à 3 mots, inséparable).
\item \eng{créé} $\rightarrow$ \eng{have set up} (present perfect, \eng{set up} séparable).
\item \eng{défendre} $\rightarrow$ \eng{stand up for} (inséparable).
\item \eng{tolérer} $\rightarrow$ \eng{put up with} (inséparable).
\item \eng{provoquer} $\rightarrow$ \eng{bring about} (inséparable).
\item \eng{travailler ensemble} $\rightarrow$ \eng{work together} (verbe + adverbe, pas un phrasal verb lexicalisé ici, mais acceptable). On pourrait aussi utiliser \eng{team up}.
\end{itemize}

\courssub{B. Error correction (5 marks — 1 mark par phrase corrigée)}
\begin{enumerate}
\item \eng{He gave \textbf{up it} before the examination.} $\rightarrow$ \textbf{Correction :} \eng{He gave \textbf{it up} before the examination.} \\
\emph{Erreur :} Place du pronom avec phrasal verb séparable (\eng{give up}).
\item \eng{We are looking forward to \textbf{meet} the minister.} $\rightarrow$ \textbf{Correction :} \eng{We are looking forward to \textbf{meeting} the minister.} \\
\emph{Erreur :} \eng{to} est une préposition ici $\rightarrow$ gérondif (\eng{-ing}).
\item \eng{She put \textbf{off with} the noise for three hours.} $\rightarrow$ \textbf{Correction :} \eng{She \textbf{put up with} the noise for three hours.} \\
\emph{Erreur :} Confusion de particules. \eng{put off} = reporter ; \eng{put up with} = tolérer.
\item \eng{The police \textbf{looked the case into} yesterday.} $\rightarrow$ \textbf{Correction :} \eng{The police \textbf{looked into the case} yesterday.} \\
\emph{Erreur :} \eng{look into} est inséparable (verbe + particule + préposition), l'objet suit l'ensemble.
\item \eng{They \textbf{carried on the survey out} last month.} $\rightarrow$ \textbf{Correction :} \eng{They \textbf{carried out the survey} last month.} (ou \eng{carried the survey out}.) \\
\emph{Erreur :} Ordre des mots incorrect. \eng{carry out} est séparable ; l'objet nominal se place après la particule ou entre les deux, mais pas \eng{on [obj] out}.
\end{enumerate}


\begin{corrige}[Exercice 11 — Type Bac : Essay practice]
\end{corrige}

\courssub{Proposition de rédaction modèle (environ 73 mots)}
Young Cameroonians can \uline{\textbf{bring about}} significant change if they \uline{\textbf{stand up for}} their rights and \uline{\textbf{take part in}} civic life. First, they should \uline{\textbf{set up}} youth associations to \uline{\textbf{carry out}} community projects like cleaning campaigns or voter education. Second, they must not \uline{\textbf{give in}} to apathy or corruption but \uline{\textbf{work out}} solutions with local leaders. By refusing to \uline{\textbf{put up with}} injustice, the youth become the driving force of a better future.

\courssub{Analyse par rapport aux consignes}
\begin{itemize}
\item \textbf{Mots :} $\approx$ 73 mots (dans la fourchette 60--80). \checkmark
\item \textbf{Phrasal verbs utilisés (8, dont 5 requis) :} \uline{\textbf{bring about}}, \uline{\textbf{stand up for}}, \uline{\textbf{take part in}}, \uline{\textbf{set up}}, \uline{\textbf{carry out}}, \uline{\textbf{give in}}, \uline{\textbf{work out}}, \uline{\textbf{put up with}}. \checkmark
\item \textbf{Opinion claire :} La jeunesse est le moteur du changement. \checkmark
\item \textbf{Deux arguments :} 1. Créer des associations pour agir concrètement. 2. Refuser la résignation et collaborer avec les autorités. \checkmark
\item \textbf{Organisation :} Connecteurs logiques (\eng{First, Second, By refusing...}). \checkmark
\item \textbf{Correction grammaticale :} Modaux (\eng{can, should, must not}), gérondif (\eng{refusing to}), voix active/passive correctes. \checkmark
\end{itemize}

\courssub{Points de vigilance pour l'élève}
\begin{itemize}
\item \textbf{Souligner} visiblement les phrasal verbs dans la copie (surligneur ou trait de soulignement).
\item Ne pas dépasser 80 mots : aller à l'essentiel, phrases concises.
\item Vérifier l'accord des temps (présent pour faits généraux, modaux pour recommandations).
\item Éviter la répétition du sujet \eng{Young Cameroonians} / \eng{They} ; utiliser des synonymes ou des pronoms.
\end{itemize}


\newpage

\coursec{Fiche bilan}

\courssub{Carte mentale : Phrasal Verbs}

\begin{popfigure}
\begin{tikzpicture}[
    mindmap,
    grow cyclic,
    every node/.style={concept, execute at begin node=\hskip0pt},
    concept color=popblue,
    level 1/.style={level distance=4.5cm, sibling angle=60, concept color=popblue!70},
    level 2/.style={level distance=3cm, sibling angle=45, concept color=popgreen!70},
    text=white,
    font=\small\bfseries
]
\node[concept, minimum size=2.8cm, align=center] (center){Phrasal Verbs\\ \textit{(Verbes à particule)}}
    child[concept color=poporange] { node[concept, minimum size=2.2cm, align=center]{1. Structure\\ Verb + Particle\\ (Prep/Adv)} }
    child[concept color=poppurple] { node[concept, minimum size=2.2cm, align=center]{2. Types\\ Separable / Inseparable\\ Three-word} }
    child[concept color=poppink] { node[concept, minimum size=2.2cm, align=center]{3. Grammar Rules\\ Object position\\ Pronouns} }
    child[concept color=popgreen] { node[concept, minimum size=2.2cm, align=center]{4. Citizenship Lexicon\\ stand up for, speak out,\\ crack down on} }
    child[concept color=popgold] { node[concept, minimum size=2.2cm, align=center]{5. French Traps\\ Literal translation\\ False friends} }
    child[concept color=popdark] { node[concept, minimum size=2.2cm, align=center]{6. Exam Skills\\ Transformation\\ Gap-fill / Writing} };
\end{tikzpicture}
\legende{Synthèse visuelle de la leçon 41 : Phrasal Verbs (Citoyenneté et Droits de l'Homme)}
\end{popfigure}

\courssub{Fiche Bilan : Les Essentiels à Retenir}

\begin{aretenir}[Lexique clé — Citizenship \& Human Rights]
\begin{center}
\begin{tabularx}{\linewidth}{|X|X|X|}
\hline
\rowcolor{popblueL}
\textbf{Phrasal Verb} & \textbf{Signification} & \textbf{Exemple contextuel} \\
\hline
\eng{stand up for} & défendre (une cause, qqn) & \eng{We must stand up for human rights.} « Nous devons défendre les droits de l'homme. » \\
\hline
\eng{speak out (against)} & s'exprimer publiquement (contre) & \eng{Journalists speak out against corruption.} « Les journalistes dénoncent la corruption. » \\
\hline
\eng{crack down on} & sévir contre, réprimer & \eng{The government cracked down on illegal logging.} « Le gouvernement a sévi contre l'exploitation illégale. » \\
\hline
\eng{call for} & réclamer, demander & \eng{NGOs call for more transparency.} « Les ONG réclament plus de transparence. » \\
\hline
\eng{bring about} & provoquer, entraîner (un changement) & \eng{The protest brought about a new law.} « La manifestation a entraîné une nouvelle loi. » \\
\hline
\eng{look into} & enquêter sur, examiner & \eng{The commission will look into the allegations.} « La commission enquêtera sur les allégations. » \\
\hline
\eng{put up with} & tolérer, supporter & \eng{Citizens should not put up with injustice.} « Les citoyens ne devraient pas tolérer l'injustice. » \\
\hline
\eng{come up with} & proposer, trouver (une solution) & \eng{We need to come up with a peace plan.} « Nous devons trouver un plan de paix. » \\
\hline
\eng{go against} & aller à l'encontre de & \eng{That decision goes against the Constitution.} « Cette décision va à l'encontre de la Constitution. » \\
\hline
\eng{live up to} & être à la hauteur de, honorer & \eng{Leaders must live up to their promises.} « Les dirigeants doivent honorer leurs promesses. » \\
\hline
\end{tabularx}
\end{center}
\end{aretenir}

\begin{propriete}[Règles de grammaire essentielles]
\begin{center}
\begin{tabularx}{\linewidth}{|X|X|X|}
\hline
\rowcolor{popblueL}
\textbf{Type} & \textbf{Règle de structure} & \textbf{Exemple / Piège} \\
\hline
\textbf{Séparables} (Transitifs) & Objet \textbf{nom} : entre les deux OU après.\newline Objet \textbf{pronom} : \textbf{OBLIGATOIREMENT entre les deux}. & \eng{Turn off the light.} / \eng{Turn the light off.}\newline \eng{Turn \textbf{it} off.} (\textcolor{red}{Pas *Turn off it}) \\
\hline
\textbf{Inséparables} (Transitifs) & Verbe + Particule \textbf{jamais séparés}. Objet (nom ou pronom) \textbf{toujours après}. & \eng{Look into the matter.} / \eng{Look into \textbf{it}.} (\textcolor{red}{Pas *Look the matter into}) \\
\hline
\textbf{Intransitifs} & Pas d'objet. Verbe + Particule inséparables. & \eng{The meeting broke up early.} \\
\hline
\textbf{À trois mots} (Verb + Part + Prep) & \textbf{Jamais séparables}. Suivis d'un objet (nom/pronom/\eng{-ing}). & \eng{look forward to \textbf{it/seeing you}} \\
\hline
\textbf{Passif} & Possible seulement pour \textbf{séparables} et \textbf{inséparables transitifs}. Particule reste collée au verbe. & \eng{The meeting was \textbf{called off}.} (Separable: \eng{call sth off}) \\
\hline
\end{tabularx}
\end{center}
\end{propriete}

\begin{methode}[Identifier le type d'un Phrasal Verb]
\begin{enumerate}
\item Consulte le dictionnaire : \eng{v.i.} (intransitif) ou \eng{v.t.} (transitif).
\item Si \eng{v.t.} : teste avec un pronom (\eng{it}, \eng{them}). \eng{Put it off} $\rightarrow$ Séparable. \eng{Look into it} $\rightarrow$ Inséparable.
\item Compte les mots : 2 mots (Verb + Part) ou 3 mots (Verb + Part + Prep).
\end{enumerate}
\end{methode}

\begin{methode}[Transformer Actif $\rightarrow$ Passif (Bac Writing)]
\begin{enumerate}
\item Vérifie que le PV est \textbf{transitif} (séparable ou inséparable).
\item Fais devenir l'objet le sujet de la phrase passive.
\item Conjugue \eng{BE} + \textbf{Participe Passé} du verbe principal. \textbf{Garde la particule collée}.
\item Exemple : \eng{They \textbf{called off} the strike.} $\rightarrow$ \eng{The strike \textbf{was called off}.}
\end{enumerate}
\end{methode}

\begin{attention}[Pièges « Spécial Francophones »]
\begin{itemize}
\item \textbf{Traduction littérale :} \eng{Make up} $\neq$ « fabriquer en haut » $\rightarrow$ \textbf{inventer} (une histoire) / \textbf{se réconcilier} / \textbf{se maquiller}.
\item \textbf{Séparation forcée du pronom :} \textcolor{red}{*She turned off it.} $\rightarrow$ \textbf{She turned \eng{it} off.}
\item \textbf{Oubli de la préposition finale (3 mots) :} \textcolor{red}{*I look forward to see you.} $\rightarrow$ \textbf{I look forward to \eng{seeing} you.} (\eng{to} = préposition $\rightarrow$ \eng{-ing}).
\item \textbf{Confusion \eng{make up} / \eng{make out} :} \eng{Make out} = distinguer à peine / remplir (un chèque) / s'embrasser (US).
\item \textbf{Registre :} Les PV sont souvent \textbf{informels / neutres}. À l'écrit formel (lettre officielle, dissertation), préfère le verbe latin : \eng{discover} (vs \eng{find out}), \eng{tolerate} (vs \eng{put up with}), \eng{investigate} (vs \eng{look into}).
\end{itemize}
\end{attention}

\courssub{Checklist avant le Bac}

\begin{aretenir}[Checklist Phrasal Verbs — Terminale A]
\begin{itemize}
\item $\square$ Je sais définir un phrasal verb : \eng{Verb + Particle (adverbe/préposition)} changeant le sens du verbe de base.
\item $\square$ Je distingue les \textbf{3 types} : Séparables, Inséparables (transitifs), Intransitifs, et à \textbf{3 mots}.
\item $\square$ Je place \textbf{correctement l'objet pronom} : \textbf{toujours au milieu} pour les séparables (\eng{turn it off}).
\item $\square$ Je ne sépare \textbf{jamais} les inséparables (\eng{look into it}) ni les verbes à 3 mots (\eng{put up with it}).
\item $\square$ Je connais \textbf{10 phrasal verbs du thème Citoyenneté/DDH} (\eng{stand up for, crack down on, speak out, call for...}).
\item $\square$ Je sais former le \textbf{passif} : \eng{The bill was passed / The issue was looked into}.
\item $\square$ J'évite le piège \eng{to + V-ing} après les verbes à 3 mots (\eng{look forward to hearing}).
\item $\square$ Je sais choisir le \textbf{registre approprié} : PV à l'oral/informel $\leftrightarrow$ Verbe latin (simple) à l'écrit formel.
\item $\square$ Je réussis l'exercice de \textbf{transformation} (Remplacer le verbe souligné par un PV correct / Actif $\rightarrow$ Passif).
\item $\square$ Je relis mes phrases pour vérifier l'\textbf{orthographe des particules} (\eng{give up} $\neq$ \eng{give-up}).
\item $\square$ Je connais la prononciation : accent tonique souvent sur la \textbf{particule} (\eng{turn \textbf{OFF}}, \eng{LOOK into}).
\end{itemize}
\end{aretenir}

\findecours

\end{document}
